% Éradication des préjugés autour de l'apparition des anomalies et/ou des caractères nouveaux au sein des familles — SVT Tle D — Cours signé OnBuch+
% Programme officiel MINESEC. Compiler avec : tectonic N05-mutations-anomalies-geniques-chromosomiques.tex
\newcommand{\DOCMATIERE}{SVT}
\newcommand{\DOCNIVEAU}{Tle D}
\newcommand{\DOCMODULE}{Module I : Le Monde Vivant — Origine de nouveaux allèles}
\newcommand{\DOCLECON}{N05}
\newcommand{\DOCTITRE}{Éradication des préjugés autour de l'apparition des anomalies et/ou des caractères nouveaux au sein des familles}
% OnBuch+ — préambule « cours signés » ENRICHI (compilé avec tectonic / XeLaTeX)
% Système visuel « Pop » d'OnBuch+ : fond crème (jamais blanc), bordures encre
% épaisses, ombres « dures » décalées, titres Archivo Black en capitales.
% Version MATHÉMATIQUES : graphes (pgfplots/tikz), boîtes pédagogiques
% (définition, propriété, méthode, attention, le savais-tu, exercice résolu,
% exercice, TP, bilan), page de garde et signature OnBuch+ ancrée en bas.
%
% Macros attendues dans main.tex AVANT \input{preamble} :
%   \DOCMATIERE  \DOCNIVEAU  \DOCMODULE  \DOCLECON (numéro)  \DOCTITRE
\documentclass[11pt]{article}

\usepackage{fontspec}
\usepackage[french]{babel}
\usepackage[a4paper,margin=2cm,top=3.4cm,bottom=2.8cm,headheight=1.4cm,footskip=1.3cm]{geometry}
\usepackage{amsmath,amssymb,amsfonts}
\usepackage{mathtools} % \xmapsto, \coloneqq... souvent utilisés par les modèles
\usepackage{cancel} % simplifications barrées (bilans de forces)
% \abs{z} (module, valeur absolue) et \norm{u} (norme), utilisés par les modèles
\providecommand{\abs}{}\let\abs\relax\DeclarePairedDelimiter{\abs}{\lvert}{\rvert}
\providecommand{\norm}{}\let\norm\relax\DeclarePairedDelimiter{\norm}{\lVert}{\rVert}
\usepackage[table]{xcolor} % [table] : fournit aussi \rowcolors (alternance de lignes)
\usepackage{pagecolor}
\usepackage{tikz}
\usetikzlibrary{arrows.meta,positioning,calc,shapes.geometric,shapes.misc,shapes.arrows,decorations.pathmorphing,decorations.pathreplacing,decorations.markings,patterns,shadows,mindmap,trees,fit,backgrounds,matrix,intersections,angles,quotes}
\usepackage{pgfplots}
\usepackage[version=4]{mhchem} % équations nucléaires \ce{^{235}_{92}U}
\usepackage[european,siunitx]{circuitikz} % schémas électriques
\pgfplotsset{compat=1.18}
\usepgfplotslibrary{fillbetween,groupplots} % aires entre courbes (calcul intégral)
\usepackage{enumitem}
\usepackage{fancyhdr}
% [most] charge toutes les bibliothèques tcolorbox (listings, minted, raster,
% documentation...) et gonfle inutilement la mémoire TeX sur un document long
% (~300 boîtes) : on ne charge que ce qui est réellement utilisé.
\usepackage[skins,breakable]{tcolorbox}
\usepackage{graphicx}
\usepackage{colortbl}
\usepackage{booktabs}
\usepackage{tabularx}
\usepackage{diagbox} % case d'en-tête diagonale des tableaux à double entrée
% Colonnes extensibles centrée / gauche / droite (C, L, R), souvent utilisées par les modèles
\newcolumntype{C}{>{\centering\arraybackslash}X}
\newcolumntype{L}{>{\raggedright\arraybackslash}X}
\newcolumntype{R}{>{\raggedleft\arraybackslash}X}
\usepackage{array}
\usepackage{multicol}
\usepackage{siunitx}
% parse-numbers=false : accepte aussi \SI{2,5 \times 10^{-3}}{...}
\sisetup{locale=FR,output-decimal-marker={,},group-minimum-digits=4,parse-numbers=false}
\DeclareSIUnit\ohm{\ensuremath{\symup{\Omega}}} % U+2126 absent de la police
\DeclareSIUnit\division{div} % réglages d'oscilloscope (ms/div, V/div)
\DeclareSIUnit\tr{tr} % tours (vitesse de rotation, ex. \SI{1500}{\tr\per\minute})
\DeclareSIUnit\molar{M} % concentration molaire (ex. \SI{3}{\molar}, \SI{1}{\micro\molar})
\DeclareSIUnit\an{an} % année (le modèle confond parfois avec \year, un compteur TeX) — ex. \SI{5}{\kilo\meter\per\an}
\DeclareSIUnit\FCFA{FCFA} % franc CFA (ex. \SI{500}{\FCFA\per\kilo\gram})
\DeclareSIUnit\degreeBrix{\textdegree Brix} % teneur en sucre d'un sirop/jus (réfractométrie)
\DeclareSIUnit\month{mois} % le modèle utilise parfois \month, un compteur TeX, comme unité de durée
\usepackage{qrcode}
% \degree : commande fréquemment utilisée par les modèles pour un angle ou une
% température en dehors de \SI{}{} (non fournie par les paquets chargés).
\providecommand{\degree}{^{\circ}}
% \qed (fin de démonstration), utilisé par les modèles sans amsthm
\providecommand{\qed}{\hfill$\square$}
% \barre : double barre des valeurs interdites dans un tableau de variations (tkz-tab)
\providecommand{\barre}{\ensuremath{\|}}
% environnement proof (amsthm non chargé) utilisé par les modèles
\ifdefined\proof\else\newenvironment{proof}[1][Démonstration]{\par\noindent\textit{#1.}\ }{\qed\par}\fi
\usepackage{lastpage}
\usepackage{hyperref}
\hypersetup{hidelinks}

% ============================================================
% Palette « Pop » OnBuch+ — mêmes hex que la classe P (plus_ui.dart)
% ============================================================
\definecolor{poporange}{HTML}{F59321}
\definecolor{popdark}{HTML}{E07A0C}
\definecolor{popcream}{HTML}{FFF6E8}
\definecolor{popink}{HTML}{1C1714}
\definecolor{poppurple}{HTML}{7A5AE0}
\definecolor{popgreen}{HTML}{2E9E6A}
\definecolor{poppink}{HTML}{E0457B}
\definecolor{popblue}{HTML}{2F7DE1}
\definecolor{popgold}{HTML}{C98A12}
\definecolor{popmuted}{HTML}{8A7F76}
\definecolor{popline}{HTML}{EADFD2}
\definecolor{popgray}{HTML}{8A7F76}
\colorlet{popgrey}{popgray}
% Alias tolérants (le modèle nomme parfois les couleurs différemment)
\colorlet{popwhite}{white}
\colorlet{popblack}{popink}
\colorlet{popred}{poppink}
\colorlet{popyellow}{popgold}
% Teintes claires (fonds de tableaux/encadrés) — le modèle nomme parfois
% les couleurs avec un suffixe L (ex. popblueL) sans les avoir définies.
\colorlet{poporangeL}{poporange!20}
\colorlet{popdarkL}{popdark!20}
\colorlet{poppurpleL}{poppurple!20}
\colorlet{popgreenL}{popgreen!20}
\colorlet{poppinkL}{poppink!20}
\colorlet{popblueL}{popblue!20}
\colorlet{popgoldL}{popgold!20}
\colorlet{popredL}{popred!20}
\colorlet{popgrayL}{popgray!20}
% Teintes claires pour fonds de tableaux / zones
\colorlet{popblueL}{popblue!10!white}
\colorlet{poporangeL}{poporange!14!white}
\colorlet{popgreenL}{popgreen!12!white}
\colorlet{poppurpleL}{poppurple!12!white}
\colorlet{poppinkL}{poppink!12!white}
\colorlet{popgoldL}{popgold!14!white}

\pagecolor{popcream}

% ============================================================
% Typographie
% ============================================================
\defaultfontfeatures{Ligatures=TeX}
\setmainfont{PlusJakartaSans-Medium.ttf}[Path=../../../fonts/,BoldFont=PlusJakartaSans-Bold.ttf]
% Maths en sans-serif (Fira Math) avec chiffres et lettres droites de Plus
% Jakarta, pour que les formules se fondent dans le texte.
\usepackage[math-style=ISO,bold-style=ISO]{unicode-math}
\setmathfont{FiraMath-Regular.otf}[Path=../../../fonts/]
\setmathfont{PlusJakartaSans-Medium.ttf}[Path=../../../fonts/,range={up/num,up/{Latin,latin},\mathup}]
\usepackage{icomma} % virgule décimale sans espace en mode math
% Caractères mathématiques Unicode absents des polices de texte (Plus Jakarta,
% Archivo Black) : rendus par leur équivalent mathématique, en texte comme en
% formule — sinon ils s'affichent en carré vide (ex. « Arithmétique dans ℤ »).
\usepackage{newunicodechar}
\newunicodechar{ℝ}{\ensuremath{\mathbb{R}}}
\newunicodechar{ℤ}{\ensuremath{\mathbb{Z}}}
\newunicodechar{ℂ}{\ensuremath{\mathbb{C}}}
\newunicodechar{ℕ}{\ensuremath{\mathbb{N}}}
\newunicodechar{ℚ}{\ensuremath{\mathbb{Q}}}
\newunicodechar{𝔻}{\ensuremath{\mathbb{D}}}
\newunicodechar{α}{\ensuremath{\alpha}}
\newunicodechar{β}{\ensuremath{\beta}}
\newunicodechar{γ}{\ensuremath{\gamma}}
\newunicodechar{δ}{\ensuremath{\delta}}
\newunicodechar{ε}{\ensuremath{\varepsilon}}
\newunicodechar{θ}{\ensuremath{\theta}}
\newunicodechar{λ}{\ensuremath{\lambda}}
\newunicodechar{μ}{\ensuremath{\mu}}
\newunicodechar{σ}{\ensuremath{\sigma}}
\newunicodechar{τ}{\ensuremath{\tau}}
\newunicodechar{φ}{\ensuremath{\varphi}}
\newunicodechar{ψ}{\ensuremath{\psi}}
\newunicodechar{ω}{\ensuremath{\omega}}
\newunicodechar{Σ}{\ensuremath{\Sigma}}
% majuscules produites par \MakeUppercase dans les titres (α-aminés -> Α)
\newunicodechar{Α}{\ensuremath{\alpha}}
\newunicodechar{Β}{\ensuremath{\beta}}
\newunicodechar{Γ}{\ensuremath{\Gamma}}
\newunicodechar{Δ}{\ensuremath{\Delta}}
\newunicodechar{Ε}{\ensuremath{\varepsilon}}
\newunicodechar{Θ}{\ensuremath{\Theta}}
\newunicodechar{Λ}{\ensuremath{\Lambda}}
\newunicodechar{Μ}{\ensuremath{\mu}}
\newunicodechar{Τ}{\ensuremath{\tau}}
\newunicodechar{Φ}{\ensuremath{\Phi}}
\newunicodechar{Ψ}{\ensuremath{\Psi}}
\newunicodechar{Ω}{\ensuremath{\Omega}}
\newunicodechar{↦}{\ensuremath{\mapsto}}
\newunicodechar{∈}{\ensuremath{\in}}
\newunicodechar{∉}{\ensuremath{\notin}}
\newunicodechar{∧}{\ensuremath{\wedge}}
\newunicodechar{⇔}{\ensuremath{\Leftrightarrow}}
\newunicodechar{⟺}{\ensuremath{\Longleftrightarrow}}
\newunicodechar{⇒}{\ensuremath{\Rightarrow}}
\newunicodechar{⟹}{\ensuremath{\Longrightarrow}}
\newunicodechar{∘}{\ensuremath{\circ}}
\newunicodechar{∪}{\ensuremath{\cup}}
\newunicodechar{∩}{\ensuremath{\cap}}
\newunicodechar{≡}{\ensuremath{\equiv}}
\newunicodechar{⊂}{\ensuremath{\subset}}
\newunicodechar{⊥}{\ensuremath{\perp}}
\newunicodechar{∃}{\ensuremath{\exists}}
\newunicodechar{∀}{\ensuremath{\forall}}
\newunicodechar{∛}{\ensuremath{\sqrt[3]{\,}}}
\newunicodechar{∜}{\ensuremath{\sqrt[4]{\,}}}
\newunicodechar{⃗}{\ensuremath{{}^{\to}}}
\newunicodechar{ⁿ}{\ifmmode^{n}\else\textsuperscript{\ensuremath{n}}\fi}
\newunicodechar{ˣ}{\ifmmode^{x}\else\textsuperscript{\ensuremath{x}}\fi}
\newunicodechar{ᵇ}{\ifmmode^{b}\else\textsuperscript{\ensuremath{b}}\fi}
\newunicodechar{ʳ}{\ifmmode^{r}\else\textsuperscript{\ensuremath{r}}\fi}
\newunicodechar{ᵘ}{\ifmmode^{u}\else\textsuperscript{\ensuremath{u}}\fi}
\newunicodechar{ʸ}{\ifmmode^{y}\else\textsuperscript{\ensuremath{y}}\fi}
\newunicodechar{ᵃ}{\ifmmode^{a}\else\textsuperscript{\ensuremath{a}}\fi}
\newunicodechar{ᵉ}{\ifmmode^{e}\else\textsuperscript{\ensuremath{e}}\fi}
\newunicodechar{ᵏ}{\ifmmode^{k}\else\textsuperscript{\ensuremath{k}}\fi}
\newunicodechar{ᵖ}{\ifmmode^{p}\else\textsuperscript{\ensuremath{p}}\fi}
\newunicodechar{ᴺ}{\ifmmode^{N}\else\textsuperscript{\ensuremath{N}}\fi}
\newunicodechar{ᶜ}{\ifmmode^{c}\else\textsuperscript{\ensuremath{c}}\fi}
\newunicodechar{ʰ}{\ifmmode^{h}\else\textsuperscript{\ensuremath{h}}\fi}
\newunicodechar{ᵠ}{\ifmmode^{q}\else\textsuperscript{\ensuremath{q}}\fi}
\newunicodechar{ₙ}{\ifmmode_{n}\else\textsubscript{\ensuremath{n}}\fi}
\newunicodechar{ₐ}{\ifmmode_{a}\else\textsubscript{\ensuremath{a}}\fi}
\newunicodechar{ᵢ}{\ifmmode_{i}\else\textsubscript{\ensuremath{i}}\fi}
\newunicodechar{ᵣ}{\ifmmode_{r}\else\textsubscript{\ensuremath{r}}\fi}
\newunicodechar{ₖ}{\ifmmode_{k}\else\textsubscript{\ensuremath{k}}\fi}
\newunicodechar{ᵦ}{\ifmmode_{\beta}\else\textsubscript{\ensuremath{\beta}}\fi}
\newunicodechar{ₓ}{\ifmmode_{x}\else\textsubscript{\ensuremath{x}}\fi}
\newfontfamily\popdisplay{ArchivoBlack-Regular.ttf}[Path=../../../fonts/]
\newfontfamily\poplogo{SpaceGrotesk-Bold.ttf}[Path=../../../fonts/]
\newfontfamily\popsemi{PlusJakartaSans-SemiBold.ttf}[Path=../../../fonts/]
\newfontfamily\popxbold{PlusJakartaSans-ExtraBold.ttf}[Path=../../../fonts/]
\newfontfamily\popxboldls{PlusJakartaSans-ExtraBold.ttf}[Path=../../../fonts/,LetterSpace=3.0]

\setlist[enumerate]{leftmargin=1.7em,itemsep=5pt,topsep=4pt}
\setlist[itemize]{leftmargin=1.7em,itemsep=4pt,topsep=4pt,label={\color{poporange}\textbullet}}
\setlength{\parindent}{0pt}
\setlength{\parskip}{5pt}
\renewcommand{\arraystretch}{1.25}

% pgfplots : style Pop par défaut
\pgfplotsset{
  every axis/.append style={axis line style={popink,line width=1pt},tick style={popink},
    grid style={popline,line width=0.5pt},label style={font=\small},tick label style={font=\footnotesize}},
  popcurve/.style={poporange,line width=1.8pt,no markers,smooth},
}
% Même style disponible dans un \draw TikZ simple (hors pgfplots)
\tikzset{popcurve/.style={poporange,line width=1.8pt}}
\tikzset{popgrid/.style={popline,very thin}}
\colorlet{popcurve}{poporange} % le nom du style est parfois utilisé comme couleur

% ============================================================
% Pill / squiggle
% ============================================================
\newcommand{\poppill}[3]{%
  \tikz[baseline=-0.55ex]\node[draw=popink,line width=1.2pt,rounded corners=2.6pt,%
    fill=#2,inner xsep=6.5pt,inner ysep=3pt]{%
    \popxboldls\fontsize{7}{7}\selectfont\color{#3}\MakeUppercase{#1}};%
}
\newcommand{\popsquiggle}[1]{%
  \tikz[baseline=-0.4ex]{
    \draw[#1,line width=2.6pt,line cap=round]
      plot [smooth] coordinates {(0,0.09) (0.34,0.01) (0.68,0.21) (1.02,0.01) (1.36,0.10)};
  }%
}
% Mot-clé surligné (marqueur orange sous le texte)
\newcommand{\cle}[1]{{\popxbold\color{popdark}#1}}

% ============================================================
% En-tête / pied
% ============================================================
\pagestyle{fancy}
\fancyhf{}
\renewcommand{\headrulewidth}{0pt}
\renewcommand{\footrulewidth}{0pt}
\lhead{\raisebox{-0.15ex}{\poplogo\fontsize{13}{13}\selectfont\color{popink}OnBuch\color{poporange}+}}
\rhead{\poppill{\DOCMATIERE\ \textbullet\ \DOCNIVEAU\ \textbullet\ Leçon \DOCLECON}{white}{popink}}
\lfoot{\popxbold\fontsize{7.6}{7.6}\selectfont\color{popmuted}\MakeUppercase{Cours signé OnBuch+}\ \textbullet\ {\popsemi\DOCTITRE}}
\rfoot{\tikz[baseline=-0.6ex]\node[rounded corners=8.5pt,fill=popink,minimum height=17pt,minimum width=17pt,inner xsep=5pt,inner sep=0]{\popxbold\fontsize{8}{8}\selectfont\color{popcream}\thepage\,/\,\pageref*{LastPage}};}

% ============================================================
% Titres
% ============================================================
\newcommand{\chaptitre}[2]{%
  \begin{tcolorbox}[enhanced,colback=poporange,colframe=popink,boxrule=2.6pt,arc=11pt,
      drop shadow=popink,left=15pt,right=15pt,top=12pt,bottom=11pt,before skip=3pt,after skip=18pt]
    \poppill{#2}{popink}{popcream}\par\vspace{9pt}
    {\raggedright\hyphenpenalty=10000\exhyphenpenalty=10000\popdisplay\fontsize{22}{24}\selectfont\color{popcream}\MakeUppercase{#1}\par}
  \end{tcolorbox}
}
% Section numérotée : pastille orange avec le numéro + titre Archivo
\newcounter{popsec}
\newcommand{\coursec}[1]{%
  \stepcounter{popsec}\setcounter{popsub}{0}%
  \par\vspace{16pt}\noindent
  \begin{minipage}{\linewidth}
  \tikz[baseline=-0.7ex]\node[draw=popink,line width=1.6pt,fill=poporange,rounded corners=4pt,minimum size=22pt,inner sep=0]{\popdisplay\fontsize{12}{12}\selectfont\color{popcream}\thepopsec};%
  \hspace{8pt}{\popdisplay\fontsize{15}{17}\selectfont\color{popink}\MakeUppercase{#1}}\par
  \vspace{4pt}\noindent{\color{poporange}\rule{36pt}{3.6pt}}
  \end{minipage}\par\nopagebreak\vspace{8pt}
}
\newcounter{popsub}
\newcommand{\courssub}[1]{%
  \stepcounter{popsub}%
  \par\vspace{9pt}\noindent{\popxbold\fontsize{11.5}{13}\selectfont\color{popink}{\color{poporange}\thepopsec.\thepopsub}\hspace{6pt}#1}\par\nopagebreak\vspace{4pt}
}

% ============================================================
% Boîtes pédagogiques — même recette : fond blanc, bordure épaisse colorée,
% ombre dure encre, badge plein. Titre optionnel [#1].
% ============================================================
\tcbset{popbase/.style={breakable,enhanced,colback=white,boxrule=2.3pt,arc=9pt,
  shadow={3pt}{-3pt}{0pt}{popink},left=14pt,right=14pt,top=11pt,bottom=11pt,before skip=11pt,after skip=15pt,
  fontupper=\popsemi\fontsize{10.6}{14.5}\selectfont\color{popink},
  fontlower=\popsemi\fontsize{10.6}{14.5}\selectfont\color{popink}}}
% Coche dessinée (le glyphe ✓ n'existe pas dans les polices du document)
\newcommand{\popcheck}[1]{\tikz[baseline=-0.5ex]\draw[#1,line width=1.6pt,line cap=round,line join=round](-0.13,0.0)--(-0.04,-0.09)--(0.13,0.1);}
\providecommand{\coche}{\popcheck{popgreen}}
\providecommand{\resultat}[1]{\par\smallskip\noindent\fcolorbox{popgreen}{popgreenL}{\parbox{\dimexpr\linewidth-2\fboxsep-2\fboxrule\relax}{\textbf{Résultat.} #1}}\par\smallskip}
\newcommand{\popboxhead}[3]{\poppill{#1}{#2}{white}\ifx\relax#3\relax\else\hspace{6pt}{\popxbold\fontsize{10.6}{12}\selectfont\color{popink}#3}\fi\par\vspace{7pt}}

\newtcolorbox{aretenir}[1][]{popbase,colframe=popgreen,before upper={\popboxhead{À retenir}{popgreen}{#1}}}
\newtcolorbox{exemplebox}[1][]{popbase,colframe=poporange,before upper={\popboxhead{Exemple}{poporange}{#1}}}
\newtcolorbox{definition}[1][]{popbase,colframe=popblue,before upper={\popboxhead{Définition}{popblue}{#1}}}
\newtcolorbox{propriete}[1][]{popbase,colframe=poppurple,before upper={\popboxhead{Propriété}{poppurple}{#1}}}
\newtcolorbox{methode}[1][]{popbase,colframe=popgold,before upper={\popboxhead{Méthode}{popgold}{#1}}}
\newtcolorbox{attention}[1][]{popbase,colframe=poppink,before upper={\popboxhead{Attention}{poppink}{#1}}}
\newtcolorbox{experience}[1][]{popbase,colframe=popblue,colback=popblueL,before upper={\popboxhead{Expérience}{popblue}{#1}}}
% « Le savais-tu ? » : carte violette pleine (contexte, culture, Cameroun)
\newtcolorbox{savaistu}[1][]{popbase,colback=poppurple,colframe=popink,
  fontupper=\popsemi\fontsize{10.4}{14.2}\selectfont\color{white},
  before upper={\poppill{Le savais-tu\,?}{white}{poppurple}\ifx\relax#1\relax\else\hspace{6pt}{\popxbold\color{white}#1}\fi\par\vspace{7pt}}}
% Exercice résolu : énoncé en haut, \tcblower puis la solution sur fond crème
\newcounter{popexo}\newcounter{popexr}
\newtcolorbox{exoresolu}[1][]{popbase,colframe=poporange,colbacklower=popcream,
  segmentation style={popink,line width=1.2pt,dashed},
  before upper={\stepcounter{popexr}\popboxhead{Exercice résolu \thepopexr}{poporange}{#1}},
  before lower={\poppill{Solution}{popink}{popcream}\par\vspace{6pt}}}
% Exercice (non résolu en place) — la correction est dans la partie Corrigés
\newtcolorbox{exercice}[1][]{popbase,colframe=popink,
  before upper={\stepcounter{popexo}\popboxhead{Exercice \thepopexo}{popink}{#1}}}
% Corrigé
\newtcolorbox{corrige}[1][]{popbase,colframe=popgreen,colback=popgreenL,
  before upper={\popboxhead{Corrigé}{popgreen}{#1}}}
% Objectifs / prérequis / bilan
\newtcolorbox{objectifs}{popbase,colframe=popink,colback=poporangeL,
  before upper={\popboxhead{Objectifs de la leçon}{popink}{}}}
\newtcolorbox{prerequis}{popbase,colframe=popink,colback=popblueL,
  before upper={\popboxhead{Prérequis}{popblue}{}}}
\newtcolorbox{bilan}{popbase,colframe=popink,colback=white,boxrule=2.8pt,
  before upper={\popboxhead{Fiche bilan}{poporange}{L'essentiel à connaître pour l'examen}}}
% Figure encadrée avec légende
\newtcolorbox{popfigure}[1][]{enhanced,breakable=false,colback=white,colframe=popline,boxrule=1.4pt,arc=8pt,
  left=10pt,right=10pt,top=10pt,bottom=8pt,before skip=10pt,after skip=12pt,halign=center,
  lower separated=false,halign lower=center,
  fontlower=\popsemi\fontsize{9}{11.5}\selectfont\color{popmuted},
  #1}
\newcommand{\legende}[1]{\par\vspace{6pt}{\popsemi\fontsize{9}{11.5}\selectfont\color{popmuted}#1}}

% ============================================================
% Signature OnBuch+
% ============================================================
\newcommand{\poplogoinline}[1]{{\poplogo\fontsize{#1}{#1}\selectfont\color{popink}OnBuch\color{poporange}+}}
\newcommand{\signatureonbuch}{%
  \begin{tcolorbox}[enhanced,colback=popink,colframe=popink,boxrule=2.2pt,arc=10pt,
      drop shadow=poporange,left=14pt,right=14pt,top=10pt,bottom=10pt,before skip=0pt,after skip=0pt]
    \tikz[baseline=-0.7ex]\node[circle,fill=popgreen,draw=popcream,line width=1.3pt,minimum size=22pt,inner sep=0]{\popcheck{white}};%
    \hspace{9pt}\begin{minipage}[c]{0.62\linewidth}
      {\popxbold\fontsize{12}{13}\selectfont\color{popcream}Signé\ \textbullet\ {\poplogo OnBuch\color{poporange}+}}\par\vspace{2pt}
      {\raggedright\popsemi\fontsize{8}{10.5}\selectfont\color{popline}\MakeUppercase{Équipe pédagogique OnBuch+}\\\MakeUppercase{Conforme au programme officiel MINESEC}\par}
    \end{minipage}\hfill
    {\poplogo\fontsize{22}{22}\selectfont\color{popcream}OnBuch\color{poporange}+}
  \end{tcolorbox}%
}

% ============================================================
% Page de garde
% #1 titre · #2 matière · #3 niveau · #4 module · #5 n° leçon · #6 durée ·
% #7 description / objectifs (texte ou itemize)
% ============================================================
\newcommand{\pagegarde}[7]{%
  \thispagestyle{empty}
  \newgeometry{a4paper,margin=1.7cm,top=1.6cm,bottom=1.4cm}%
  \begin{tikzpicture}[remember picture,overlay]
    \fill[poporange!18] ([xshift=-3.2cm,yshift=-3.6cm]current page.north east) circle (4.6cm);
    \fill[poppurple!14] ([xshift=2.4cm,yshift=7.6cm]current page.south west) circle (3.8cm);
  \end{tikzpicture}%
  {\poplogo\fontsize{30}{30}\selectfont\color{popink}OnBuch\color{poporange}+}\hfill
  \raisebox{0.6ex}{\poppill{Cours signé \textbullet\ Premium}{popink}{popcream}}\par
  \vspace{1pt}\popsquiggle{poppurple}\par
  \vspace{8pt}
  \begin{tcolorbox}[enhanced,colback=poporange,colframe=popink,boxrule=2.8pt,arc=13pt,
      drop shadow=popink,left=18pt,right=18pt,top=12pt,bottom=13pt,after skip=10pt]
    \poppill{#2 \textbullet\ #3}{popink}{popcream}\hspace{5pt}\poppill{#4}{white}{popink}\par\vspace{8pt}
    {\popdisplay\fontsize{14}{14}\selectfont\color{popink}LEÇON #5}\par\vspace{4pt}
    {\raggedright\hyphenpenalty=10000\exhyphenpenalty=10000\popdisplay\fontsize{25}{27}\selectfont\color{popcream}\MakeUppercase{#1}\par}
    \vspace{8pt}
    {\popxbold\fontsize{9.5}{9.5}\selectfont\color{popink}\MakeUppercase{Durée conseillée : #6}}
  \end{tcolorbox}
  \begin{tcolorbox}[enhanced,colback=white,colframe=popink,boxrule=2.2pt,arc=10pt,
      drop shadow=popink,left=16pt,right=16pt,top=10pt,bottom=10pt,after skip=8pt]
    \poppill{Dans cette leçon}{poppurple}{white}\par\vspace{5pt}
    \popsemi\fontsize{9.3}{12.6}\selectfont\color{popink}#7
  \end{tcolorbox}
  \noindent\begin{minipage}[t]{0.487\linewidth}
    \begin{tcolorbox}[enhanced,colback=popgreen,colframe=popink,boxrule=2.2pt,arc=10pt,
        drop shadow=popink,left=11pt,right=11pt,top=10pt,bottom=10pt,height=3.1cm,valign=center]
      \tcbox[colback=white,colframe=popink,boxrule=1.3pt,arc=5pt,boxsep=0pt,nobeforeafter,
        left=3pt,right=3pt,top=3pt,bottom=3pt,valign=center]{\qrcode[height=1.9cm]{https://play.google.com/store/apps/details?id=com.luvvix.onbuch&hl=fr}}%
      \hspace{8pt}\begin{minipage}[c]{0.5\linewidth}
        {\popdisplay\fontsize{11}{12.5}\selectfont\color{white}\MakeUppercase{Télécharge OnBuch}}\par\vspace{4pt}
        {\raggedright\popsemi\fontsize{8.4}{11}\selectfont\color{white}Gratuit \textbullet\ Google Play\\Tuteur IA, annales, concours, résultats\par}
      \end{minipage}
    \end{tcolorbox}
  \end{minipage}\hfill
  \begin{minipage}[t]{0.487\linewidth}
    \begin{tcolorbox}[enhanced,colback=poppurple,colframe=popink,boxrule=2.2pt,arc=10pt,
        drop shadow=popink,left=11pt,right=11pt,top=10pt,bottom=10pt,height=3.1cm,valign=center]
      \tikz[baseline=-0.7ex]\node[circle,fill=white,draw=popink,line width=1.3pt,minimum size=1.6cm,inner sep=0]{\popdisplay\fontsize{24}{24}\selectfont\color{poppurple}+};%
      \hspace{8pt}\begin{minipage}[c]{0.6\linewidth}
        {\popdisplay\fontsize{11}{12.5}\selectfont\color{white}\MakeUppercase{Passe à OnBuch+}}\par\vspace{4pt}
        {\raggedright\popsemi\fontsize{8.4}{11}\selectfont\color{white}Tous les cours signés, Tuteur IA illimité, fiches PDF — onglet OnBuch+ de l'app\par}
      \end{minipage}
    \end{tcolorbox}
  \end{minipage}
  \vspace{8pt}
  \popcontenu
  \vfill
  \signatureonbuch
  \restoregeometry
  \newpage
}

% Rangée « Ce cours contient » (page de garde)
\newcommand{\poptile}[3]{% #1 couleur #2 gros texte #3 légende
  \begin{tcolorbox}[enhanced,colback=#1,colframe=popink,boxrule=2pt,arc=9pt,shadow={2.5pt}{-2.5pt}{0pt}{popink},
      left=6pt,right=6pt,top=9pt,bottom=9pt,halign=center,valign=center,height=1.75cm,width=\linewidth,nobeforeafter]
    {\popdisplay\fontsize{12.5}{13}\selectfont\color{white}\MakeUppercase{#2}}\par\vspace{3pt}
    {\centering\popsemi\fontsize{7.8}{9.5}\selectfont\color{white}#3\par}
  \end{tcolorbox}}
\newcommand{\popcontenu}{%
  \par\noindent{\popxboldls\fontsize{7.5}{7.5}\selectfont\color{popmuted}CE COURS SIGNÉ CONTIENT}\par\vspace{7pt}
  \noindent\begin{minipage}[t]{0.235\linewidth}\poptile{popblue}{Cours}{complet, illustré, pas à pas}\end{minipage}\hfill
  \begin{minipage}[t]{0.235\linewidth}\poptile{poporange}{Méthodes}{et exercices résolus}\end{minipage}\hfill
  \begin{minipage}[t]{0.235\linewidth}\poptile{poppink}{Exercices}{type examen + corrigés}\end{minipage}\hfill
  \begin{minipage}[t]{0.235\linewidth}\poptile{popgreen}{Fiche bilan}{l'essentiel à retenir}\end{minipage}\par}

% Sommaire
\newtcolorbox{sommaire}{enhanced,colback=white,colframe=popink,boxrule=2.2pt,arc=10pt,
  drop shadow=popink,left=18pt,right=18pt,top=13pt,bottom=13pt,before skip=6pt,after skip=20pt,
  before upper={\poppill{Sommaire}{poppurple}{white}\par\vspace{9pt}\popsemi\fontsize{10.6}{14.5}\selectfont\color{popink}}}

% Signature de fin de cours — poussée tout en bas de la dernière page
\newcommand{\findecours}{%
  \par\vfill\nopagebreak
  \begin{center}\popsquiggle{poporange}\end{center}\vspace{4pt}
  \signatureonbuch
}
\AtBeginDocument{\def\llbracket{\mathopen{[\mkern-3.5mu[}}\def\rrbracket{\mathclose{]\mkern-3.5mu]}}} % intervalles d'entiers (glyphe ⟦ absent de la police)
% Glyphes absents de Fira Math : redéfinis à partir de glyphes disponibles
\AtBeginDocument{%
  \def\setminus{\mathbin{\backslash}}\let\smallsetminus\setminus
  \def\vdots{\mathord{\vbox{\baselineskip3.2pt\lineskiplimit0pt\kern4pt\hbox{.}\hbox{.}\hbox{.}}}}%
  \def\ddots{\mathinner{\mkern1mu\raise6.5pt\hbox{.}\mkern2mu\raise3.6pt\hbox{.}\mkern2mu\raise0.7pt\hbox{.}\mkern1mu}}%
  \def\triangle{\mathord{\scalebox{0.9}{$\symup{\Delta}$}}}%
  \def\nabla{\mathord{\raisebox{\depth}{\scalebox{1}[-1]{$\symup{\Delta}$}}}}%
  \def\checkmark{\popcheck{popdark}}%
  \def\star{\mathbin{\ast}}\def\bigstar{\mathord{\ast}}%
  \def\diamond{\mathbin{\scalebox{0.6}{$\Diamond$}}}%
  \def\bigcup{\mathop{\mathchoice{\raisebox{-0.3em}{\scalebox{1.7}{$\cup$}}}{\scalebox{1.25}{$\cup$}}{\cup}{\cup}}\displaylimits}%
  \def\bigcap{\mathop{\mathchoice{\raisebox{-0.3em}{\scalebox{1.7}{$\cap$}}}{\scalebox{1.25}{$\cap$}}{\cap}{\cap}}\displaylimits}%
  \def\lesseqqgtr{\mathrel{\lessgtr}}\def\bigoplus{\mathop{\oplus}\displaylimits}%
}
\providecommand{\ssi}{si et seulement si} % abréviation fréquente
\AtBeginDocument{\providecommand{\wideparen}{\overparen}} % arc de cercle

%% FIGFIX-BEGIN v4 — correctifs globaux des figures (docs/onbuch-generation/tools/figfix)
\usepackage{adjustbox}\usepackage{etoolbox}
% toute figure TikZ/pgfplots plus large que la ligne est réduite à la largeur disponible
\BeforeBeginEnvironment{tikzpicture}{\begin{adjustbox}{max width=\linewidth}}
\AfterEndEnvironment{tikzpicture}{\end{adjustbox}}
% légendes pgfplots lisibles par-dessus les courbes
\pgfplotsset{every axis/.append style={legend style={fill=white,fill opacity=0.92,text opacity=1,draw=black!15,inner sep=3pt}}}
% étiquettes d'arêtes (syntaxe edge["..."]) avec fond blanc
\tikzset{every edge quotes/.append style={fill=white,inner sep=1.2pt}}
%% FIGFIX-END
\begin{document}

\pagegarde{Éradication des préjugés autour de l'apparition des anomalies et/ou des caractères nouveaux au sein des familles}{SVT}{Tle D}{Module I : Le Monde Vivant — Origine de nouveaux allèles}{N05}{8 h}{Cette leçon explique l'origine scientifique des caractères nouveaux et des anomalies observés au sein des familles : mutations géniques, mutations chromosomiques et brassage des allèles. Elle fournit les outils pour identifier ces anomalies à partir de documents (séquences d'ADN, caryotypes) et pour combattre les préjugés familiaux et sociaux.
\vspace{4pt}
\begin{multicols}{2}\small
\begin{itemize}
\item À la fin de cette leçon, je sais définir une mutation et distinguer mutation génique et mutation chromosomique.
\item À la fin de cette leçon, je sais identifier les types de mutations géniques (substitution, insertion, délétion) et prévoir leurs conséquences sur la protéine.
\item À la fin de cette leçon, je sais identifier une anomalie génique (drépanocytose, hémophilie) par analyse comparée de séquences d'ADN.
\item À la fin de cette leçon, je sais identifier une anomalie chromosomique (trisomie 21, syndromes de Turner et de Klinefelter) à partir d'un caryotype.
\item À la fin de cette leçon, je sais expliquer le mécanisme de la non-disjonction chromosomique à l'origine des trisomies et monosomies.
\item À la fin de cette leçon, je sais expliquer comment le brassage des allèles (méiose et fécondation) assure la diversité génétique et l'unicité des individus.
\end{itemize}\end{multicols}}

\begin{sommaire}
\begin{multicols}{2}
\begin{enumerate}
\item I. Les mutations : origine des allèles nouveaux
\item II. Anomalies géniques : exemple de la drépanocytose
\item III. Autres anomalies géniques : hémophilie et aperçus
\item IV. Anomalies chromosomiques : trisomie 21, Turner, Klinefelter
\item V. Brassage des allèles : diversité génétique et unicité des individus
\item VI. Éradication des préjugés : éducation et sensibilisation
\item Activité d'intégration
\item Méthodes et exercices résolus
\item Exercices
\item Corrigés des exercices
\item Fiche bilan
\end{enumerate}
\end{multicols}
\end{sommaire}

\begin{objectifs}
\begin{itemize}
\item À la fin de cette leçon, je sais définir une mutation et distinguer mutation génique et mutation chromosomique.
\item À la fin de cette leçon, je sais identifier les types de mutations géniques (substitution, insertion, délétion) et prévoir leurs conséquences sur la protéine.
\item À la fin de cette leçon, je sais identifier une anomalie génique (drépanocytose, hémophilie) par analyse comparée de séquences d'ADN.
\item À la fin de cette leçon, je sais identifier une anomalie chromosomique (trisomie 21, syndromes de Turner et de Klinefelter) à partir d'un caryotype.
\item À la fin de cette leçon, je sais expliquer le mécanisme de la non-disjonction chromosomique à l'origine des trisomies et monosomies.
\item À la fin de cette leçon, je sais expliquer comment le brassage des allèles (méiose et fécondation) assure la diversité génétique et l'unicité des individus.
\item À la fin de cette leçon, je sais argumenter scientifiquement pour combattre les préjugés familiaux et sociaux sur l'origine des anomalies.
\end{itemize}
\end{objectifs}

\begin{prerequis}
\begin{itemize}
\item Structure de l'ADN, gène, allèle et code génétique (Première)
\item Chromosomes, caryotype et formule chromosomique humaine (2n = 46)
\item Méiose et fécondation : production de gamètes haploïdes
\item Notions de dominance, récessivité et transmission des caractères héréditaires
\item Relation gène → protéine → caractère
\end{itemize}
\end{prerequis}

\newpage

\coursec{I. Les mutations : origine des allèles nouveaux}

Dans un village de la région du Centre, une famille accueille un enfant atteint de drépanocytose : certains proches accusent la mère de sorcellerie ou d'infidélité, tandis qu'un couple voisin, dont le bébé présente une trisomie 21, est rejeté par la belle-famille. Pourtant, ni les parents ni les grands-parents ne présentent ces caractères. D'où viennent alors ces \cle{anomalies géniques} et ces \cle{anomalies chromosomiques} qui apparaissent parfois au sein des familles ? Sont-ils le fruit d'une malédiction, d'une faute, ou obéissent-ils à des mécanismes biologiques précis et explicables ? Cette leçon apporte les réponses scientifiques pour éradiquer ces préjugés.

\textbf{Problématique :} comment des caractères nouveaux, parfois anormaux, peuvent-ils apparaître dans une lignée où ils n'existaient pas ? Quelle est leur origine réelle, et pourquoi ne peut-on en accuser personne ?

\courssub{Constat : l'apparition spontanée de caractères nouveaux dans les lignées}

L'observation attentive des populations animales et végétales montre que des individus présentant des caractères inédits apparaissent parfois spontanément, sans que leurs parents possèdent ces caractères.

\begin{exemplebox}[Le mouton Ancon : un caractère nouveau apparu spontanément]
En 1791, dans le troupeau du fermier américain Seth Wright, naît un agneau mâle aux pattes anormalement courtes, alors que tous les autres moutons du troupeau sont normaux. Ce caractère « pattes courtes » est transmis à la descendance de cet agneau : il donne naissance à une race entière, le mouton \emph{Ancon}, dont les pattes courtes empêchent les sauts par-dessus les clôtures — un avantage pour l'éleveur. Un caractère totalement nouveau est donc apparu \emph{spontanément} chez un individu, puis s'est révélé \emph{héréditaire}.
\end{exemplebox}

\begin{exemplebox}[Les mutants de la drosophile observés par Morgan]
Au début du \textsc{xx}\textsuperscript{e} siècle, Thomas Hunt Morgan élève des milliers de drosophiles (mouches du vinaigre) au corps gris et aux yeux rouges. En 1910, il observe soudain, dans une souche pure, un mâle aux \emph{yeux blancs}. Ce caractère, absent de tous les ancêtres, se transmet ensuite selon des proportions prévisibles : il correspond à un nouvel allèle du gène de la couleur des yeux, localisé sur le chromosome $X$. Par la suite, Morgan et son équipe décrivent des dizaines de caractères nouveaux apparus de la même façon (ailes atrophiées, corps ébène, yeux en barre\ldots).
\end{exemplebox}

Ces observations conduisent à une question fondatrice : \emph{quel mécanisme biologique peut créer, de façon stable et transmissible, un caractère nouveau ?} La réponse est la \cle{mutation}.

\courssub{Définition de la mutation}

\begin{definition}[Mutation]
Une \cle{mutation} est une modification \textbf{stable} et \textbf{héréditaire} du matériel génétique. On distingue :
\begin{itemize}
\item les \cle{mutations géniques}, qui modifient la séquence de nucléotides d'un gène (sur l'ADN) ;
\item les \cle{mutations chromosomiques}, qui modifient la \textbf{structure} d'un chromosome ou le \textbf{nombre} de chromosomes.
\end{itemize}
La version nouvelle du gène créée par mutation constitue un \cle{nouvel allèle}. La mutation est donc la \textbf{source ultime de la diversité des allèles}.
\end{definition}

Deux précisions essentielles découlent de cette définition :

\begin{itemize}
\item la modification est \textbf{stable} : elle persiste au cours des divisions cellulaires qui suivent son apparition, car l'ADN muté est répliqué tel quel ;
\item elle est \textbf{héréditaire} au niveau cellulaire : toute cellule fille issue de la division d'une cellule mutée porte la mutation. En revanche, sa transmission à la \emph{descendance} de l'individu dépend de la localisation de la cellule mutée (voir ci-dessous).
\end{itemize}

\courssub{Mutations somatiques et mutations germinales}

Le corps d'un organisme pluricellulaire comporte deux grandes catégories de cellules : les \cle{cellules somatiques} (cellules du corps : peau, muscle, foie\ldots) et les \cle{cellules germinales} (lignée produisant les gamètes). Le destin d'une mutation dépend entièrement de la catégorie de cellule qu'elle affecte.

\begin{propriete}[Transmission des mutations — admise]
Seules les mutations affectant les \textbf{cellules germinales} (gamètes ou leurs précurseurs) sont \textbf{transmissibles à la descendance}. Une mutation survenue dans une cellule \textbf{somatique} ne concerne que l'individu lui-même (elle peut provoquer localement un dysfonctionnement, par exemple un cancer) et \textbf{disparaît avec lui} : elle n'est jamais transmise à ses enfants.
\end{propriete}

\begin{exemplebox}[Illustration]
Un coup de soleil intense provoque une mutation dans une cellule de la peau : cette mutation somatique peut dégénérer en cancer de la peau chez la personne exposée, mais ses enfants ne naîtront pas avec cette mutation. En revanche, si une mutation survient dans une cellule de la lignée germinale d'un parent, le gamète correspondant la porte, et l'enfant issu de ce gamète possédera la mutation dans \emph{toutes} ses cellules.
\end{exemplebox}

\begin{attention}[Piège classique]
Ne pas confondre « héréditaire » au sens cellulaire (transmise aux cellules filles lors des mitoses) et « transmissible à la descendance » (via les gamètes). Une mutation somatique est héréditaire pour les cellules filles, mais \textbf{non transmissible} aux générations suivantes.
\end{attention}

\courssub{Les mutations géniques : substitution, insertion, délétion}

Une \cle{mutation génique} (ou mutation ponctuelle lorsqu'elle ne concerne qu'un nucléotide) affecte la séquence d'un gène. On en distingue trois types fondamentaux :

\begin{itemize}
\item la \cle{substitution} : un nucléotide est remplacé par un autre (ex. A remplacé par T) ;
\item l'\cle{insertion} : un ou plusieurs nucléotides sont ajoutés dans la séquence ;
\item la \cle{délétion} : un ou plusieurs nucléotides sont supprimés de la séquence.
\end{itemize}

\begin{popfigure}
\begin{center}
\begin{tikzpicture}[font=\small, scale=0.9]
% Styles
\tikzset{base/.style={draw=popblue, fill=popblueL, minimum width=0.55cm, minimum height=0.55cm, font=\small\bfseries},
         mut/.style={draw=poporange, fill=poporangeL, minimum width=0.55cm, minimum height=0.55cm, font=\small\bfseries},
         del/.style={draw=poporange, fill=white, dashed, minimum width=0.55cm, minimum height=0.55cm}}
% Séquence normale
\node[anchor=west, font=\bfseries\small, popdark] at (0,3.8) {Séquence normale :};
\foreach \i/\b in {1/A,2/T,3/G,4/C,5/A,6/T,7/G,8/C,9/A}{
  \node[base] at (\i*0.65,3.2) {\b};}
% Substitution
\node[anchor=west, font=\bfseries\small, popdark] at (0,2.5) {Substitution :};
\foreach \i/\b in {1/A,2/T,3/G,4/C,5/A,6/T,7/G,8/C,9/A}{
  \node[base] at (\i*0.65,1.9) {\b};}
\node[mut] at (4*0.65,1.9) {A};
% Insertion
\node[anchor=west, font=\bfseries\small, popdark] at (0,1.2) {Insertion :};
\foreach \i/\b in {1/A,2/T,3/G}{
  \node[base] at (\i*0.65,0.6) {\b};}
\node[mut] at (4*0.65,0.6) {T};
\foreach \i/\b in {5/C,6/A,7/T,8/G,9/C,10/A}{
  \node[base] at (\i*0.65,0.6) {\b};}
% Délétion
\node[anchor=west, font=\bfseries\small, popdark] at (0,-0.1) {Délétion :};
\foreach \i/\b in {1/A,2/T,3/G}{
  \node[base] at (\i*0.65,-0.7) {\b};}
\node[del] at (4*0.65,-0.7) {};
\foreach \i/\b in {5/A,6/T,7/G,8/C,9/A}{
  \node[base] at (\i*0.65,-0.7) {\b};}
% Légende
\node[mut] at (7.5,1.9) {};
\node[anchor=west] at (7.8,1.9) {nucléotide modifié / ajouté};
\node[del] at (7.5,0.6) {};
\node[anchor=west] at (7.8,0.6) {nucléotide perdu (case vide)};
\end{tikzpicture}
\end{center}
\legende{Comparaison des trois types de mutations géniques sur une même séquence d'ADN de référence (9 paires de bases). La substitution remplace un nucléotide (position 4 : C $\to$ A) ; l'insertion en ajoute un (position 4 : T) décalant la suite ; la délétion en supprime un (position 4 : C manquant).}
\end{popfigure}

\courssub{Conséquences des mutations géniques sur la protéine}

Un gène code une protéine via le code génétique : la séquence d'ARNm est lue par \cle{codons} (triplets de nucléotides), chaque codon spécifiant un acide aminé. L'effet d'une mutation dépend donc de la façon dont elle modifie cette lecture.

\begin{itemize}
\item \textbf{Mutation silencieuse} : la substitution change un codon en un autre codon codant \emph{le même} acide aminé (le code génétique est dégénéré : plusieurs codons peuvent coder un même acide aminé). La protéine est inchangée : la mutation est \textbf{sans effet phénotypique}.
\item \textbf{Mutation faux-sens} (ou \emph{missense}) : la substitution change un codon en un codon codant un \emph{autre} acide aminé. La protéine possède un acide aminé différent à une position : son activité peut être modifiée ou abolie. C'est le cas de la drépanocytose (section II).
\item \textbf{Mutation non-sens} : la substitution transforme un codon en \emph{codon stop}. La traduction s'arrête prématurément : la protéine est \textbf{tronquée}, le plus souvent non fonctionnelle.
\item \textbf{Décalage du cadre de lecture} (\emph{frameshift}) : une insertion ou une délétion d'un nombre de nucléotides \emph{non multiple de 3} décale le regroupement en codons à partir du site de la mutation. \textbf{Tous} les acides aminés situés en aval sont modifiés, et un codon stop apparaît généralement rapidement : la protéine est profondément altérée. C'est la conséquence la plus grave.
\end{itemize}

\begin{aretenir}[Hiérarchie des effets]
Substitution silencieuse (souvent neutre) $<$ faux-sens (effet variable) $<$ non-sens (protéine tronquée) $<$ décalage de cadre (protéine bouleversée sur toute sa partie aval). Mais attention : une faux-sens peut être très grave si elle touche un acide aminé essentiel au repliement ou au site actif de la protéine.
\end{aretenir}

\begin{methode}[Identifier le type d'une mutation génique à partir de séquences]
\begin{enumerate}
\item \textbf{Aligner} la séquence mutante avec la séquence de référence, nucléotide par nucléotide.
\item \textbf{Repérer le premier nucléotide différent} entre les deux séquences.
\item \textbf{Comparer les longueurs} : si les deux séquences ont la même longueur, il s'agit d'une \emph{substitution} ; si la mutante est plus longue, d'une \emph{insertion} ; si elle est plus courte, d'une \emph{délétion}.
\item \textbf{Vérifier le cadre de lecture} : regrouper les nucléotides en codons à partir du début du gène (codon initiateur ATG). Si le nombre de nucléotides insérés ou délétés n'est pas un multiple de 3, le cadre de lecture est décalé (\emph{frameshift}) et tous les codons en aval sont modifiés.
\item \textbf{Conclure} sur la conséquence protéique : même acide aminé (silencieuse), acide aminé différent (faux-sens), codon stop (non-sens) ou séquence bouleversée (décalage de cadre).
\end{enumerate}
\end{methode}

\begin{exoresolu}[Analyse comparée de deux séquences d'ADN]
On donne un fragment du brin transcrit d'un gène normal et celui de l'allèle muté (le cadre de lecture commence au premier nucléotide) :

\textbf{Séquence normale :} \texttt{ATG GCT TAC GAA TGA}

\textbf{Séquence mutante :} \texttt{ATG GCT TAC CGA ATG A\ldots}

1. Identifier le type de mutation. 2. Préciser sa conséquence sur la protéine.
\tcblower
\textbf{1. Type de mutation.} En alignant les séquences, on constate que la mutante possède un nucléotide \textbf{C inséré} entre le 9\textsuperscript{e} et le 10\textsuperscript{e} nucléotide (après \texttt{TAC}). La séquence mutante est plus longue d'un nucléotide : c'est une \textbf{insertion}.

\textbf{2. Conséquence.} Un seul nucléotide est inséré : or 1 n'est pas un multiple de 3, donc le cadre de lecture est \textbf{décalé} à partir du site de l'insertion. Le 4\textsuperscript{e} codon passe de \texttt{GAA} à \texttt{CGA} (acide aminé différent), et tous les codons suivants sont modifiés ; de plus, le codon stop \texttt{TGA} a disparu, ce qui prolonge anormalement la traduction. Il s'agit d'une \textbf{mutation par décalage du cadre de lecture}, produisant une protéine profondément modifiée, très probablement non fonctionnelle.
\end{exoresolu}

\courssub{Les mutations chromosomiques}

Les mutations chromosomiques affectent non plus un nucléotide, mais des \textbf{fragments entiers de chromosomes} (mutations de structure) ou le \textbf{nombre} de chromosomes (mutations de nombre).

\textbf{a) Mutations de structure.} Elles résultent de cassures chromosomiques suivies de réparations incorrectes :

\begin{itemize}
\item la \cle{délétion} : perte d'un fragment de chromosome ;
\item la \cle{duplication} : un fragment est présent en deux exemplaires sur le même chromosome ;
\item l'\cle{inversion} : un fragment se détache puis se réinsère en sens inverse ;
\item la \cle{translocation} : un fragment d'un chromosome est transféré sur un autre chromosome (non homologue).
\end{itemize}

\begin{popfigure}
\begin{center}
\begin{tikzpicture}[font=\small, scale=0.85]
\tikzset{seg/.style={minimum width=0.7cm, minimum height=0.45cm, draw=popdark, font=\footnotesize\bfseries, outer sep=0pt}}
% Coordonnées de base pour les 4 types
\def\dx{0.75}
% --- DELÉTION (colonne 1) ---
\node[anchor=south west, font=\bfseries\small, popdark] at (0,4.2) {Délétion};
\foreach \i/\col/\lab in {1/popblueL/A, 2/popgreenL/B, 3/poporangeL/C, 4/poppinkL/D} {
  \node[seg, fill=\col] (d\i) at (\i*\dx, 3.6) {\lab};
}
\draw[->, thick, popmuted] (1.5*\dx, 3.2) -- (1.5*\dx, 2.8);
\foreach \i/\col/\lab in {1/popblueL/A, 2/popgreenL/B, 4/poppinkL/D} {
  \node[seg, fill=\col] at (\i*\dx - 0.75, 2.2) {\lab}; % shift left for missing C
}
% Correction position après délétion: A B D
\node[seg, fill=popblueL] at (1*\dx, 2.2) {A};
\node[seg, fill=popgreenL] at (2*\dx, 2.2) {B};
\node[seg, fill=poppinkL] at (3*\dx, 2.2) {D};

% --- DUPLICATION (colonne 2, décalage x = 5cm) ---
\def\ox{5.0}
\node[anchor=south west, font=\bfseries\small, popdark] at (\ox,4.2) {Duplication};
\foreach \i/\col/\lab in {1/popblueL/A, 2/popgreenL/B, 3/poporangeL/C} {
  \node[seg, fill=\col] at (\ox+\i*\dx, 3.6) {\lab};
}
\draw[->, thick, popmuted] (\ox+1.5*\dx, 3.2) -- (\ox+1.5*\dx, 2.8);
\foreach \i/\col/\lab in {1/popblueL/A, 2/popgreenL/B, 3/poporangeL/C, 4/poporangeL/C} {
  \node[seg, fill=\col] at (\ox+\i*\dx, 2.2) {\lab};
}

% --- INVERSION (colonne 3, décalage x = 10cm) ---
\def\oxii{10.0}
\node[anchor=south west, font=\bfseries\small, popdark] at (\oxii,4.2) {Inversion};
\foreach \i/\col/\lab in {1/popblueL/A, 2/popgreenL/B, 3/poporangeL/C, 4/poppinkL/D} {
  \node[seg, fill=\col] at (\oxii+\i*\dx, 3.6) {\lab};
}
\draw[->, thick, popmuted] (\oxii+1.5*\dx, 3.2) -- (\oxii+1.5*\dx, 2.8);
\foreach \i/\col/\lab in {1/popblueL/A, 2/poporangeL/C, 3/popgreenL/B, 4/poppinkL/D} {
  \node[seg, fill=\col] at (\oxii+\i*\dx, 2.2) {\lab};
}

% --- TRANSLOCATION (ligne du bas, deux chromosomes) ---
\node[anchor=north west, font=\bfseries\small, popdark] at (0,1.0) {Translocation (échange entre chromosomes non homologues)};
% Chromosome 1 (bleu) avant
\node[seg, fill=popblueL] at (1*\dx, 0.4) {A};
\node[seg, fill=popgreenL] at (2*\dx, 0.4) {B};
\node[seg, fill=poporangeL] at (3*\dx, 0.4) {C};
% Chromosome 2 (or) avant
\node[seg, fill=popgoldL] at (5*\dx, 0.4) {P};
\node[seg, fill=popgoldL] at (6*\dx, 0.4) {Q};
\draw[->, thick, popmuted] (3.5*\dx, 0.0) -- (3.5*\dx, -0.4);
% Après
\node[seg, fill=popblueL] at (1*\dx, -1.0) {A};
\node[seg, fill=popgreenL] at (2*\dx, -1.0) {B};
\node[seg, fill=popgoldL] at (3*\dx, -1.0) {Q};
\node[seg, fill=popgoldL] at (5*\dx, -1.0) {P};
\node[seg, fill=poporangeL] at (6*\dx, -1.0) {C};
\end{tikzpicture}
\end{center}
\legende{Les quatre mutations chromosomiques de structure. Chaque lettre colorée représente un fragment de chromosome. \textbf{Délétion} : perte du fragment C. \textbf{Duplication} : le fragment C est en double exemplaire. \textbf{Inversion} : l'ordre des fragments B et C est inversé. \textbf{Translocation} : échange de fragments (C et Q) entre deux chromosomes non homologues (bleu et doré).}
\end{popfigure}

\textbf{b) Mutations de nombre : les aneuploïdies.} Elles résultent le plus souvent d'une \cle{non-disjonction} des chromosomes (ou des chromatides) au cours de la méiose : une paire de chromosomes homologues (méiose I) ou des chromatides sœurs (méiose II) ne se sépare pas, produisant des gamètes avec un chromosome en trop ($n+1$) ou en moins ($n-1$). Après fécondation par un gamète normal ($n$), on obtient :

\begin{itemize}
\item une \cle{trisomie} : $2n+1$ chromosomes (un chromosome en trois exemplaires) — exemple : la trisomie 21 (syndrome de Down), étudiée en section IV ;
\item une \cle{monosomie} : $2n-1$ chromosomes (un chromosome en un seul exemplaire) — exemple : le syndrome de Turner ($45,X$).
\end{itemize}

\begin{attention}[Critère de distinction au Bac — Ne pas confondre]
Une \textbf{mutation génique} est invisible au microscope optique : elle touche un ou quelques nucléotides et ne modifie ni l'aspect ni le nombre des chromosomes sur un caryotype. Une \textbf{mutation chromosomique de nombre} (aneuploïdie) est visible sur un caryotype (nombre de chromosomes anormal). Une \textbf{mutation chromosomique de structure} (délétion, translocation\ldots) est souvent visible sur un caryotype à bande (changement de morphologie). C'est ce critère d'observation (caryotype vs séquençage) qui permet, à l'examen, de classer une anomalie à partir des documents fournis.
\end{attention}

\courssub{Fréquence des mutations et agents mutagènes}

Les mutations spontanées sont des événements \textbf{rares} : la fréquence de mutation d'un gène est de l'ordre de \num{1e-5} à \num{1e-6} par gène et par génération. Cette rareté s'explique par l'extrême fidélité de la réplication de l'ADN (activité de relecture des ADN polymérases) et par l'existence de mécanismes de réparation de l'ADN.

Cependant, certains facteurs de l'environnement, appelés \cle{agents mutagènes}, augmentent fortement la fréquence des mutations :

\begin{itemize}
\item les \textbf{rayonnements} : rayons ultraviolets (UV du soleil), rayons X, rayonnements ionisants (radioactivité) ;
\item les \textbf{substances chimiques} : goudrons du tabac, benzène, certains pesticides, colorants azoïques, moutarde azotée\ldots ;
\item certains \textbf{virus} oncogènes capables de s'insérer dans le génome hôte (ex. VPH, VHB).
\end{itemize}

\begin{experience}[Mise en évidence expérimentale de l'effet mutagène (Muller, 1927)]
\textbf{Protocole.} On soumet deux lots de drosophiles issues d'une même souche pure à deux conditions : lot témoin non irradié, lot expérimental exposé aux rayons X.

\textbf{Observations.} La fréquence d'apparition de caractères nouveaux (mutants visibles : yeux blancs, ailes courtes\ldots) est multipliée par plus de 100 dans le lot irradié par rapport au lot témoin.

\textbf{Interprétation.} Les rayons X endommagent l'ADN (cassures des brins, modifications de bases) ; les réparations imparfaites créent des mutations stables. Les agents mutagènes \emph{augmentent la fréquence} des mutations mais n'en déterminent \emph{ni la nature ni la localisation} : les mutations obtenues restent aléatoires.
\end{experience}

\begin{attention}[Une mutation n'est pas forcément nuisible]
Contrairement à une idée reçue, une mutation n'est pas systématiquement une « maladie ». Elle peut être \textbf{neutre} (mutation silencieuse, ou mutation dans une région non codante sans fonction régulatrice connue), \textbf{défavorable} (anomalies génétiques, cancers), ou même \textbf{avantageuse} dans un environnement donné (ex. résistance d'un moustique à un insecticide, résistance au paludisme pour les hétérozygotes drépanocytaires). La mutation est la \textbf{source ultime de la diversité des allèles} : sans mutations, aucune variabilité génétique nouvelle ne pourrait apparaître, et l'évolution serait impossible.
\end{attention}

\courssub{La mutation : un événement aléatoire et non dirigé}

\begin{aretenir}[Savoir admis — fondamental pour lutter contre les préjugés]
La mutation est un événement \textbf{aléatoire}, \textbf{non dirigé} et \textbf{indépendant de la volonté des individus}. Elle survient par accident lors de la réplication de l'ADN ou sous l'effet d'agents mutagènes, sans aucun rapport avec le comportement, la moralité ou les souhaits des parents. Aucun être humain ne choisit les mutations qu'il transmet, et personne ne peut être tenu pour responsable de la naissance d'un enfant porteur d'une anomalie génétique ou chromosomique.
\end{aretenir}

\begin{savaistu}[Le mouton Ancon : une seule mutation, une race entière]
L'agneau né en 1791 dans le troupeau de Seth Wright (Massachusetts, États-Unis) portait une mutation d'un gène du développement des os (gène \emph{Chondrodysplasia}), survenue spontanément dans la lignée germinale de l'un de ses parents. Ce caractère « pattes courtes », héréditaire de façon dominante, a été sélectionné par les éleveurs car il empêchait les moutons de franchir les clôtures : il a fondé la race \emph{Ancon}. C'est la preuve historique qu'\textbf{une seule mutation peut créer un caractère nouveau, stable et transmissible}, sans aucune intervention surnaturelle. Le même principe explique l'apparition de caractères nouveaux — y compris des anomalies — dans nos propres familles.
\end{savaistu}

\begin{exercice}[Pour vérifier tes acquis]
1. Définir une mutation et distinguer mutation génique et mutation chromosomique.

2. Un homme développe un cancer de la peau après une exposition prolongée au soleil sans protection. Ses enfants risquent-ils de naître avec la mutation responsable ? Justifier.

3. La séquence normale d'un fragment d'ADN (brin codant) est \texttt{ATG CCA GTT}. La séquence mutante est \texttt{ATG CCA GCT}. Identifier le type de mutation et prévoir sa conséquence la plus probable sur la protéine (on précise que \texttt{GTT} et \texttt{GCT} codent deux acides aminés différents).

4. Citer deux agents mutagènes et expliquer leur mode d'action général.
\end{exercice}

\begin{corrige}[Exercice — Pour vérifier tes acquis]
\textbf{1.} Une mutation est une modification stable et héréditaire du matériel génétique. La mutation génique modifie la séquence de nucléotides d'un gène (substitution, insertion, délétion) ; la mutation chromosomique modifie la structure d'un chromosome (délétion, duplication, inversion, translocation) ou le nombre de chromosomes (trisomie, monosomie).

\textbf{2.} Non. La mutation causée par les UV affecte des cellules \textbf{somatiques} (kératinocytes de la peau). Seules les mutations des cellules \textbf{germinales} (dans les testicules pour l'homme) sont transmissibles à la descendance. Les enfants de cet homme ne naîtront donc pas avec cette mutation somatique acquise (ils pourront seulement hériter d'une éventuelle prédisposition génétique constitutionnelle indépendante, comme un allèle de susceptibilité au mélanome).

\textbf{3.} Les deux séquences ont la même longueur (9 nucléotides) et ne diffèrent que par le 8\textsuperscript{e} nucléotide (\texttt{T} $\to$ \texttt{C}) : c'est une \textbf{substitution}. Le codon \texttt{GTT} devient \texttt{GCT} : comme ces codons codent deux acides aminés différents (Val $\to$ Ala), il s'agit d'une mutation \textbf{faux-sens} ; la protéine possédera un acide aminé différent à cette position, ce qui peut modifier son repliement ou son activité.

\textbf{4.} Les rayons ultraviolets (UV) du soleil et les hydrocarbures aromatiques polycycliques (goudrons) de la fumée de tabac. Ils endommagent l'ADN en provoquant des dimères de pyrimidine (UV) ou des adduits covalents sur les bases (tabac) ; lorsque les mécanismes de réparation sont défaillants ou saturés, ces lésions deviennent des mutations stables lors de la réplication.
\end{corrige}

\textbf{Transition.} Nous savons désormais \emph{comment} naissent les allèles nouveaux par mutation. Il reste à comprendre comment une mutation génique précise provoque une maladie concrète : c'est l'objet de la section suivante, consacrée à la drépanocytose, première « maladie moléculaire » élucidée par la science.

\coursec{II. Anomalies géniques : exemple de la drépanocytose}

Dans la section précédente, nous avons établi que les mutations — modifications accidentelles de la séquence d'ADN — sont à l'origine de tous les allèles nouveaux. Mais quelles sont les conséquences concrètes d'une mutation sur un individu ? Pour le comprendre, rien de tel qu'un cas réel, fréquent dans nos familles camerounaises : la \cle{drépanocytose}. Cette maladie, longtemps entourée de superstitions et de préjugés (« l'enfant est ensorcelé », « la mère a mangé quelque chose de mauvais »), s'explique aujourd'hui parfaitement par la génétique moléculaire. Suivons la démarche des chercheurs : de l'observation clinique jusqu'au nucléotide fautif.

\courssub{Présentation clinique : une maladie du sang fréquente au Cameroun}

\textbf{Observation.} Dans les hôpitaux de Yaoundé, de Douala ou de Garoua, certains enfants présentent des signes récurrents : une \cle{anémie} chronique (pâleur, fatigue, teint jaunâtre), des \cle{crises douloureuses} brutales et intenses (douleurs osseuses, abdominales), des infections fréquentes et un retard de croissance. Au microscope, leurs globules rouges (hématies) présentent une forme anormale : au lieu d'être des disques biconcaves souples, ils prennent, en condition de faible oxygénation, une forme de \cle{faucille} ou de croissant, rigide et fragile.

\begin{definition}[La drépanocytose]
La \cle{drépanocytose} (ou anémie falciforme) est une \textbf{anomalie génique autosomique récessive} due à une \textbf{substitution ponctuelle} dans le gène de la $\beta$-globine, entraînant la production d'une hémoglobine anormale, l'\textbf{hémoglobine S} (HbS). En condition de désoxygénation, l'HbS polymérise et déforme les hématies en faucilles, qui se détruisent prématurément (anémie hémolytique) et obstruent les petits vaisseaux (crises vaso-occlusives douloureuses).
\end{definition}

\textbf{Données épidémiologiques.} La drépanocytose est la maladie génétique la plus répandue au monde, avec une forte fréquence en \textbf{Afrique subsaharienne} : au Cameroun, environ \num{20} à \SI{25}{\%} de la population porte l'allèle muté, et on estime que plusieurs milliers d'enfants naissent chaque année avec la forme grave de la maladie. Cette fréquence exceptionnelle n'est pas un hasard : nous verrons qu'elle s'explique par un avantage sélectif vis-à-vis du paludisme.

\courssub{Mise en évidence expérimentale : l'électrophorèse de l'hémoglobine}

Comment passer du symptôme à la cause ? La première étape historique (1949, travaux de Linus Pauling) fut biochimique : comparer l'hémoglobine des malades à celle des sujets sains.

\begin{experience}[Électrophorèse de l'hémoglobine]
\textbf{Principe.} L'\cle{électrophorèse} sépare les protéines selon leur charge électrique : placées dans un champ électrique, sur un gel à pH alcalin (pH \approx 8,6), les protéines migrent vers l'anode ($+$) à des vitesses dépendant de leur charge nette négative.

\textbf{Protocole.} On extrait l'hémoglobine de trois individus : un sujet sain (AA), un sujet porteur sain (AS) et un malade (SS). On dépose chaque extrait au même niveau sur le gel (puits), puis on applique le champ électrique.

\textbf{Observations.}
\begin{itemize}
\item Chez le sujet AA : une seule bande, l'\textbf{HbA} (hémoglobine normale), qui migre rapidement vers l'anode.
\item Chez le malade SS : une seule bande, l'\textbf{HbS}, qui migre \emph{moins vite} que l'HbA (reste plus près du puits).
\item Chez le porteur AS : \emph{deux} bandes distinctes, l'une au niveau de l'HbA, l'autre au niveau de l'HbS, d'intensité comparable.
\end{itemize}

\textbf{Interprétation.} L'HbS migre moins vite que l'HbA : elle possède donc une \textbf{charge électrique négative moins forte}, ce qui traduit une \textbf{composition en acides aminés différente}. L'anomalie est donc \emph{moléculaire} : c'est la protéine hémoglobine elle-même qui est modifiée. Le porteur AS synthétise les deux formes : il possède donc un allèle normal \emph{et} un allèle muté, tous deux exprimés (codominance au niveau moléculaire).
\end{experience}

\begin{popfigure}
\begin{center}
\begin{tikzpicture}[scale=0.95]
% Gel background
\fill[popblueL] (-0.5,-0.5) rectangle (10.5,6.0);
\draw[thick,popdark] (-0.5,-0.5) rectangle (10.5,6.0);
% Electrodes
\node[popdark,font=\bfseries\small] at (5.0,6.4) {Cathode ($-$) : dépôt des échantillons};
\node[popdark,font=\bfseries\small] at (5.0,-1.0) {Anode ($+$)};
% Wells (puits)
\foreach \x/\t in {2.0/AA, 5.0/AS, 8.0/SS}{
  \fill[white] (\x-0.6,5.2) rectangle (\x+0.6,5.7);
  \draw[popdark,thick] (\x-0.6,5.2) rectangle (\x+0.6,5.7);
  \node[popdark,font=\bfseries\small] at (\x,5.95) {\t};
}
% Bands: HbA (lower, y~1.5), HbS (higher, y~3.0)
% AA lane: only HbA
\fill[popgreen!80!black] (1.45,1.3) rectangle (2.55,1.7);
% AS lane: HbA and HbS
\fill[popgreen!80!black] (4.45,1.3) rectangle (5.55,1.7);
\fill[poporange!80!black] (4.45,2.8) rectangle (5.55,3.2);
% SS lane: only HbS
\fill[poporange!80!black] (7.45,2.8) rectangle (8.55,3.2);
% Band labels on right
\node[popgreen!60!black,font=\small\bfseries,anchor=west] at (10.7,1.5) {HbA (rapide)};
\node[poporange!80!black,font=\small\bfseries,anchor=west] at (10.7,3.0) {HbS (lente)};
% Migration arrows
\foreach \x in {2.0, 5.0, 8.0}{
  \draw[-{Stealth[length=3mm]},thick,popdark] (\x,5.0) -- (\x,3.8);
}
% Direction arrow
\draw[-{Stealth[length=4mm]},very thick,popdark] (10.0,4.0) -- (10.0,2.0) node[midway,right,font=\small\bfseries] {Migration};
\end{tikzpicture}
\end{center}
\legende{Électrophorèse de l'hémoglobine sur gel (pH alcalin). L'HbA (bande verte), plus chargée négativement (perte d'un groupement carboxyle), migre plus loin vers l'anode ($+$) que l'HbS (bande orange). Le sujet sain $AA$ ne possède que de l'HbA ; le malade $SS$ que de l'HbS ; le porteur sain $AS$ possède les deux bandes, en quantités comparables (codominance moléculaire).}
\end{popfigure}

\courssub{Analyse comparée des séquences : un seul acide aminé modifié}

L'électrophorèse montre que la protéine est différente, mais pas \emph{où}. L'étape suivante (Vernon Ingram, 1956) consiste à comparer les séquences peptidiques. L'hémoglobine est un tétramère de deux chaînes $\alpha$ et deux chaînes $\beta$ ; la chaîne $\beta$ compte 146 acides aminés. La comparaison des chaînes $\beta$ de l'HbA et de l'HbS donne :

\begin{center}
\begin{tabularx}{\linewidth}{|c|X|X|}\hline
\rowcolor{popblueL} \textbf{Position (chaîne $\beta$)} & \textbf{HbA (normale)} & \textbf{HbS (drépanocytaire)} \\ \hline
1 & Val & Val \\ \hline
2 & His & His \\ \hline
3 & Leu & Leu \\ \hline
4 & Thr & Thr \\ \hline
5 & Pro & Pro \\ \hline
\rowcolor{poporangeL} \textbf{6} & \textbf{Glu (acide glutamique)} & \textbf{Val (valine)} \\ \hline
7 & Glu & Glu \\ \hline
8 & Lys & Lys \\ \hline
\end{tabularx}
\end{center}

\textbf{Conclusion.} Sur 146 acides aminés, \textbf{un seul} diffère : en \textbf{position 6}, l'\textbf{acide glutamique} (acide aminé hydrophile, chargé négativement à pH physiologique) est remplacé par la \textbf{valine} (acide aminé hydrophobe, neutre). Ce changement minime suffit à expliquer la différence de charge observée en électrophorèse et, surtout, la propriété dramatique de l'HbS : la valine en position 6 crée un « point d'accroche » hydrophobe à la surface de la chaîne $\beta$ qui fait s'agréger les molécules d'HbS désoxygénées en longues fibres polymères, déformant l'hématie en faucille.

\courssub{Remontée à l'ADN : la mutation ponctuelle}

Pourquoi la valine remplace-t-elle l'acide glutamique ? Il faut remonter au gène de la $\beta$-globine (\emph{HBB}), situé sur le chromosome 11 (locus 11p15,5). Comparons les séquences d'ADN du \textbf{brin codant} (sens, identique à l'ARNm hormis T/U) autour du codon n\textsuperscript{o} 6 :

\begin{center}
\begin{tabularx}{\linewidth}{|l|X|}\hline
\rowcolor{popblueL} \textbf{Allèle} & \textbf{Séquence codons 4 à 8 (brin codant 5' $\rightarrow$ 3')} \\ \hline
Normal ($A$) & ... ACT -- CCT -- \textbf{GAG} -- GAG -- AAG ... \\ \hline
Muté ($S$) & ... ACT -- CCT -- \textbf{GTG} -- GAG -- AAG ... \\ \hline
\end{tabularx}
\end{center}

D'après le code génétique (tableau fourni au Bac), le codon \textbf{GAG} code l'acide glutamique (Glu), tandis que le codon \textbf{GTG} code la valine (Val). La différence tient donc à la \textbf{substitution d'un seul nucléotide} : une \textbf{adénine (A) est remplacée par une thymine (T)} en deuxième position du codon n\textsuperscript{o} 6. C'est une \cle{mutation ponctuelle par substitution} de type \textbf{faux-sens} (le sens du codon est modifié).

\begin{aretenir}[La chaîne causale complète : de l'ADN au phénotype]
\[ \underbrace{\text{A} \xrightarrow{\text{substitution}} \text{T}}_{\text{Mutation ADN (brin codant)}} \;\Longrightarrow\; \underbrace{\text{GAG} \xrightarrow{\text{transcription}} \text{GUG}}_{\text{ARNm modifié (U remplace T)}} \;\Longrightarrow\; \underbrace{\text{Glu} \xrightarrow{\text{traduction}} \text{Val}}_{\text{Protéine modifiée (HbS)}} \;\Longrightarrow\; \underbrace{\text{Polymérisation HbS}}_{\text{Niveau moléculaire}} \;\Longrightarrow\; \underbrace{\text{Hématies en faucilles}}_{\text{Niveau cellulaire}} \;\Longrightarrow\; \underbrace{\text{Anémie + Crises vaso-occlusives}}_{\text{Phénotype (maladie)}} \]
Une seule « lettre » d'ADN modifiée sur les 3 milliards du génome suffit à provoquer une maladie grave. C'est la démonstration la plus célèbre du lien \textbf{gène $\rightarrow$ protéine $\rightarrow$ caractère}.
\end{aretenir}

\begin{methode}[Identifier une anomalie génique à partir de séquences d'ADN (Méthode Bac)]
Face à un exercice de comparaison de séquences (très fréquent au Baccalauréat), procède en étapes rigoureuses :
\begin{enumerate}
\item \textbf{Aligner} la séquence du patient avec la séquence de référence (sujet sain), codon par codon (groupes de 3 nucléotides).
\item \textbf{Localiser} la ou les différences : \textbf{substitution} (remplacement d'un nucléotide), \textbf{délétion} (perte d'un ou plusieurs nucléotides) ou \textbf{insertion} (addition).
\item \textbf{Traduire} le codon de référence et le codon muté à l'aide du tableau du code génétique fourni dans l'énoncé (attention : le tableau utilise souvent l'alphabet ADN avec T, ou ARN avec U).
\item \textbf{Déduire} l'effet sur la protéine :
      \begin{itemize}
      \item Acide aminé changé $\rightarrow$ \textbf{mutation faux-sens} (ex: Glu $\rightarrow$ Val).
      \item Codon STOP prématuré (UAA, UAG, UGA) $\rightarrow$ \textbf{mutation non-sens} (protéine tronquée, souvent non fonctionnelle).
      \item Acide aminé inchangé $\rightarrow$ \textbf{mutation silencieuse} (redondance du code génétique).
      \item Décalage du cadre de lecture (délétion/insertion non multiple de 3) $\rightarrow$ \textbf{mutation frameshift} (protéine totalement modifiée en aval).
      \end{itemize}
\item \textbf{Conclure} en reliant la modification protéique au phénotype observé (perte de fonction, gain de fonction toxique, etc.).
\end{enumerate}
\end{methode}

\courssub{Transmission : une anomalie autosomique récessive}

Notons $A$ l'allèle normal et $S$ l'allèle muté du gène \emph{HBB} (chromosome 11). Trois génotypes sont possibles :

\begin{center}
\begin{tabularx}{\linewidth}{|c|X|X|}\hline
\rowcolor{popblueL} \textbf{Génotype} & \textbf{Phénotype} & \textbf{Hémoglobine présente} \\ \hline
$A//A$ (ou $AA$) & Sain (non porteur) & 100 \% HbA \\ \hline
$A//S$ (ou $AS$) & \textbf{Porteur sain} (trait drépanocytaire), asymptomatique en conditions normales & HbA $\approx$ 60 \%, HbS $\approx$ 40 \% \\ \hline
$S//S$ (ou $SS$) & \textbf{Malade} drépanocytaire (anémie falciforme majeure) & quasi 100 \% HbS \\ \hline
\end{tabularx}
\end{center}

La maladie est dite \textbf{récessive} car l'allèle $S$ ne provoque la maladie qu'à l'état homozygote ($S//S$) ; l'allèle $A$ est dominant sur le plan phénotypique (santé). Elle est dite \textbf{autosomique} car le gène est porté par un autosome (chromosome 11), donc la maladie touche indifféremment garçons et filles (pas de liaison au sexe).

\begin{exemplebox}[Croisement de deux porteurs sains : $AS \times AS$]
Chaque parent produit deux types de gamètes : $1/2$ porteurs de $A$ et $1/2$ porteurs de $S$. L'échiquier de croisement (tableau de Punnett) donne :

\begin{center}
\begin{tabular}{|c|c|c|}\hline
\rowcolor{popblueL} \textbf{Gamètes $\downarrow$ / $\rightarrow$} & $A$ ($1/2$) & $S$ ($1/2$) \\ \hline
$A$ ($1/2$) & $AA$ ($1/4$) : Sain & $AS$ ($1/4$) : Porteur sain \\ \hline
$S$ ($1/2$) & $AS$ ($1/4$) : Porteur sain & $SS$ ($1/4$) : Malade \\ \hline
\end{tabular}
\end{center}

\textbf{Bilan :} $1/4$ d'enfants sains $AA$, $1/2$ de porteurs sains $AS$, $1/4$ de malades $SS$. \textbf{Deux parents sains porteurs ont donc, à chaque grossesse, 25 \% de risque d'avoir un enfant drépanocytaire} — et 50 \% de risque d'avoir un enfant porteur sain. Ces probabilités sont indépendantes d'une grossesse à l'autre.
\end{exemplebox}

\begin{popfigure}
\begin{center}
\begin{tikzpicture}[
  male/.style={draw,thick,popblue,fill=white,minimum size=0.6cm,inner sep=2pt},
  female/.style={draw,thick,poppink,fill=white,circle,minimum size=0.6cm,inner sep=2pt},
  maleS/.style={draw,thick,popblue,fill=popblue!40,minimum size=0.6cm,inner sep=2pt},
  femaleS/.style={draw,thick,poppink,fill=poppink!40,circle,minimum size=0.6cm,inner sep=2pt},
  maleM/.style={draw,thick,popblue,fill=popdark,minimum size=0.6cm,inner sep=2pt},
  femaleM/.style={draw,thick,poppink,fill=popdark,circle,minimum size=0.6cm,inner sep=2pt},
  lbl/.style={font=\small,popdark,below=2pt}
]
% Génération I
\node[maleS,label={[lbl]I-1 $AS$}] (pere) at (0,0) {};
\node[femaleS,label={[lbl]I-2 $AS$}] (mere) at (3,0) {};
\draw[thick,popdark] (pere) -- (mere) node[midway,above=2pt,font=\small] {Union};
\coordinate (union) at (1.5,0);
\draw[thick,popdark] (union) -- (1.5,-0.8);
\draw[thick,popdark] (-1.5,-0.8) -- (4.5,-0.8);
% Génération II
\node[male,label={[lbl]II-1 $AA$}] (e1) at (-1.5,-1.6) {};
\node[femaleS,label={[lbl]II-2 $AS$}] (e2) at (0.2,-1.6) {};
\node[maleS,label={[lbl]II-3 $AS$}] (e3) at (1.9,-1.6) {};
\node[femaleM,label={[lbl]II-4 $SS$}] (e4) at (3.6,-1.6) {};
\foreach \e in {e1,e2,e3,e4}{ \draw[thick,popdark] (\e.north) -- (\e.north |- 0,-0.8); }
% Légende
\begin{scope}[xshift=5.5cm, yshift=0.5cm]
\node[male,label={[lbl,right]Homme sain $AA$}] at (0,0.5) {};
\node[femaleS,label={[lbl,right]Femme porteuse saine $AS$}] at (0,-0.5) {};
\node[femaleM,label={[lbl,right]Femme malade $SS$}] at (0,-1.5) {};
\node[draw,thick,popblue,fill=popblue!40,minimum size=0.6cm,label={[lbl,right]Homme porteur sain $AS$}] at (0,-2.5) {};
\node[draw,thick,popblue,fill=popdark,minimum size=0.6cm,label={[lbl,right]Homme malade $SS$}] at (0,-3.5) {};
\end{scope}
\end{tikzpicture}
\end{center}
\legende{Arbre généalogique d'une famille drépanocytaire. Carré = homme, cercle = femme ; symbole blanc = sain $AA$ ; symbole à moitié rempli (hachuré) = porteur sain $AS$ ; symbole entièrement rempli (noir) = malade $SS$. Les parents I-1 et I-2, tous deux porteurs sains, ont engendré 4 enfants illustrant les 3 génotypes possibles, dont une fille malade (II-4).}
\end{popfigure}

\begin{attention}[Pièges à éviter et idées reçues -- \textbf{À connaître pour le Bac et la vie}]
\begin{itemize}
\item \textbf{Le porteur sain ($AS$) n'est PAS un malade.} Il vit normalement, fait du sport, a une espérance de vie normale. Il ne faut \textbf{ni le stigmatiser, ni lui interdire le mariage}. La seule précaution est l'\textbf{information génétique} : un couple $AS \times AS$ doit connaître le risque de 25 \% à chaque grossesse pour bénéficier d'un suivi (conseil génétique, diagnostic prénatal possible).
\item \textbf{La drépanocytose n'est PAS contagieuse.} Elle ne se transmet ni par le sang (sauf transfusion), ni par la salive, ni par contact. Elle se transmet \textbf{uniquement par les gamètes}, des parents aux enfants.
\item \textbf{Les parents ne sont pas « coupables ».} Avoir un enfant $SS$ est la réalisation d'un risque probabiliste (1/4), non une faute ou une malédiction. La culpabilisation maternelle est un préjugé nuisible à combattre.
\item \textbf{Attention au vocabulaire :} « Trait drépanocytaire » = phénotype du porteur sain $AS$ (asymptomatique). « Drépanocytose » ou « Anémie falciforme » = phénotype du malade $SS$. Ne pas confondre.
\item \textbf{Électrophorèse :} L'HbS migre \emph{moins vite} (plus près du puits/départ) que l'HbA car elle est \emph{moins négative} (perte d'une charge $-$) $\rightarrow$ attraction plus faible vers l'anode ($+$).
\end{itemize}
\end{attention}

\begin{exoresolu}[Analyse d'un document de séquences -- Type Bac]
\textbf{Énoncé.} On compare un fragment du gène de la $\beta$-globine chez un sujet sain et chez un enfant anémique admis au service de pédiatrie de l'hôpital de district de Bafoussam :

Sujet sain (référence) : ... CCT \textbf{GAG} GAG ... \\
Enfant malade : \hphantom{Sujet sain (référence) : } ... CCT \textbf{GTG} GAG ...

On donne la correspondance (extrait du code génétique) : GAG = Acide glutamique (Glu) ; GTG = Valine (Val).

\textbf{Questions :}
\begin{enumerate}
\item Identifier le type de mutation survenu au niveau de l'ADN. Préciser sa nature.
\item Expliquer l'origine moléculaire et cellulaire de la maladie de l'enfant.
\item Les deux parents de l'enfant sont phénotypiquement sains. Est-ce compatible avec la génétique de la maladie ? Justifier en précisant leur génotype probable.
\end{enumerate}

\tcblower
\textbf{Solution détaillée.}

(1) \emph{Identification de la mutation.}
L'alignement des séquences codon par codon montre une unique différence au niveau du deuxième codon présenté (codon 6 de la $\beta$-globine) :
\begin{itemize}
\item Séquence référence : \textbf{GAG}
\item Séquence mutée : \textbf{GTG}
\end{itemize}
Le deuxième nucléotide du codon est passé de \textbf{A} (adénine) à \textbf{T} (thymine). Il s'agit d'une \textbf{mutation ponctuelle par substitution} (un seul nucléotide remplacé). Comme le codon muté code un acide aminé différent (Val au lieu de Glu), c'est une mutation de type \textbf{faux-sens}.

(2) \emph{Origine de la maladie (chaîne causale).}
\begin{itemize}
\item \textbf{Niveau ADN :} Substitution A $\rightarrow$ T sur le brin codant du gène \emph{HBB} (chromosome 11).
\item \textbf{Niveau ARNm :} Le codon GAG devient GUG (U remplaçant T).
\item \textbf{Niveau Protéine :} L'acide glutamique (Glu, hydrophile, chargé $-$) en position 6 de la chaîne $\beta$ est remplacé par la valine (Val, hydrophobe, neutre).
\item \textbf{Niveau Moléculaire :} Cette substitution crée un patch hydrophobe à la surface de la chaîne $\beta$. En condition de désoxygénation, l'HbS polymérise en fibres rigides (polymères de 14 brins).
\item \textbf{Niveau Cellulaire :} Les fibres déforment l'hématie en faucille (drépanocyte), la rigidifient et la fragilisent.
\item \textbf{Niveau Organisme :} Destruction prématurée des hématies (anémie hémolytique chronique, ictère) + obstruction des capillaires par les drépanocytes rigides (crises vaso-occlusives douloureuses, nécroses, infections).
\end{itemize}
L'enfant est atteint de \textbf{drépanocytose majeure} ; son génotype est $S//S$ (homozygote muté).

(3) \emph{Génotype des parents.}
Oui, c'est \textbf{parfaitement compatible} et même \textbf{obligatoire} pour une maladie autosomique récessive.
\begin{itemize}
\item L'enfant $SS$ a nécessairement reçu un allèle $S$ de son père et un allèle $S$ de sa mère.
\item Comme les parents sont sains (phénotype normal), ils ne peuvent pas être $SS$. Ils sont donc nécessairement hétérozygotes \textbf{$AS$} (porteurs sains).
\item Le croisement $AS \times AS$ donne, à chaque grossesse, 1/4 de risque d'enfant $SS$. La naissance de cet enfant malade correspond à la réalisation de ce risque statistique (25 \%), et non à une anomalie de transmission.
\end{itemize}
\end{exoresolu}

\courssub{Pourquoi l'allèle S est-il si fréquent ? L'avantage sélectif face au paludisme}

Un paradoxe saute aux yeux : si l'allèle $S$ provoque une maladie grave à l'état homozygote, pourquoi la sélection naturelle ne l'a-t-elle pas éliminé au fil des générations ? Pourquoi touche-t-il jusqu'à 1 personne sur 4 ou 5 au Cameroun ?

\textbf{Observation épidémiologique.} La carte mondiale de fréquence de l'allèle $S$ se superpose remarquablement à la carte des zones d'\textbf{endémie palustre} à \emph{Plasmodium falciparum} (Afrique subsaharienne, bassin méditerranéen, Moyen-Orient, Inde).

\textbf{Interprétation scientifique (Haldane, Allison, 1954).} Le parasite du paludisme se multiplie à l'intérieur des hématies. Or, les hématies des porteurs $AS$ (contenant $\approx$ 40 \% d'HbS) offrent un environnement défavorable au développement du parasite : elles se déforment plus facilement en condition de faible oxygène (milieu interne du parasite), sont éliminées plus vite par la rate, et présentent des altérations métaboliques inhibant la croissance de \emph{Plasmodium}.
\begin{itemize}
\item Sujets $AA$ : Sensibles au paludisme grave (forte mortalité infantile historique).
\item Sujets $SS$ : Victimes de la drépanocytose (forte mortalité sans soins modernes).
\item Sujets $AS$ (hétérozygotes) : \textbf{Résistants aux formes graves du paludisme} (meilleure survie).
\end{itemize}
C'est un cas classique d'\textbf{avantage sélectif de l'hétérozygote} (ou surdominance) : la sélection naturelle \emph{maintient} l'allèle $S$ dans la population au prix de la naissance d'enfants $SS$, car le gain de survie des $AS$ en zone palustre compense la perte des $SS$. Cet équilibre évolutif explique la fréquence élevée du trait drépanocytaire là où le paludisme sévit depuis des millénaires.

\begin{savaistu}[La drépanocytose au Cameroun : dépistage et prise en charge]
Au Cameroun, environ \textbf{20 à 25 \% de la population} est porteuse du trait drépanocytaire ($AS$), soit plusieurs millions de personnes. La maladie concerne des milliers de nouveau-nés par an.
\begin{itemize}
\item \textbf{Dépistage prénuptial recommandé :} Un simple test sanguin (test de désoxygénation au métabisulfite de sodium -- « test d'Emmel » -- ou électrophorèse de l'hémoglobine, réalisés dans la plupart des laboratoires d'analyses et lors de campagnes gratuites) permet à chaque futur couple de connaître son statut ($AA$, $AS$, $SS$, ou autres variants comme $AC$, $SC$).
\item \textbf{Conseil génétique :} Un couple $AS \times AS$ informé peut choisir un suivi de grossesse avec diagnostic prénatal (biopsie de trophoblaste ou amniocentèse) pour connaître le génotype du fœtus.
\item \textbf{Prise en charge moderne :} Hydratation rigoureuse, prophylaxie anti-paludique, vaccination (pneumocoque, \emph{Haemophilus}, méningocoque), acide folique, traitement antalgique des crises. L'\textbf{hydroxyurée} ( hydroxycarbamide) est le traitement de référence : elle stimule la production d'hémoglobine fœtale (HbF, $\alpha_2\gamma_2$) qui inhibe la polymérisation de l'HbS, réduisant drastiquement la fréquence des crises et la mortalité. Des consultations spécialisées existent (ex: CHU Yaoundé, Hôpital Laquintinie Douala).
\end{itemize}
\end{savaistu}

\begin{aretenir}[L'essentiel de la section II -- \textbf{Fiche de révision Bac}]
\begin{itemize}
\item \textbf{Définition :} La drépanocytose est une anomalie génique \textbf{autosomique récessive} (chromosome 11, gène \emph{HBB}).
\item \textbf{Mutation causale :} Substitution ponctuelle A $\rightarrow$ T (brin codant) au codon 6 $\rightarrow$ Glu $\rightarrow$ Val sur la chaîne $\beta$-globine $\rightarrow$ HbS.
\item \textbf{Conséquences moléculaires/cellulaires :} HbS polymérise en condition désoxygénée $\rightarrow$ hématies en faucilles (drépanocytes) rigides et fragiles $\rightarrow$ anémie hémolytique + crises vaso-occlusives.
\item \textbf{Diagnostic biologique :} Électrophorèse de l'hémoglobine : $AA$ (1 bande HbA), $AS$ (2 bandes HbA+HbS), $SS$ (1 bande HbS). Codominance moléculaire.
\item \textbf{Transmission :} Croisement $AS \times AS$ $\rightarrow$ 1/4 $AA$, 1/2 $AS$, 1/4 $SS$ à chaque grossesse. Risque de 25 \% d'enfant malade pour un couple de porteurs.
\item \textbf{Fréquence élevée :} Avantage sélectif de l'hétérozygote $AS$ (résistance au paludisme grave) en zone d'endémie. Équilibre mutation-sélection.
\item \textbf{Éradication des préjugés :} Cause strictement génétique (hasard méiotique). Ni contagion, ni sorcellerie, ni faute parentale. Porteurs sains = individus en bonne santé. Information et dépistage = prévention efficace.
\end{itemize}
\end{aretenir}

\coursec{III. Autres anomalies géniques : hémophilie et aperçus}

\noindent Après avoir analysé la drépanocytose, exemple paradigmatique d'anomalie \cle{autosomique récessive} touchant indifféremment les deux sexes, nous allons étudier une pathologie dont la transmission obéit à une logique différente : l'hémophilie. Cette maladie a marqué l'histoire des familles royales européennes et a servi de modèle historique pour établir les lois de la génétique liée au sexe. Nous en profiterons pour dresser un panorama d'autres anomalies géniques au programme, afin de renforcer une conviction scientifique centrale : ces maladies sont la conséquence de \cle{mutations géniques} héréditaires, nullement le fruit d'une faute morale, d'un « mauvais sort » ou d'une malédiction familiale.

\courssub{L'hémophilie : un trouble de la coagulation lié au chromosome X}

\noindent L'hémophilie est un ensemble de maladies hémorragiques héréditaires caractérisées par un défaut de la \cle{coagulation sanguine}. Rappelons que la coagulation est une cascade enzymatique complexe aboutissant à la formation d'un caillot de fibrine. Chez l'hémophile, l'un des facteurs de cette cascade est absent ou dysfonctionnel.

\begin{definition}[Hémophilie et facteurs de coagulation]
L'hémophilie est une maladie génétique due à une mutation affectant un gène codant pour un facteur de coagulation. Les deux formes principales sont :
\begin{itemize}
    \item L'\textbf{hémophilie A} (85 \% des cas) : déficit en \cle{facteur VIII} (anti-hémophile A).
    \item L'\textbf{hémophilie B} (15 \% des cas) : déficit en \cle{facteur IX} (anti-hémophile B ou facteur de Christmas).
\end{itemize}
La gravité (mineure, modérée, sévère) corrèle avec le taux résiduel d'activité du facteur circulant (ex. : $< 1\ \%$ pour la forme sévère).
\end{definition}

\noindent Le gène du facteur VIII (\emph{F8}) et celui du facteur IX (\emph{F9}) sont tous deux localisés sur le \cle{chromosome X} (bande Xq28 pour \emph{F8}, Xq27 pour \emph{F9}). Cette localisation a des conséquences majeures sur le mode de transmission, que nous allons établir par une démarche d'investigation.

\begin{experience}[Analyse de la transmission de l'hémophilie dans une fratrie]
\textbf{Observation (document clinique) :} On observe une famille où le père est sain, la mère est phénotypiquement saine mais a un père hémophile. Ils ont quatre enfants : deux garçons (l'un hémophile, l'autre sain) et deux filles (toutes deux phénotypiquement saines).
\textbf{Question :} Comment expliquer que seuls les garçons soient atteints et que la mère, bien que saine, transmette la maladie ?
\textbf{Hypothèses :}
\begin{enumerate}
    \item Le gène responsable est porté par un autosome (transmission autosomique).
    \item Le gène responsable est porté par le chromosome X (transmission gonosomique liée à l'X).
\end{enumerate}
\textbf{Interprétation :}
\begin{itemize}
    \item Si transmission autosomique récessive : la mère (saine) serait hétérozygote $A/a$. Le père (sain) serait $A/A$. Croisement $A/a \times A/A \rightarrow$ 1/2 $A/A$ (sains), 1/2 $A/a$ (porteurs sains). \textbf{Aucun enfant malade attendu.} Contredit l'observation (un garçon malade).
    \item Si transmission récessive liée à l'X : Notons $X^H$ l'allèle normal, $X^h$ l'allèle muté. Le père sain est $X^H Y$. La mère, fille d'un père hémophile ($X^h Y$), a obligatoirement reçu le $X^h$ de son père. Elle est donc \cle{conductrice} (porteuse saine) de génotype $X^H X^h$. Croisement $X^H X^h \times X^H Y$.
\end{itemize}
\textbf{Conclusion :} Seule l'hypothèse « gène sur le chromosome X » rend compte de l'apparition de garçons malades nés de parents phénotypiquement sains (pour la mère) et de l'absence de transmission père-fils. L'hémophilie est une maladie \cle{récessive liée au chromosome X}.
\end{experience}

\noindent Cette analyse expérimentale (sur document) nous permet de poser la propriété fondamentale de la transmission liée à l'X.

\begin{propriete}[Transmission récessive liée au chromosome X]
\begin{itemize}
    \item L'allèle muté ($X^h$) est porté par le chromosome X. L'allèle normal ($X^H$) est dominant.
    \item \textbf{Hommes (XY) :} N'ayant qu'un seul chromosome X, ils expriment le phénotype muté dès qu'ils portent l'allèle $X^h$ (génotype $X^h Y$). Ils sont \textbf{hémizygotes}.
    \item \textbf{Femmes (XX) :} Elles ne sont malades que si elles sont homozygotes $X^h X^h$ (très rare, nécessite un père malade et une mère conductrice). L'hétérozygote $X^H X^h$ est une \textbf{femme conductrice} (porteuse saine), généralement asymptomatique grâce à l'inactivation aléatoire d'un X (lyonisation), bien que certaines puissent présenter des symptômes légers (taux de facteur $\approx 50\ \%$).
    \item \textbf{Transmission père-fils : IMPOSSIBLE.} Le père transmet son chromosome Y à ses fils.
    \item \textbf{Transmission père-filles : SYSTÉMATIQUE.} Le père transmet son unique chromosome X (muté s'il est malade) à toutes ses filles, qui deviennent conductrices.
\end{itemize}
\end{propriete}

\begin{popfigure}
\begin{tikzpicture}[scale=0.9, every node/.style={font=\small}]
    % Styles
    \tikzset{
        male/.style={rectangle, draw=popdark, thick, minimum width=1.2cm, minimum height=0.8cm, align=center},
        female/.style={circle, draw=popdark, thick, minimum width=1.2cm, minimum height=0.8cm, align=center},
        affected/.style={fill=popred!80},
        carrier/.style={fill=poporange!50, pattern=north east lines, pattern color=poporange},
        normal/.style={fill=popgreen!20},
        gamete/.style={rectangle, draw=popblue, fill=popblue!10, rounded corners, minimum width=1cm, minimum height=0.6cm, align=center},
        arrow/.style={->, thick, popdark},
        label/.style={font=\tiny, text=popdark}
    }
    % Parents
    \node[female, carrier, align=center] (mother) at (0, 3){Mère\\$X^H X^h$\\Conductrice};
    \node[male, normal, align=center] (father) at (5, 3){Père\\$X^H Y$\\Sain};
    % Gametes Mother
    \node[gamete, align=center] (gM1) at (-1.5, 1.5){Ovule\\$X^H$};
    \node[gamete, align=center] (gM2) at (1.5, 1.5){Ovule\\$X^h$};
    % Gametes Father
    \node[gamete, align=center] (gF1) at (3.5, 1.5){Spermatozoïde\\$X^H$};
    \node[gamete, align=center] (gF2) at (6.5, 1.5){Spermatozoïde\\$Y$};
    % Arrows Gametes
    \draw[arrow] (mother.south west) -- (gM1.north);
    \draw[arrow] (mother.south east) -- (gM2.north);
    \draw[arrow] (father.south west) -- (gF1.north);
    \draw[arrow] (father.south east) -- (gF2.north);
    % Offspring (Punnett Square results)
    % Row 1: XH (Mother) x XH (Father) -> XH XH (Girl normal)
    \node[female, normal, align=center] (c11) at (-1.5, 0){Fille\\$X^H X^H$\\Saine non porteuse};
    % Row 1: XH (Mother) x Y (Father) -> XH Y (Boy normal)
    \node[male, normal, align=center] (c12) at (1.5, 0){Garçon\\$X^H Y$\\Sain};
    % Row 2: Xh (Mother) x XH (Father) -> XH Xh (Girl carrier)
    \node[female, carrier, align=center] (c21) at (3.5, 0){Fille\\$X^H X^h$\\Conductrice};
    % Row 2: Xh (Mother) x Y (Father) -> Xh Y (Boy affected)
    \node[male, affected, align=center] (c22) at (6.5, 0){Garçon\\$X^h Y$\\Hémophile};
    % Arrows Fertilization
    \draw[arrow] (gM1.south) -- (c11.north);
    \draw[arrow] (gF1.south west) -- (c11.north east);
    \draw[arrow] (gM1.south east) -- (c12.north west);
    \draw[arrow] (gF2.south west) -- (c12.north);
    \draw[arrow] (gM2.south west) -- (c21.north);
    \draw[arrow] (gF1.south east) -- (c21.north east);
    \draw[arrow] (gM2.south east) -- (c22.north west);
    \draw[arrow] (gF2.south east) -- (c22.north);
    % Probabilities
    \node[label, align=center] at (-1.5, -1.2) {1/4 (25 \%)};
    \node[label, align=center] at (1.5, -1.2) {1/4 (25 \%)};
    \node[label, align=center] at (3.5, -1.2) {1/4 (25 \%)};
    \node[label, align=center] at (6.5, -1.2) {1/4 (25 \%)};
    % Legend
    \node[label, align=left, anchor=north west] at (-2.5, -1.8) {
        \textbf{Légende :} \textcolor{popgreen!60!black}{\rule{3mm}{3mm}} Sain non porteur \quad
        \textcolor{poporange!70!black}{\rule{3mm}{3mm}} Conductrice (fille) \quad
        \textcolor{popred!80}{\rule{3mm}{3mm}} Atteint (garçon)
    };
\end{tikzpicture}
\legende{Échiquier de croisement (carré de Punnett) pour le couple : Femme conductrice ($X^H X^h$) $\times$ Homme sain ($X^H Y$). À chaque grossesse, la probabilité d'avoir un garçon hémophile est de 25 \%, une fille conductrice 25 \%, un enfant sain non porteur 50 \%.}
\end{popfigure}

\begin{exoresolu}[Calcul des risques pour un couple : femme conductrice $\times$ homme sain]
\textbf{Énoncé :} Une femme conductrice d'hémophilie A ($X^H X^h$) et un homme sain ($X^H Y$) souhaitent avoir un enfant. Déterminer les probabilités génotypiques et phénotypiques à la naissance.
\textbf{Solution détaillée :}
\begin{enumerate}
    \item \textbf{Gamètes maternels :} Méiose $\rightarrow$ 1/2 ovules $X^H$, 1/2 ovules $X^h$.
    \item \textbf{Gamètes paternels :} Méiose $\rightarrow$ 1/2 spermatozoïdes $X^H$, 1/2 spermatozoïdes $Y$.
    \item \textbf{Fécondation aléatoire (croisement) :} On combine les gamètes (voir figure ci-dessus).
    \begin{align*}
        X^H (\text{mère}) \times X^H (\text{père}) &\rightarrow X^H X^H \quad \text{(Fille saine non porteuse)} \quad p = 1/4 \\
        X^H (\text{mère}) \times Y (\text{père}) &\rightarrow X^H Y \quad \text{(Garçon sain)} \quad p = 1/4 \\
        X^h (\text{mère}) \times X^H (\text{père}) &\rightarrow X^H X^h \quad \text{(Fille conductrice)} \quad p = 1/4 \\
        X^h (\text{mère}) \times Y (\text{père}) &\rightarrow X^h Y \quad \text{(Garçon hémophile)} \quad p = 1/4
    \end{align*}
    \item \textbf{Bilan phénotypique par sexe :}
    \begin{itemize}
        \item \textbf{Filles (50 \% des naissances) :} 100 \% phénotypiquement saines. Parmi elles : 1/2 non porteuses ($X^H X^H$), 1/2 conductrices ($X^H X^h$).
        \item \textbf{Garçons (50 \% des naissances) :} 1/2 sains ($X^H Y$), 1/2 hémophiles ($X^h Y$).
    \end{itemize}
    \item \textbf{Risque global par grossesse :} 25 \% d'avoir un garçon hémophile ; 25 \% d'avoir une fille conductrice ; 50 \% d'avoir un enfant (fille ou garçon) sain et non porteur de l'allèle muté.
\end{enumerate}
\end{exoresolu}

\begin{attention}[Pièges fréquents et idées reçues]
\begin{itemize}
    \item \textbf{« La mère est responsable de la maladie de son fils. »} \textbf{FAUX.} La transmission du chromosome X maternel ($X^H$ ou $X^h$) relève du \cle{hasard méiotique} (ségrégation aléatoire des chromosomes homologues en Anaphase I). La mère ne « choisit » pas quel X transmettre. Aucune faute morale n'est en cause.
    \item \textbf{« Un père hémophile transmet la maladie à ses fils. »} \textbf{FAUX.} Le père donne son chromosome Y à ses fils. Il transmet son $X^h$ à \textbf{toutes ses filles} (qui deviennent conductrices), mais \textbf{jamais à ses fils}.
    \item \textbf{« Une femme conductrice est toujours parfaitement saine. »} \textbf{PRUDENCE.} Du fait de l'inactivation aléatoire de l'un des deux X (lyonisation), certaines conductrices peuvent avoir un taux de facteur VIII/IX bas (30--50 \%) et présenter des saignements anormaux (règles abondantes, hématomes faciles). Elles ne sont pas « malades » au sens sévère, mais ne sont pas strictement asymptomatiques.
    \item \textbf{Notation :} N'oubliez jamais le chromosome sexuel dans le génotype : $X^H Y$ et non $H/$-- ; $X^h Y$ et non $h/$--.
\end{itemize}
\end{attention}

\courssub{Étude historique : la descendance de la reine Victoria — une preuve par l'arbre généalogique}

\noindent L'histoire de l'hémophilie dans les familles royales européennes au XIX$^{\text{e}}$ et au début du XX$^{\text{e}}$ siècle constitue une « expérience naturelle » à grande échelle qui a permis de valider le modèle chromosomique de l'hérédité (théorie de Morgan, 1910). La reine Victoria (1819--1901), probablement porteuse d'une \cle{mutation de novo} (apparue spontanément dans sa lignée germinale ou chez son père âgé), a transmis l'allèle $X^h$ à plusieurs de ses neuf enfants.

\begin{popfigure}
\begin{tikzpicture}[scale=0.75, every node/.style={font=\footnotesize, transform shape}]
    % Styles
    \tikzset{
        pers/.style={draw=popdark, thick, minimum width=1.4cm, minimum height=0.7cm, align=center, inner sep=1pt},
        male/.style={pers, rectangle},
        female/.style={pers, circle},
        aff/.style={fill=popred!80, text=white},
        car/.style={fill=poporange!60, pattern=north east lines, pattern color=poporange!80!black},
        nor/.style={fill=popgreen!30},
        unk/.style={fill=popgray!20},
        link/.style={thick, popdark},
        genlabel/.style={font=\bfseries\large, text=popblue, anchor=west}
    }
    % Generation 1: Victoria & Albert
    \node[female, car, align=center] (Vic) at (0, 7){Victoria\\$X^H X^h$};
    \node[male, nor, align=center] (Alb) at (4, 7){Albert\\$X^H Y$};
    \draw[link] (Vic) -- (Alb) node[midway, above, font=\tiny] {Mariage 1840};
    % Generation 2: Children (simplified: 3 relevant)
    % Alice (carrier) - left
    \node[female, car, align=center] (Alice) at (-4, 5){Alice\\$X^H X^h$};
    % Leopold (affected) - center
    \node[male, aff, align=center] (Leo) at (0, 5){Léopold\\$X^h Y$};
    % Beatrice (carrier) - right
    \node[female, car, align=center] (Beat) at (4, 5){Béatrice\\$X^H X^h$};
    % Links Gen1 -> Gen2
    \foreach \c in {Alice, Leo, Beat} \draw[link] (Vic) |- (\c);
    \foreach \c in {Alice, Leo, Beat} \draw[link] (Alb) |- (\c);
    % Generation 3: Descendants key figures
    % Alice line -> Alexandra (Tsarina) -> Alexis (Tsarevich)
    \node[female, car, align=center] (Alexandra) at (-6, 3){Alexandra\\$X^H X^h$};
    \node[male, aff, align=center] (Alexis) at (-6, 1.5){Alexis\\$X^h Y$};
    \draw[link] (Alice) -- (Alexandra);
    \draw[link] (Alexandra) -- (Alexis);
    \node[genlabel] at (-8, 3) {Russie};
    % Beatrice line -> Victoria Eugenie -> Alfonso, Juan Carlos
    \node[female, car, align=center] (VicEug) at (6, 3){Vic. Eugénie\\$X^H X^h$};
    \node[male, aff, align=center] (Alfonso) at (5, 1.5){Alphonse\\$X^h Y$};
    \node[male, aff, align=center] (Juan) at (7, 1.5){Juan Carlos\\$X^h Y$};
    \draw[link] (Beat) -- (VicEug);
    \draw[link] (VicEug) -- (Alfonso);
    \draw[link] (VicEug) -- (Juan);
    \node[genlabel] at (8.5, 3) {Espagne};
    % Leopold line -> Daughter (carrier) -> Rupert (affected) - simplified
    \node[female, car, align=center] (LeoDau) at (-1, 3){Fille de Léopold\\$X^H X^h$};
    \node[male, aff, align=center] (Rupert) at (-1, 1.5){Rupert\\$X^h Y$};
    \draw[link] (Leo) -- (LeoDau);
    \draw[link] (LeoDau) -- (Rupert);
    \node[genlabel] at (1.5, 3) {Allemagne};
    % Legend
    \begin{scope}[shift={(-9, -0.5)}]
        \node[male, nor, minimum width=0.6cm, minimum height=0.4cm] (l1) at (0,0) {};
        \node[right=2pt of l1, font=\tiny] {Homme sain};
        \node[male, aff, minimum width=0.6cm, minimum height=0.4cm] (l2) at (0,-0.6) {};
        \node[right=2pt of l2, font=\tiny] {Homme hémophile};
        \node[female, car, minimum width=0.6cm, minimum height=0.4cm] (l3) at (0,-1.2) {};
        \node[right=2pt of l3, font=\tiny] {Femme conductrice};
        \node[female, nor, minimum width=0.6cm, minimum height=0.4cm] (l4) at (0,-1.8) {};
        \node[right=2pt of l4, font=\tiny] {Femme saine non porteuse};
    \end{scope}
\end{tikzpicture}
\legende{Arbre généalogique simplifié de la descendance de la reine Victoria montrant la transmission de l'hémophilie (allèle $X^h$). En rouge : hommes atteints ($X^h Y$). En orange hachuré : femmes conductrices ($X^H X^h$). En vert : individus sains non porteurs. On note l'absence de transmission père-fils (Léopold $\rightarrow$ fille conductrice $\rightarrow$ petit-fils atteint) et la propagation via les filles conductrices vers les cours de Russie, d'Espagne et d'Allemagne.}
\end{popfigure}

\begin{methode}[Analyser un arbre généalogique pour identifier une transmission liée à l'X récessive]
Pour conclure au mode de transmission « récessif lié à l'X » à partir d'un arbre généalogique, recherchez la \textbf{signature} suivante (tous les critères doivent être vérifiés) :
\begin{enumerate}
    \item \textbf{Prédominance masculine :} La quasi-totalité des individus atteints sont des \textbf{hommes}.
    \item \textbf{Absence de transmission père $\rightarrow$ fils :} Aucun homme atteint n'a de fils atteint (le père donne Y à ses fils).
    \item \textbf{Transmission père $\rightarrow$ filles :} Tout homme atteint a \textbf{toutes ses filles conductrices} (il leur donne son unique X muté).
    \item \textbf{Saut de génération via les femmes :} La maladie « saute » une génération : homme atteint $\rightarrow$ fille conductrice (saine) $\rightarrow$ petit-fils atteint (50 \% de chances).
    \item \textbf{Femmes atteintes rares :} Une femme atteinte n'apparaît que si son père est atteint \textbf{ET} sa mère conductrice (croisement $X^h Y \times X^H X^h \rightarrow 1/2\ X^h X^h$).
\end{enumerate}
Si ces 5 critères sont remplis $\rightarrow$ Transmission récessive liée à l'X confirmée.
\end{methode}

\begin{savaistu}[L'hémophilie, « maladie des rois » et géopolitique]
La mutation de la reine Victoria s'est propagée dans les cours d'Europe par les mariages dynastiques de ses petites-filles :
\begin{itemize}
    \item \textbf{Russie :} Sa petite-fille Alexandra (fille d'Alice) épouse le Tsar Nicolas II. Leur fils unique, le Tsarévitch Alexis, est hémophile. La détresse de la famille impériale face à la maladie d'Alexis favorise l'influence du mystique Raspoutine, contribuant à la chute des Romanov en 1917.
    \item \textbf{Espagne :} Sa petite-fille Victoria-Eugénie (fille de Béatrice) épouse le roi Alphonse XIII. Deux de leurs fils (Alphonse et Gonzalve) sont hémophiles ; un troisième (Juan, père de Juan Carlos Ier) est sain.
    \item \textbf{Allemagne :} Sa fille Alice transmet à sa fille Irène, qui épouse le prince Henri de Prusse ; leurs fils Waldemar et Henri sont hémophiles.
\end{itemize}
Cette « expérience historique » a fourni à Thomas Hunt Morgan les données humaines massives confirmant que les gènes sont portés par les chromosomes (théorie chromosomique de l'hérédité, Prix Nobel 1933). Au Cameroun, bien que moins médiatisée, l'hémophilie existe ; la Fédération Camerounaise des Hémophiles (FECAH) œuvre pour le diagnostic (dosage facteur VIII/IX) et la prise en charge (concentrés de facteurs, prophylaxie).
\end{savaistu}

\courssub{Aperçus d'autres anomalies géniques au programme}

\noindent Le programme cite quatre autres anomalies géniques. Deux suffisent pour l'examen, mais les connaître toutes renforce votre culture scientifique. Toutes illustrent le même principe : une \cle{mutation ponctuelle} (substitution, insertion, délétion) dans un gène unique altère la protéine codée.

\begin{exemplebox}[Phénylcétonurie (PKU) — Autosomique récessive, Chromosome 12]
\begin{itemize}
    \item \textbf{Gène / Protéine :} \emph{PAH} (Phénylalanine Hydroxylase). Enzyme hépatique convertissant la phénylalanine (Phe, acide aminé essentiel) en tyrosine (Tyr).
    \item \textbf{Mécanisme :} Mutation $\rightarrow$ enzyme absente/inefficace $\rightarrow$ accumulation toxique de phénylalanine dans le sang et le cerveau $\rightarrow$ retard mental sévère, troubles neurologiques, hypopigmentation (cheveux blonds, yeux bleus), odeur de « souris » des urines (phénylacétate).
    \item \textbf{Dépistage néonatal (Test de Guthrie) :} Effectué au 3$^{\text{e}}$ jour de vie (talon). Mesure de la Phe sanguine. Au Cameroun, le programme national de dépistage néonatal (PNDN) commence à l'intégrer dans les grandes villes (Yaoundé, Douala).
    \item \textbf{Prise en charge :} \textbf{Régime hypoprotidique strict} (pauvre en phénylalanine) dès la naissance et à vie (ou au moins jusqu'à la puberté/adolescence). Supplémentation en tyrosine (devenue essentielle). Aspartame (édulcorant libérant de la Phe) \textbf{interdit}.
    \item \textbf{Message clé :} Une anomalie génétique grave devient \textbf{maîtrisable par l'environnement (nutrition)} si dépistée tôt. Génotype $\neq$ Destin immuable.
\end{itemize}
\end{exemplebox}

\begin{exemplebox}[Myopathie de Duchenne (DMD) — Récessive liée à l'X, Chromosome Xp21]
\begin{itemize}
    \item \textbf{Gène / Protéine :} \emph{DMD} (plus grand gène humain, 2,4 Mb). Code pour la \cle{dystrophine}, protéine cytosquelettique stabilisant la membrane des fibres musculaires (complexe dystrophine-glycoprotéines).
    \item \textbf{Mécanisme :} Délétions massives (65 \% des cas) ou mutations ponctuelles $\rightarrow$ absence de dystrophine fonctionnelle $\rightarrow$ fragilité membranaire $\rightarrow$ nécrose musculaire, remplacement par tissu fibreux/adipeux.
    \item \textbf{Phénotype :} Atteinte \textbf{quasiment exclusive des garçons} (transmission X récessive). Début marche (3--5 ans) : démarche dandinante, signe de Gowers (main sur cuisses pour se relever), pseudohypertrophie mollets. Évolution fatale (insuffisance cardio-respiratoire) vers 20--30 ans (améliorée par corticoïdes, ventilation non invasive, cardiologie).
    \item \textbf{Femmes conductrices :} Peuvent présenter une cardiomyopathie dilatée ou une faiblesse musculaire modérée (lyonisation défavorable).
    \item \textbf{Dépistage prénatal / Diagnostic préimplantatoire (DPI) :} Possible si mutation familiale connue.
\end{itemize}
\end{exemplebox}

\begin{exemplebox}[Mucoviscidose (Fibrose Kystique) — Autosomique récessive, Chromosome 7q31]
\begin{itemize}
    \item \textbf{Gène / Protéine :} \emph{CFTR} (Cystic Fibrosis Transmembrane conductance Regulator). Canal chlore (Cl$^-$) épithélial régulé par l'AMPc.
    \item \textbf{Mutation fréquente :} $\Delta$F508 (délétion 3 paires de bases $\rightarrow$ perte Phénylalanine 508) $\rightarrow$ protéine mal repliée, dégradée, n'atteint pas la membrane.
    \item \textbf{Conséquence physiologique :} Défaut d'hydratation des sécrétions $\rightarrow$ \textbf{mucus épais, visqueux, collant}.
    \item \textbf{Atteintes multiviscérales :}
    \begin{itemize}
        \item \textbf{Poumons :} Encombrement, infections bactériennes chroniques (\emph{Staphylococcus aureus}, \emph{Pseudomonas aeruginosa}), bronchiectasies $\rightarrow$ insuffisance respiratoire (cause majeure de décès).
        \item \textbf{Pancréas :} Obstruction canaux excréteurs $\rightarrow$ insuffisance exocrine (malabsorption lipidique, stéatorrhée, déficit vitamines liposolubles A, D, E, K) $\rightarrow$ traitement par extraits pancréatiques à chaque repas.
        \item \textbf{Appareil reproducteur :} Azoospermie obstructive chez l'homme (absence canaux déférents) ; mucus cervical épais chez la femme (fertilité réduite).
        \item \textbf{Test de la sueur :} Chlorures $> 60\ \text{mmol/L}$ (diagnostic de référence).
    \end{itemize}
    \item \textbf{Thérapies innovantes :} Modulateurs CFTR (ex. Trikafta\textsuperscript{\textregistered} : correcteur + potentiateurs) transforment le pronostic pour les mutations éligibles (dont $\Delta$F508). Accès encore limité au Cameroun (coût élevé).
\end{itemize}
\end{exemplebox}

\begin{aretenir}[Point commun fondamental : l'origine mutationnelle]
\begin{center}
\textbf{Toutes ces anomalies (drépanocytose, hémophilie, PKU, myopathie de Duchenne, mucoviscidose) partagent une étiologie unique :}
\end{center}
\begin{itemize}
    \item Une \textbf{mutation géniques} (altération de la séquence d'ADN d'un gène unique).
    \item Une \textbf{transmission héréditaire} suivant les lois de Mendel (autosomique récessive, dominante, ou liée à l'X).
    \item \textbf{Aucune cause surnaturelle, mystique, morale ou comportementale.} Ni la faute des parents, ni un « mauvais sort » (sorcellerie, « ndjobi », « kindoki »), ni une alimentation interdite pendant la grossesse, ni un péché des ancêtres n'en sont la cause.
    \item La \cle{variabilité génétique} (mutations + brassage) est le moteur de l'évolution ; les allèles pathologiques en sont un sous-produit statistiquement inévitable dans toute population.
\end{itemize}
Comprendre cela, c'est la base scientifique de la \cle{lutte contre les stigmatisations} familiales et sociales.
\end{aretenir}

\courssub{Éradication des préjugés : de la génétique à l'éducation et à la solidarité}

\noindent En tant que futurs citoyens, professionnels de santé, enseignants ou décideurs, vous avez un rôle clé. La science ne suffit pas si elle ne s'accompagne pas d'une éthique de la communication et de l'action.

\begin{methode}[Sensibiliser efficacement : démarche en 4 temps]
\begin{enumerate}
    \item \textbf{Expliquer le mécanisme biologique simplement :} « C'est une erreur de copie dans un manuel d'instructions (l'ADN) qui fait qu'une pièce de la machine (protéine) ne marche pas. Cela arrive au hasard lors de la formation des spermatozoïdes ou ovules. »
    \item \textbf{Démontrer l'innocence des parents :} Montrer l'échiquier de croisement. « Regardez : le père et la mère sont sains. Ils ont chacun un allèle normal et un allèle muté. Le hasard a fait que l'enfant a reçu les deux allèles mutés. Ce n'est la faute de personne. »
    \item \textbf{Valoriser la prise en charge médicale :} « Aujourd'hui, pour la drépanocytose (hydroxyurée, vaccination, suivi), l'hémophilie (facteurs recombinants, prophylaxie), la PKU (régime), la mucoviscidose (modulateurs CFTR, kiné respiratoire), la myopathie (corticoïdes, rééducation), \textbf{des traitements existent}. L'enfant peut grandir, aller à l'école, travailler. »
    \item \textbf{Promouvoir le conseil génétique (Genetic Counseling) :} Avant le mariage ou la grossesse, informer les couples à risque (ex. porteurs d'hémoglobine S, conductrices d'hémophilie) sur les risques de récidive (25 \%, 50 \%) et les options (diagnostic prénatal, DPI, adoption). Au Cameroun, le Centre de Génétique Humaine (Faculté de Médecine Yaoundé) et les centres de prise en charge de la drépanocytose (Hôpital Laquintinie Douala, Hôpital Central Yaoundé, Hôpital Général Douala) offrent ce service.
\end{enumerate}
\end{methode}

\begin{attention}[Erreurs de langage à bannir dans le discours public et familial]
\begin{itemize}
    \item \textbf{« Enfant sorcier / enfant mystère / enfant maudit. »} $\rightarrow$ Remplacer par : « Enfant atteint de drépanocytose / d'hémophilie / de mucoviscidose. »
    \item \textbf{« La femme a mangé ceci / a fait cela / a été ensorcelée. »} $\rightarrow$ Remplacer par : « La mutation est survenue aléatoirement. »
    \item \textbf{« C'est la famille de la mère / du père qui a ce sang. »} $\rightarrow$ Préciser : « L'allèle muté circule dans cette lignée. Le conseil génétique permet de le savoir. »
    \item \textbf{Cacher l'enfant, le négliger, le confier à un tradipraticien exclusif.} $\rightarrow$ Danger vital. Orienter vers le service de pédiatrie / hématologie / pneumologie / génétique le plus proche.
\end{itemize}
\end{attention}

\begin{exoresolu}[Analyse d'une situation familiale : conseil génétique]
\textbf{Énoncé :} Un couple consulte. L'homme est hémophile A ($X^h Y$). La femme est phénotypiquement saine, sans antécédent familial connu. Ils veulent connaître le risque pour leurs futurs enfants.
\textbf{Solution :}
\begin{enumerate}
    \item \textbf{Génotypes parentaux :} Père = $X^h Y$. Mère = $X^H X^H$ (probabilité très élevée car pas d'antécédent, fréquence allélique $q \approx 1/5000$ pour hémophilie A $\rightarrow$ fréquence conductrices $2pq \approx 1/2500$).
    \item \textbf{Croisement :} $X^h Y \times X^H X^H$.
    \item \textbf{Gamètes :} Père : $X^h$ ou $Y$. Mère : uniquement $X^H$.
    \item \textbf{Descendance :}
    \begin{itemize}
        \item Filles : 100 \% $X^H X^h$ $\rightarrow$ \textbf{Toutes conductrices}, phénotypiquement saines (sauf lyonisation défavorable rare).
        \item Garçons : 100 \% $X^H Y$ $\rightarrow$ \textbf{Tous sains}, non porteurs.
    \end{itemize}
    \item \textbf{Conclusion pour le couple :} \textbf{Aucun enfant ne sera hémophile.} Toutes les filles seront conductrices. Elles devront être informées plus tard pour leur propre projet parental (risque 50 \% de fils hémophiles si elles ont des enfants avec un homme sain).
    \item \textbf{Nuance importante :} Si la mère \emph{était} conductrice ($X^H X^h$) sans le savoir, le risque serait : 50 \% fils hémophiles, 50 \% filles conductrices. D'où l'intérêt du dosage du facteur VIII/IX chez la femme (taux $\approx 50\ \%$ suggère conductrice) ou du test moléculaire.
\end{enumerate}
\end{exoresolu}

\noindent En résumé, la génétique moléculaire et formelle nous donne les clés pour \textbf{nommer}, \textbf{comprendre}, \textbf{prévoir} et \textbf{soigner}. Elle nous donne aussi le devoir moral de \textbf{protéger} les familles contre l'ignorance qui engendre la culpabilité, le rejet et la souffrance évitable. La science est une lampe ; à nous de la porter dans la cité.

\coursec{IV. Anomalies chromosomiques : trisomie 21, Turner, Klinefelter}

Après avoir exploré les anomalies géniques — où la séquence d'un seul gène est altérée — nous changeons d'échelle. Les anomalies chromosomiques concernent l'ensemble du génome : ce sont des modifications du \textbf{nombre} ou de la \textbf{structure} des chromosomes, visibles au microscope optique sur un caryotype. Elles touchent des centaines, voire des milliers de gènes simultanément. Au Cameroun comme ailleurs, ces anomalies (trisomie 21, syndrome de Turner, syndrome de Klinefelter) sont souvent sources de stigmatisation familiale. Comprendre leur origine strictement biologique, aléatoire et non héréditaire dans l'immense majorité des cas est la première arme scientifique contre les préjugés.

\courssub{Le caryotype : outil d'identification des anomalies chromosomiques}

\begin{definition}[Caryotype et aneuploïdie]
Le \cle{caryotype} est le portrait chromosomique d'un individu : il représente l'ensemble de ses chromosomes, classés par paires homologues selon leur taille décroissante, la position du centromère et leur profil de bandage. Chez l'humain, l'espèce est diploïde avec $2n = 46$ chromosomes : 22 paires d'autosomes (numérotées de 1 à 22) et une paire de chromosomes sexuels (XX chez la femme, XY chez l'homme). Une \cle{aneuploïdie} est une anomalie du nombre de chromosomes par rapport à la formule diploïde normale : une \cle{trisomie} ($2n+1$, trois exemplaires d'un chromosome) ou une \cle{monosomie} ($2n-1$, un seul exemplaire).
\end{definition}

L'analyse d'un caryotype suit une procédure standardisée : culture de lymphocytes (sang périphérique) ou d'amniocytes (liquide amniotique), blocage en métaphase par la colchicine, étalement chromosomique, coloration (bande G) et photographies. Le classement permet de repérer tout chromosome en excès ou en défaut.

\begin{popfigure}
\begin{tikzpicture}[scale=0.65, every node/.style={font=\small}]
% Normal karyotype - 23 pairs arranged in 7 columns
\def\chrwidth{0.35}
\def\chrheight{1.2}
\def\pairgap{0.6}
\def\rowgap{1.0}
% Draw 22 autosome pairs
\foreach \i in {1,...,22} {
  \pgfmathsetmacro{\col}{mod(\i-1,7)}
  \pgfmathsetmacro{\row}{floor((\i-1)/7)}
  \pgfmathsetmacro{\x}{\col * (\chrwidth*2 + \pairgap)}
  \pgfmathsetmacro{\y}{-\row * (\chrheight + \rowgap)}
  % Pair label
  \node at (\x+\chrwidth+\pairgap/2, \y+0.3) {\textbf{\i}};
  % Two homologues
  \draw[popdark, thick, fill=popblue!10] (\x, \y) rectangle ++(\chrwidth, \chrheight);
  \draw[popdark, thick, fill=popblue!10] (\x+\chrwidth+\pairgap, \y) rectangle ++(\chrwidth, \chrheight);
  % Centromere
  \pgfmathsetmacro{\cent}{0.4*\chrheight}
  \fill[popdark] (\x, \y+\cent) rectangle ++(\chrwidth, 0.08*\chrheight);
  \fill[popdark] (\x+\chrwidth+\pairgap, \y+\cent) rectangle ++(\chrwidth, 0.08*\chrheight);
}
% Sex chromosomes XX (female) placed next to pair 22 (row 3, col 1)
\pgfmathsetmacro{\col}{1}
\pgfmathsetmacro{\row}{3}
\pgfmathsetmacro{\x}{\col * (\chrwidth*2 + \pairgap)}
\pgfmathsetmacro{\y}{-\row * (\chrheight + \rowgap)}
\node at (\x+\chrwidth+\pairgap/2, \y+0.3) {\textbf{XX}};
\draw[popdark, thick, fill=poppink!20] (\x, \y) rectangle ++(\chrwidth, \chrheight);
\draw[popdark, thick, fill=poppink!20] (\x+\chrwidth+\pairgap, \y) rectangle ++(\chrwidth, \chrheight);
\fill[popdark] (\x, \y+0.5*\chrheight) rectangle ++(\chrwidth, 0.08*\chrheight);
\fill[popdark] (\x+\chrwidth+\pairgap, \y+0.5*\chrheight) rectangle ++(\chrwidth, 0.08*\chrheight);

% Title
\node[anchor=south, font=\bfseries] at (7.5, 2.5) {Caryotype normal 46,XX};
\end{tikzpicture}
\legende{Caryotype humain normal (46,XX) : 22 paires d'autosomes classés par taille décroissante (paires 1 à 22) et une paire de chromosomes sexuels XX. Chaque chromosome est représenté par ses deux chromatides sœurs jointes au centromère (zone resserrée). La paire 22 et les chromosomes sexuels occupent la 4\textsuperscript{e} rangée.}
\end{popfigure}

\begin{popfigure}
\begin{tikzpicture}[scale=0.65, every node/.style={font=\small}]
% Trisomy 21 karyotype - highlight pair 21
\def\chrwidth{0.35}
\def\chrheight{1.2}
\def\pairgap{0.6}
\def\rowgap{1.0}
\foreach \i in {1,...,22} {
  \pgfmathsetmacro{\col}{mod(\i-1,7)}
  \pgfmathsetmacro{\row}{floor((\i-1)/7)}
  \pgfmathsetmacro{\x}{\col * (\chrwidth*2 + \pairgap)}
  \pgfmathsetmacro{\y}{-\row * (\chrheight + \rowgap)}
  \node at (\x+\chrwidth+\pairgap/2, \y+0.3) {\textbf{\i}};
  \draw[popdark, thick, fill=popblue!10] (\x, \y) rectangle ++(\chrwidth, \chrheight);
  \draw[popdark, thick, fill=popblue!10] (\x+\chrwidth+\pairgap, \y) rectangle ++(\chrwidth, \chrheight);
  \pgfmathsetmacro{\cent}{0.4*\chrheight}
  \fill[popdark] (\x, \y+\cent) rectangle ++(\chrwidth, 0.08*\chrheight);
  \fill[popdark] (\x+\chrwidth+\pairgap, \y+\cent) rectangle ++(\chrwidth, 0.08*\chrheight);
  % Highlight pair 21 with red border
  \ifnum\i=21
    \draw[poporange, very thick] (\x-0.05, \y-0.05) rectangle ++(\chrwidth*2+\pairgap+0.1, \chrheight+0.1);
    \node[poporange, anchor=south] at (\x+\chrwidth+\pairgap/2, \y+\chrheight+0.15) {\textbf{$\leftarrow$ Trisomie 21}};
  \fi
}
% Sex chromosomes XY (male) placed next to pair 22 (row 3, col 1)
\pgfmathsetmacro{\col}{1}
\pgfmathsetmacro{\row}{3}
\pgfmathsetmacro{\x}{\col * (\chrwidth*2 + \pairgap)}
\pgfmathsetmacro{\y}{-\row * (\chrheight + \rowgap)}
\node at (\x+\chrwidth+\pairgap/2, \y+0.3) {\textbf{XY}};
\draw[popdark, thick, fill=popblue!20] (\x, \y) rectangle ++(\chrwidth, \chrheight); % X
\draw[popdark, thick, fill=poporange!20] (\x+\chrwidth+\pairgap, \y) rectangle ++(\chrwidth, 0.7*\chrheight); % Y smaller
\fill[popdark] (\x, \y+0.5*\chrheight) rectangle ++(\chrwidth, 0.08*\chrheight);
\fill[popdark] (\x+\chrwidth+\pairgap, \y+0.35*\chrheight) rectangle ++(\chrwidth, 0.08*\chrheight);

% Title
\node[anchor=south, font=\bfseries] at (7.5, 2.5) {Caryotype trisomie 21 : 47,XY,+21};
\end{tikzpicture}
\legende{Caryotype d'un garçon trisomique 21 (47,XY,+21) : on observe 47 chromosomes au lieu de 46. La paire 21 (entourée en orange) comporte \textbf{trois} chromosomes au lieu de deux. La formule caryotypique s'écrit 47,XY,+21 (ou 47,XX,+21 pour une fille).}
\end{popfigure}

\begin{methode}[Identifier une anomalie chromosomique sur un caryotype]
\textbf{Étape 1 :} Compter le nombre total de chromosomes (46 normalement).\par
\textbf{Étape 2 :} Vérifier chaque paire d'autosomes (1 à 22) : y a-t-il 2 chromosomes par paire ? Repérer une trisomie (3 exemplaires) ou une monosomie (1 exemplaire).\par
\textbf{Étape 3 :} Examiner les chromosomes sexuels : XX (femme), XY (homme), X seul (45,X), XXY (47,XXY), XYY, etc.\par
\textbf{Étape 4 :} Écrire la formule caryotypique selon la nomenclature ISCN : nombre total, chromosomes sexuels, anomalie (ex. 47,XX,+21 ; 45,X ; 47,XXY).
\end{methode}

\courssub{Le syndrome de Down (trisomie 21) : signes, mécanisme et facteur d'âge maternel}

\begin{exemplebox}[Signes cliniques principaux du syndrome de Down]
Le syndrome de Down (trisomie 21) associe : \textbf{retard mental} (variable, souvent modéré), \textbf{dysmorphie faciale} (fentes palpébrales obliques, épicanthus, visage plat, langue proéminente), \textbf{hypotonie} musculaire, \textbf{ligne palmaire unique} (ligne simienne), \textbf{malformations cardiaques} fréquentes (communication interventriculaire, canal artériel persistant — $\approx 50\%$ des cas), troubles digestifs, risque accru de leucémie et maladie d'Alzheimer précoce. L'espérance de vie s'est considérablement améliorée (au-delà de 60 ans dans les pays à ressources médicales suffisantes) grâce à la prise en charge chirurgicale cardiaque et au suivi pluridisciplinaire.
\end{exemplebox}

La trisomie 21 résulte de la présence de \textbf{trois exemplaires du chromosome 21} au lieu de deux. Comment cet excès survient-il ? L'observation des caryotypes parentaux et la génétique formelle ont établi le mécanisme.

\begin{experience}[Origine méiotique de la trisomie 21 : la non-disjonction]
\textbf{Observation :} Chez $\approx 95\%$ des trisomiques 21, les parents ont un caryotype normal (46,XX et 46,XY). L'enfant a 47 chromosomes dont trois chromosomes 21. L'excès provient donc d'un \textbf{gamète anormal} (à 24 chromosomes au lieu de 23) qui a fusionné avec un gamète normal (23 chromosomes) $\rightarrow$ zygote à 47 chromosomes.\par
\textbf{Hypothèse :} Une erreur de ségrégation chromosomique pendant la méiose (non-disjonction) a produit un gamète disomique pour le chromosome 21 (deux exemplaires au lieu d'un).\par
\textbf{Expérience / Document :} Suivi cytogénétique des ovocytes et spermatozoïdes ; analyse de l'origine parentale par polymorphismes ADN (marqueurs microsatellites sur le chromosome 21).\par
\textbf{Interprétation :} Dans $\approx 90\%$ des cas, la non-disjonction a lieu \textbf{au cours de la méiose I maternelle} (séparation des chromosomes homologues) ; dans $\approx 5\%$ des cas, en méiose II maternelle (séparation des chromatides sœurs) ; dans $\approx 5\%$ des cas, en méiose paternelle. Le gamète résultant porte deux chromosomes 21 (disomie 21) ou aucun (nullisomie 21, létale).\par
\textbf{Conclusion :} La trisomie 21 résulte d'une \cle{non-disjonction chromosomique} lors de la méiose, accident aléatoire non héréditaire.
\end{experience}

\begin{popfigure}
\begin{tikzpicture}[scale=0.85, every node/.style={font=\small}]
% Normal Meiosis I and II
\def\chrW{0.5}
\def\chrH{1.0}
\def\gap{1.5}
% Title
\node[font=\bfseries] at (-2.5, 5.5) {Méiose NORMALE (gamétogenèse féminine)};
% 2n cell (prophase I) - replicated chromosomes
\draw[popdark, thick, fill=popblue!10] (-3.5, 4.5) rectangle ++(\chrW, \chrH);
\draw[popdark, thick, fill=popblue!10] (-3.5+\chrW+0.1, 4.5) rectangle ++(\chrW, \chrH);
\node at (-3.5+\chrW+0.05, 4.8) {2n};
\node at (-3.5+\chrW+0.05, 4.2) {Chr 21 répliqué};
% Meiosis I -> two n cells (each with replicated chr)
\draw[->, thick, popdark] (-3.5+\chrW+0.05, 4.0) -- ++(0, -0.6);
\draw[popdark, thick, fill=popblue!10] (-4.5, 3.0) rectangle ++(\chrW, \chrH);
\draw[popdark, thick, fill=popblue!10] (-2.5, 3.0) rectangle ++(\chrW, \chrH);
\node at (-4.5+\chrW/2, 3.3) {n};
\node at (-2.5+\chrW/2, 3.3) {n};
% Meiosis II -> four n cells (each with single chr)
\draw[->, thick, popdark] (-4.5+\chrW/2, 2.5) -- ++(0, -0.5);
\draw[->, thick, popdark] (-2.5+\chrW/2, 2.5) -- ++(0, -0.5);
\foreach \x in {-4.5, -2.5} {
  \draw[popdark, thick, fill=popgreen!20] (\x, 1.5) rectangle ++(\chrW, 0.6*\chrH);
  \draw[popdark, thick, fill=popgreen!20] (\x+\chrW+0.1, 1.5) rectangle ++(\chrW, 0.6*\chrH);
}
\node at (-3.5, 1.0) {4 gamètes normaux (n)};

% NON-DISJUNCTION Meiosis I
\node[font=\bfseries] at (5.5, 5.5) {NON-DISJONCTION EN MÉIOSE I};
\draw[popdark, thick, fill=popblue!10] (4, 4.5) rectangle ++(\chrW, \chrH);
\draw[popdark, thick, fill=popblue!10] (4+\chrW+0.1, 4.5) rectangle ++(\chrW, \chrH);
\node at (4+\chrW+0.05, 4.8) {2n};
\draw[->, thick, popdark] (4+\chrW+0.05, 4.0) -- ++(0, -0.6);
% Two cells: one with 2 chr21, one with 0
\draw[popdark, thick, fill=poporange!30] (3, 3.0) rectangle ++(\chrW, \chrH);
\draw[popdark, thick, fill=poporange!30] (3+\chrW+0.1, 3.0) rectangle ++(\chrW, \chrH);
\node[poporange] at (3+\chrW+0.05, 3.3) {n+1 (disomique)};
\draw[popdark, thick, fill=popgray!30] (5.5, 3.0) rectangle ++(\chrW, \chrH);
\node[popgray] at (5.5+\chrW/2, 3.3) {n-1 (nullisomique)};
% Meiosis II
\draw[->, thick, popdark] (3+\chrW+0.05, 2.5) -- ++(0, -0.5);
\draw[->, thick, popdark] (5.5+\chrW/2, 2.5) -- ++(0, -0.5);
% n+1 cell divides into two n+1
\draw[popdark, thick, fill=poporange!30] (3, 1.5) rectangle ++(\chrW, 0.6*\chrH);
\draw[popdark, thick, fill=poporange!30] (3+\chrW+0.1, 1.5) rectangle ++(\chrW, 0.6*\chrH);
\node[poporange] at (3+\chrW+0.05, 1.0) {2 gamètes n+1};
% n-1 cell divides into two n-1 (lethal)
\draw[popdark, thick, fill=popgray!30] (5.5, 1.5) rectangle ++(\chrW, 0.6*\chrH);
\draw[popdark, thick, fill=popgray!30] (5.5+\chrW+0.1, 1.5) rectangle ++(\chrW, 0.6*\chrH);
\node[popgray] at (5.5+\chrW+0.05, 1.0) {2 gamètes n-1 (létaux)};

% Legend
\node[anchor=west] at (-5.5, -0.5) {\textbf{Légende :} \textcolor{popblue!80!black}{$\blacksquare$} Chromosome 21 répliqué (2 chromatides) \quad \textcolor{popgreen!60!black}{$\blacksquare$} Chromosome 21 simple (1 chromatide) \quad \textcolor{poporange}{$\blacksquare$} Excès (disomie 21) \quad \textcolor{popgray}{$\blacksquare$} Déficit (nullisomie 21)};
\end{tikzpicture}
\legende{Schéma de la non-disjonction du chromosome 21 en méiose I (gauche : méiose normale ; droite : non-disjonction). En méiose I, les chromosomes homologues ne se séparent pas : une cellule fille reçoit les deux chromosomes 21 (n+1), l'autre n'en reçoit aucun (n-1). Après méiose II, on obtient 2 gamètes disomiques (n+1) et 2 gamètes nullisomiques (n-1, létaux). La fécondation d'un gamète n+1 par un gamète normal (n) donne un zygote trisomique (2n+1).}
\end{popfigure}

Un fait épidémiologique majeur guide le conseil génétique et le dépistage prénatal : la fréquence de la trisomie 21 à la naissance augmente fortement avec l'âge de la mère.

\begin{popfigure}
\begin{tikzpicture}
\begin{axis}[
  width=\linewidth,
  height=8cm,
  xlabel={Âge de la mère (ans)},
  ylabel={Fréquence des naissances trisomiques 21 (pour 1000 naissances)},
  xmin=15, xmax=50,
  ymin=0, ymax=35,
  xtick={15,20,25,30,35,40,45,50},
  ytick={0,5,10,15,20,25,30,35},
  grid=major,
  grid style={popline},
  axis line style={popdark},
  tick style={popdark},
  label style={font=\small, popdark},
  tick label style={font=\small, popdark},
]
% Data points approximate: 20y~0.7, 25y~1, 30y~1.5, 35y~3.5, 40y~10, 45y~30 per 1000
\addplot[popcurve, thick, mark=*, mark size=2.5, mark options={fill=poporange}] coordinates {
  (15,0.5) (20,0.7) (25,1.0) (30,1.5) (35,3.5) (37,5.5) (40,10.0) (42,15.0) (45,30.0) (48,35.0)
};
\addplot[popblue, dashed, thick] coordinates { (35,0) (35,35) };
\node[popblue, anchor=south west, font=\small\bfseries] at (35.5, 20) {Seuil de dépistage renforcé (35 ans)};
\end{axis}
\end{tikzpicture}
\legende{Évolution de la fréquence de la trisomie 21 à la naissance en fonction de l'âge maternel (données épidémiologiques de référence). La courbe montre une augmentation exponentielle après 35 ans. À 20 ans : $\approx 1/1500$ ; à 35 ans : $\approx 1/350$ ; à 40 ans : $\approx 1/100$ ; à 45 ans : $\approx 1/30$. Ce risque justifie la proposition systématique de diagnostic prénatal (amniocentèse, choriocentèse ou ADN fœtal circulant) aux femmes de 35 ans et plus.}
\end{popfigure}

\begin{exemplebox}[Risque de trisomie 21 selon l'âge maternel — données chiffrées]
\begin{center}
\begin{tabularx}{\linewidth}{|c|>{\centering\arraybackslash}X|>{\centering\arraybackslash}X|}\hline
\rowcolor{popblueL}
\textbf{Âge maternel} & \textbf{Risque approximatif à la naissance} & \textbf{1 naissance sur...} \\ \hline
20 ans & 0,07 \% & 1 / 1\,500 \\ \hline
25 ans & 0,10 \% & 1 / 1\,000 \\ \hline
30 ans & 0,15 \% & 1 / 670 \\ \hline
35 ans & 0,30 \% & 1 / 330 \\ \hline
37 ans & 0,55 \% & 1 / 180 \\ \hline
40 ans & 1,0 \% & 1 / 100 \\ \hline
42 ans & 1,5 \% & 1 / 65 \\ \hline
45 ans & 3,0 \% & 1 / 33 \\ \hline
\end{tabularx}
\end{center}
Ces chiffres expliquent pourquoi le suivi obstétrical au Cameroun (comme dans les recommandations internationales) propose un \textbf{dépistage combiné du premier trimestre} (marqueurs sériques PAPP-A, $\beta$-hCG libre + clarté nucale fœtale par échographie) à toutes les femmes, et un \textbf{diagnostic prénatal invasif} (caryotype sur liquide amniotique ou villosités choriales) ou \textbf{test ADN fœtal circulant} (NIPT) aux femmes de $\ge 35$ ans ou à risque élevé.
\end{exemplebox}

\begin{attention}[Ne pas confondre anomalie chromosomique et anomalie génique]
\begin{itemize}
\item \textbf{Anomalie chromosomique} : visible au microscope sur caryotype (ex. trisomie 21, 47,XX,+21). Concerne un chromosome entier $\rightarrow$ des centaines de gènes affectés. Phénotype souvent complexe (syndrome polymalformatif).
\item \textbf{Anomalie génique (mutation ponctuelle)} : invisible au microscope, détectable seulement par séquençage de l'ADN (ex. drépanocytose : mutation ponctuelle $A \to T$ au codon 6 du gène $\beta$-globine). Concerne un seul gène. Phénotype souvent plus ciblé.
\end{itemize}
\textbf{Piège fréquent au Bac :} écrire « mutation du chromosome 21 » pour la trisomie 21. \textbf{Non !} C'est une \textbf{anomalie de nombre} (non-disjonction), pas une mutation de la séquence d'ADN du chromosome 21 (qui reste structurellement normal dans la trisomie 21 libre).
\end{attention}

\courssub{Autres anomalies des chromosomes sexuels : syndromes de Turner et de Klinefelter}

Les chromosomes sexuels (X et Y) sont aussi sujets aux non-disjonctions. Contrairement aux autosomes, les anomalies de nombre des chromosomes sexuels sont souvent viables et compatibles avec une vie longue, car l'inactivation du X (lyonisation) compense le dosage génique.

\begin{propriete}[Syndrome de Turner (monosomie X) — 45,X]
\textbf{Caryotype :} 45 chromosomes au total, un seul chromosome X (monosomie X). Formule : 45,X.\par
\textbf{Origine :} Non-disjonction des chromosomes sexuels lors de la méiose paternelle (perte du Y) ou maternelle (perte d'un X) $\rightarrow$ gamète nullique pour le chromosome sexuel (0) + gamète normal (X) $\rightarrow$ zygote 45,X. Parfois mosaïcisme 45,X/46,XX.\par
\textbf{Phénotype :} \textbf{Femme phénotypique} (gonades indifférenciées $\rightarrow$ ovaires), mais \textbf{ovaires dysgénétiques} (streak gonads) $\rightarrow$ absence de puberté spontanée, infertilité (sauf don d'ovocytes), petite taille ($\approx \SI{145}{\centi\meter}--\SI{150}{\centi\meter}$ sans traitement), plis cutanés au cou (ptyérygion), malformations cardiaques (coarctation de l'aorte), rénales, troubles auditifs. Intelligence normale. \textbf{Prise en charge :} hormone de croissance (enfance), œstrogènes/progestatifs (puberté), suivi cardiaque/rénal.
\end{propriete}

\begin{propriete}[Syndrome de Klinefelter — 47,XXY]
\textbf{Caryotype :} 47 chromosomes, deux X et un Y. Formule : 47,XXY (parfois mosaïques 46,XY/47,XXY).\par
\textbf{Origine :} Non-disjonction des chromosomes sexuels en méiose I ou II maternelle (XX) ou paternelle (XY) $\rightarrow$ gamète disomique (XX ou XY) + gamète normal $\rightarrow$ zygote 47,XXY.\par
\textbf{Phénotype :} \textbf{Homme phénotypique} (présence du Y $\rightarrow$ déterminisme testiculaire), mais \textbf{testicules petits et fermes}, \textbf{azoospermie} (infertilité quasi constante), \textbf{hypogonadisme} (taux de testostérone bas, FSH/LH élevés), gynécomastie, taille grande (jambes longues), intelligence normale (parfois troubles du langage/apprentissage). \textbf{Prise en charge :} testostérone à vie (dès la puberté), soutien psychopédagogique, techniques de procréation assistée (TESE/ICSI) possibles chez certains mosaïques.
\end{propriete}

\begin{savaistu}[Pourquoi la trisomie 21 est-elle la seule trisomie autosomique viable à long terme ?]
Les chromosomes 13, 18, 21 sont les plus petits autosomes. Pourtant, les trisomies 13 (syndrome de Patau) et 18 (syndrome d'Edwards) entraînent des malformations majeures (cerveau, cœur, reins) et un décès dans les premiers mois de vie dans $> 90\%$ des cas. Le chromosome 21 est le \textbf{plus petit autosome} ($\approx \SI{48}{Mb}$, $\approx 200-300$ gènes) : son excès en trois exemplaires déséquilibre moins le transcriptome global que les chromosomes 13 ou 18 (plus gros, plus de gènes). De plus, certains gènes du chromosome 21 échappent partiellement à la régulation de dosage. C'est pourquoi la trisomie 21 est compatible avec une vie de plusieurs décennies, alors que les autres trisomies autosomiques complètes sont létales in utero ou néonatalement.
\end{savaistu}

\courssub{Origine méiotique aléatoire : ni faute, ni hérédité — base de l'éradication des préjugés}

\begin{aretenir}[Points clés pour combattre les préjugés]
\begin{enumerate}
\item \textbf{Accident méiotique aléatoire :} La non-disjonction est un événement stochastique survenant lors de la formation des gamètes. Elle n'est \textbf{pas causée} par l'alimentation, le mode de vie, les « sorts », la colère des ancêtres, une malédiction familiale ou une faute des parents.
\item \textbf{Non héréditaire (généralement) :} Dans $> 95\%$ des trisomies 21, les parents ont un caryotype normal. Le risque de récurrence pour une grossesse ultérieure est faible ($\approx 1\%$ s'ajoutant au risque lié à l'âge maternel), sauf cas rares de translocation robertsonienne parentale (voir cours sur mutations chromosomiques structurelles) ou mosaïcisme gonadique.
\item \textbf{Aucun lien avec l'ethnie, la religion, le statut social :} La trisomie 21 touche toutes les populations humaines avec une fréquence de base similaire ($\approx 1/700$ naissances vivantes). Seule l'âge maternel module le risque.
\item \textbf{Au Cameroun :} La prise en charge médicale (chirurgie cardiaque à l'Hôpital Général de Yaoundé, à Laquintinie Douala, suivi pédiatrique, orthophonie, éducation inclusive) permet une intégration sociale réussie. Les associations de parents (ex. Association Camerounaise pour le Bien-être des Trisomiques) jouent un rôle clé dans la sensibilisation.
\end{enumerate}
\end{aretenir}

\begin{exoresolu}[Analyse d'un caryotype et conseil génétique]
\textbf{Énoncé :} Un couple consulte après la naissance d'un enfant présentant un syndrome de Down. Le caryotype de l'enfant est 47,XY,+21. Les caryotypes des parents sont normaux : père 46,XY ; mère 46,XX. La mère a 38 ans. Le couple craint une récidive pour une future grossesse et s'interroge sur une éventuelle responsabilité.\par
\tcblower
\textbf{Solution détaillée :}
\begin{enumerate}
\item \textbf{Identification de l'anomalie :} L'enfant a 47 chromosomes dont trois chromosomes 21 $\rightarrow$ trisomie 21 libre (non-disjonction), non une translocation.
\item \textbf{Origine :} Parents caryotypiquement normaux $\rightarrow$ anomalie \textbf{dé novo} survenue lors de la méiose (probablement méiose I maternelle vu l'âge de 38 ans).
\item \textbf{Risque de récidive :} Risque de base lié à l'âge maternel (à 38 ans $\approx 1/200$ à la naissance) + risque de récidive empirique faible ($\approx 1\%$) $\rightarrow$ risque total $\approx 1,5\%$. Ce n'est \textbf{pas} un risque de 25\% ou 50\% comme dans une maladie autosomique récessive.
\item \textbf{Conseil :} Rassurer le couple : « Ce n'est la faute de personne. C'est un accident chromosomique aléatoire. Pour une future grossesse, nous vous proposerons un diagnostic prénatal précoce (NIPT ou amniocentèse) pour vous informer, sans obligation d'interruption de grossesse. »
\end{enumerate}
\end{exoresolu}

\begin{attention}[Pièges à éviter dans la rédaction au Bac]
\begin{itemize}
\item \textbf{Écrire :} « La trisomie 21 est une mutation du chromosome 21. » \textbf{Faux.} C'est une anomalie de \textbf{nombre} (aneuploïdie), le chromosome 21 lui-même est structurellement normal.
\item \textbf{Écrire :} « Le syndrome de Turner touche les hommes (45,X). » \textbf{Faux.} 45,X est un phénotype \textbf{féminin} (absence de Y $\rightarrow$ pas de déterminisme testiculaire).
\item \textbf{Écrire :} « Le syndrome de Klinefelter (47,XXY) est une trisomie du chromosome X. » \textbf{Imprécis.} C'est une aneuploïdie des chromosomes sexuels ; on parle de « disomie X » chez un homme. Le terme « trisomie » est réservé aux autosomes.
\item \textbf{Oublier :} Préciser que la non-disjonction peut se produire en méiose I \textbf{ou} II, maternelle \textbf{ou} paternelle (mais majoritairement maternelle en méiose I pour la trisomie 21).
\end{itemize}
\end{attention}

\coursec{V. Brassage des allèles : diversité génétique et unicité des individus}

Nous venons d'étudier les anomalies chromosomiques (trisomie 21, syndromes de Turner et de Klinefelter) qui résultent d'erreurs de ségrégation lors de la méiose. Mais au-delà de ces accidents, la méiose accomplit une fonction fondamentale et parfaitement réglée : elle \emph{brasse} les allèles existants pour créer des combinaisons inédites. C'est ce brassage, allié à la fécondation aléatoire, qui explique pourquoi deux enfants d'un même couple (hors vrais jumeaux) sont toujours génétiquement distincts, et pourquoi un caractère récessif rare peut « réapparaître » soudainement dans une famille après des générations de silence. Comprenons les mécanismes de cette loterie génétique.

\courssub{Le problème : l'unicité génétique des frères et sœurs}

\begin{experience}[Observation : des frères et sœurs tous différents]
Dans une famille camerounaise de cinq enfants, on observe : l'aîné a le groupe sanguin A et est grand ; la cadette est du groupe O, plus petite, et drépanocytaire (SS) alors que ses parents sont asymptomatiques (AS) ; le troisième a le groupe B et les yeux clairs comme sa grand-mère paternelle ; les deux derniers sont jumeaux dizygotes (faux jumeaux) : l'un a la peau claire, l'autre très foncée.
\par\textbf{Question :} Comment un même couple peut-il engendrer une telle diversité phénotypique ? Les allèles des enfants sont-ils \emph{identiques} à ceux des parents ou \emph{répartis différemment} ?
\end{experience}

\begin{methode}[Démarche d'investigation : comparer les génotypes parentaux et filiaux]
\begin{enumerate}
    \item On génotype les parents et les enfants pour plusieurs gènes (ex. : \emph{ABO}, \emph{HBB} pour l'hémoglobine, \emph{MC1R} pour la pigmentation).
    \item On vérifie que chaque enfant possède \textbf{un allèle paternel et un allèle maternel} pour chaque gène (respect de la mendélien).
    \item On compare les \textbf{combinaisons d'allèles} entre les enfants : sont-elles identiques ou différentes ?
    \item On relie ces combinaisons au comportement des chromosomes lors de la méiose (ségrégation, crossing-over).
\end{enumerate}
\end{methode}

\begin{exoresolu}[Analyse de génotypes familiaux — Groupe ABO et HBB]
\textbf{Énoncé.} Un couple a pour génotypes : Père $[I^A//i]$ (groupe A) et $[A//S]$ (AS drépanocytose) ; Mère $[I^B//i]$ (groupe B) et $[A//S]$ (AS drépanocytose). Ils ont 4 enfants. Déterminer les génotypes possibles pour ces deux gènes chez les gamètes et chez les enfants.
\tcblower
\textbf{Solution.}
\begin{enumerate}
    \item \textbf{Gamètes paternels possibles} (loi de l'assortiment indépendant si gènes sur chromosomes différents ou éloignés) : 
    $[I^A; A]$, $[I^A; S]$, $[i; A]$, $[i; S]$ $\rightarrow$ 4 types.
    \item \textbf{Gamètes maternels possibles} : 
    $[I^B; A]$, $[I^B; S]$, $[i; A]$, $[i; S]$ $\rightarrow$ 4 types.
    \item \textbf{Zygotes possibles} : $4 \times 4 = 16$ combinaisons génétiques distinctes pour ces deux gènes seulement.
    \item \textbf{Exemples d'enfants possibles} :
    \begin{itemize}
        \item Enfant 1 : $[I^A//I^B]$ (groupe AB) et $[A//A]$ (AA sain) $\leftarrow$ gamètes $[I^A;A]$ (père) + $[I^B;A]$ (mère).
        \item Enfant 2 : $[i//i]$ (groupe O) et $[S//S]$ (SS drépanocytaire) $\leftarrow$ gamètes $[i;S]$ (père) + $[i;S]$ (mère).
        \item Enfant 3 : $[I^A//i]$ (groupe A) et $[A//S]$ (AS) $\leftarrow$ gamètes $[I^A;A]$ (père) + $[i;S]$ (mère).
    \end{itemize}
    \item \textbf{Conclusion :} Chaque enfant reçoit un \emph{mélange unique} d'allèles paternels et maternels. La drépanocytose (SS) « apparaît » chez l'enfant 2 bien que les parents soient sains, car l'allèle $S$ était présent à l'état hétérozygote chez chacun.
\end{enumerate}
\end{exoresolu}

\begin{aretenir}[Interprétation]
Les enfants ne sont pas des « photocopies » de leurs parents ni des « demi-photocopies » identiques entre eux. La méiose redistribue les allèles parentaux de manière \textbf{aléatoire} à chaque formation de gamète. C'est le \textbf{brassage génétique}.
\end{aretenir}

\courssub{Les deux niveaux du brassage méiotique}

La méiose réalise deux brassages successifs, sources de combinaisons nouvelles d'allèles \emph{existants}.

\begin{definition}[Brassage génétique (ou recombinaison génétique)]
Ensemble des mécanismes (brassage interchromosomique, brassage intrachromosomique, fécondation aléatoire) qui redistribuent aléatoirement les allèles portés par les chromosomes homologues lors de la méiose et de la fécondation, créant de nouvelles combinaisons alléliques sans modifier la séquence des allèles eux-mêmes.
\end{definition}

\textbf{2.1. Brassage interchromosomique (ou ségrégation indépendante)}\par

Lors de l'anaphase I, les paires de chromosomes homologues se séparent indépendamment les unes des autres. Pour chaque paire, le chromosome d'origine paternelle ou maternelle peut migrer vers l'un ou l'autre pôle.

\begin{popfigure}
\begin{tikzpicture}[scale=0.9, every node/.style={font=\small}]
    % Couleurs
    \colorlet{pat}{popblue!80!popdark}
    \colorlet{mat}{poppink!80!popdark}
    \colorlet{gam}{popgreen!70!popdark}
    
    % --- Niveau 1 : Cellule mère (2n=4) ---
    \node[draw, thick, ellipse, minimum width=3.5cm, minimum height=2.5cm, fill=popcream] (cell) at (0,3) {};
    \node[above=0.2cm of cell, font=\bfseries] {Cellule mère (2n=4)};
    
    % Chromosomes homologues : Paire 1 (Grand), Paire 2 (Petit)
    % Paire 1
    \draw[pat, line width=3pt, line cap=round] (-1.2, 3.5) -- (-1.2, 2.5) node[midway, left=2mm] {Chr 1\textsuperscript{pat}};
    \draw[mat, line width=3pt, line cap=round] (1.2, 3.5) -- (1.2, 2.5) node[midway, right=2mm] {Chr 1\textsuperscript{mat}};
    % Paire 2
    \draw[pat, line width=2pt, line cap=round] (-0.4, 3.8) .. controls (-0.4, 3.3) and (0.4, 3.3) .. (0.4, 2.8) node[midway, above=1mm] {Chr 2\textsuperscript{pat}};
    \draw[mat, line width=2pt, line cap=round] (0.4, 3.8) .. controls (0.4, 3.3) and (-0.4, 3.3) .. (-0.4, 2.8) node[midway, above=1mm] {Chr 2\textsuperscript{mat}};
    
    % Flèche méiose I
    \draw[->, thick, popdark] (0, 1.5) -- (0, 0.8) node[midway, right] {Méiose I (Anaphase I)};
    
    % --- Niveau 2 : Deux cellules haploïdes (n=2) ---
    % Cellule A
    \node[draw, thick, ellipse, minimum width=2.5cm, minimum height=1.5cm, fill=popblueL!30] (cA) at (-3, -0.5) {};
    \node[above=0.1cm of cA] {Cellule A};
    \draw[pat, line width=3pt] (-3.8, -0.2) -- (-3.8, -0.8);
    \draw[mat, line width=2pt] (-2.2, -0.2) .. controls (-2.2, -0.5) and (-1.8, -0.5) .. (-1.8, -0.8);
    \node[below=0.1cm of cA] {Chr 1\textsuperscript{pat} ; Chr 2\textsuperscript{mat}};
    
    % Cellule B
    \node[draw, thick, ellipse, minimum width=2.5cm, minimum height=1.5cm, fill=poppinkL!30] (cB) at (3, -0.5) {};
    \node[above=0.1cm of cB] {Cellule B};
    \draw[mat, line width=3pt] (3.8, -0.2) -- (3.8, -0.8);
    \draw[pat, line width=2pt] (2.2, -0.2) .. controls (2.2, -0.5) and (1.8, -0.5) .. (1.8, -0.8);
    \node[below=0.1cm of cB] {Chr 1\textsuperscript{mat} ; Chr 2\textsuperscript{pat}};
    
    % Flèche méiose II
    \draw[->, thick, popdark] (-3, -1.5) -- (-3, -2.2) node[midway, right] {Méiose II};
    \draw[->, thick, popdark] (3, -1.5) -- (3, -2.2) node[midway, right] {Méiose II};
    
    % --- Niveau 3 : 4 Gamètes ---
    % Gamètes de A
    \node[draw, circle, minimum size=1.2cm, fill=gam!20] (g1) at (-4.2, -3.2) {};
    \draw[pat, line width=2pt] (-4.2, -2.8) -- (-4.2, -3.6);
    \draw[mat, line width=1.5pt] (-4.6, -2.8) .. controls (-4.6, -3.2) and (-4.2, -3.2) .. (-4.2, -3.6);
    \node[below=0.1cm of g1] {G1: 1\textsuperscript{pat}, 2\textsuperscript{mat}};
    
    \node[draw, circle, minimum size=1.2cm, fill=gam!20] (g2) at (-1.8, -3.2) {};
    \draw[pat, line width=2pt] (-1.8, -2.8) -- (-1.8, -3.6);
    \draw[mat, line width=1.5pt] (-2.2, -2.8) .. controls (-2.2, -3.2) and (-1.8, -3.2) .. (-1.8, -3.6);
    \node[below=0.1cm of g2] {G2: 1\textsuperscript{pat}, 2\textsuperscript{mat}};
    
    % Gamètes de B
    \node[draw, circle, minimum size=1.2cm, fill=gam!20] (g3) at (1.8, -3.2) {};
    \draw[mat, line width=2pt] (1.8, -2.8) -- (1.8, -3.6);
    \draw[pat, line width=1.5pt] (2.2, -2.8) .. controls (2.2, -3.2) and (1.8, -3.2) .. (1.8, -3.6);
    \node[below=0.1cm of g3] {G3: 1\textsuperscript{mat}, 2\textsuperscript{pat}};
    
    \node[draw, circle, minimum size=1.2cm, fill=gam!20] (g4) at (4.2, -3.2) {};
    \draw[mat, line width=2pt] (4.2, -2.8) -- (4.2, -3.6);
    \draw[pat, line width=1.5pt] (3.8, -2.8) .. controls (3.8, -3.2) and (4.2, -3.2) .. (4.2, -3.6);
    \node[below=0.1cm of g4] {G4: 1\textsuperscript{mat}, 2\textsuperscript{pat}};
    
    % Légende
    \node[align=left, anchor=north west] at (-5, 4.2) {
        \textbf{Légende :}\\
        \textcolor{pat}{\rule{5mm}{2mm}} Chromosome paternel\\
        \textcolor{mat}{\rule{5mm}{2mm}} Chromosome maternel\\
        \textbf{2 paires} $\rightarrow$ \textbf{2\textsuperscript{2} = 4} types de gamètes
    };
\end{tikzpicture}
\legende{Brassage interchromosomique simplifié (2n=4). La ségrégation indépendante des deux paires de chromosomes homologues (Paire 1 : grand ; Paire 2 : petit) produit 4 combinaisons chromosomiques différentes dans les gamètes. Chez l'Homme (23 paires), cela donne $2^{23}$ combinaisons.}
\end{popfigure}

\begin{propriete}[Nombre de combinaisons interchromosomiques — Admis au Bac]
Pour une espèce à $n$ paires de chromosomes homologues, le brassage interchromosomique permet de former $2^n$ types de gamètes différents \emph{au niveau chromosomique}.
\par Chez l'Homme ($n=23$) : $2^{23} = 8\,388\,608 \approx 8,4 \times 10^6$ combinaisons chromosomiques possibles par individu.
\end{propriete}

\textbf{2.2. Brassage intrachromosomique (Crossing-over)}\par

Il se produit en prophase I (stade pachytène), entre chromatides non-sœurs de chromosomes homologues (tétravalent). Des échanges physiques de fragments d'ADN créent des chromosomes \textbf{recombinés} portant des associations d'allèles inédites sur une même molécule d'ADN.

\begin{popfigure}
\begin{tikzpicture}[scale=1.1, every node/.style={font=\small}]
    % Couleurs
    \colorlet{pat}{popblue}
    \colorlet{mat}{poppink}
    \colorlet{rec}{poporange}
    
    % Chromosomes homologues appariés (Tétravalent)
    % Chaque chromosome = 2 chromatides sœurs
    % Chromosome Paternel (haut) - Bleu
    \draw[pat, line width=4pt] (-4, 2) -- (4, 2) node[right] {Chr Paternel};
    \draw[pat, line width=4pt] (-4, 1.2) -- (4, 1.2);
    % Centromères paternels
    \fill[pat] (0, 2) rectangle (0.3, 1.2);
    
    % Chromosome Maternel (bas) - Rose
    \draw[mat, line width=4pt] (-4, -1.2) -- (4, -1.2) node[right] {Chr Maternel};
    \draw[mat, line width=4pt] (-4, -2) -- (4, -2);
    % Centromères maternels
    \fill[mat] (0, -1.2) rectangle (0.3, -2);
    
    % Allèles sur chromatides (représentés par petits carrés colorés)
    % Paternel haut : A - B
    \node[pat, fill=pat, minimum size=3pt] at (-2.5, 2) {}; \node[above] at (-2.5, 2) {$A$};
    \node[pat, fill=pat, minimum size=3pt] at (-1, 2) {}; \node[above] at (-1, 2) {$B$};
    % Paternel bas : A - B
    \node[pat, fill=pat, minimum size=3pt] at (-2.5, 1.2) {}; \node[below] at (-2.5, 1.2) {$A$};
    \node[pat, fill=pat, minimum size=3pt] at (-1, 1.2) {}; \node[below] at (-1, 1.2) {$B$};
    
    % Maternel haut : a - b
    \node[mat, fill=mat, minimum size=3pt] at (-2.5, -1.2) {}; \node[below] at (-2.5, -1.2) {$a$};
    \node[mat, fill=mat, minimum size=3pt] at (-1, -1.2) {}; \node[below] at (-1, -1.2) {$b$};
    % Maternel bas : a - b
    \node[mat, fill=mat, minimum size=3pt] at (-2.5, -2) {}; \node[below] at (-2.5, -2) {$a$};
    \node[mat, fill=mat, minimum size=3pt] at (-1, -2) {}; \node[below] at (-1, -2) {$b$};
    
    % Point de crossing-over (Chiasme)
    \draw[rec, very thick, decorate, decoration={zigzag, segment length=2mm, amplitude=1mm}] (-1.5, 1.2) -- (-1.5, -1.2);
    \node[rec, right] at (-1.5, 0) {\textbf{Chiasme (Crossing-over)}};
    
    % Flèche vers résultat
    \draw[->, very thick, popdark] (5, 0) -- (6.5, 0) node[midway, above] {Ségrégation (Anaphase I \& II)};
    
    % --- RÉSULTAT : 4 Chromatides filles ---
    \begin{scope}[xshift=8cm]
        % 1. Paternelle non recombinée (A-B)
        \draw[pat, line width=3pt] (0, 2.5) -- (3, 2.5);
        \node[pat, fill=pat, minimum size=3pt] at (0.5, 2.5) {}; \node[above] at (0.5, 2.5) {$A$};
        \node[pat, fill=pat, minimum size=3pt] at (1.5, 2.5) {}; \node[above] at (1.5, 2.5) {$B$};
        \node[left] at (0, 2.5) {1. };
        
        % 2. Paternelle recombinée (A-b) -> Orange pour partie recombinée
        \draw[pat, line width=3pt] (0, 1.2) -- (1, 1.2);
        \draw[rec, line width=3pt] (1, 1.2) -- (3, 1.2);
        \node[pat, fill=pat, minimum size=3pt] at (0.5, 1.2) {}; \node[above] at (0.5, 1.2) {$A$};
        \node[rec, fill=rec, minimum size=3pt] at (2, 1.2) {}; \node[above] at (2, 1.2) {$b$};
        \node[left] at (0, 1.2) {2. };
        
        % 3. Maternelle recombinée (a-B)
        \draw[mat, line width=3pt] (0, -0.1) -- (1, -0.1);
        \draw[rec, line width=3pt] (1, -0.1) -- (3, -0.1);
        \node[mat, fill=mat, minimum size=3pt] at (0.5, -0.1) {}; \node[below] at (0.5, -0.1) {$a$};
        \node[rec, fill=rec, minimum size=3pt] at (2, -0.1) {}; \node[below] at (2, -0.1) {$B$};
        \node[left] at (0, -0.1) {3. };
        
        % 4. Maternelle non recombinée (a-b)
        \draw[mat, line width=3pt] (0, -1.4) -- (3, -1.4);
        \node[mat, fill=mat, minimum size=3pt] at (0.5, -1.4) {}; \node[below] at (0.5, -1.4) {$a$};
        \node[mat, fill=mat, minimum size=3pt] at (1.5, -1.4) {}; \node[below] at (1.5, -1.4) {$b$};
        \node[left] at (0, -1.4) {4. };
        
        \node[above] at (1.5, 3.2) {\textbf{Chromatides filles (Gamètes)}};
        \node[rec, align=center] at (1.5, -2.2) {\textbf{2 chromatides parentales}\\(non recombinées)\\\textcolor{pat}{A-B} / \textcolor{mat}{a-b}\\\textbf{2 chromatides recombinées}\\\textcolor{rec}{A-b} / \textcolor{rec}{a-B}};
    \end{scope}
\end{tikzpicture}
\legende{Crossing-over (brassage intrachromosomique) entre deux gènes liés $A/a$ et $B/b$. L'échange de segments entre chromatides non-sœurs au niveau du chiasme produit 2 chromatides recombinées (nouvelles combinaisons $A-b$ et $a-B$) et 2 chromatides parentales ($A-B$, $a-b$).}
\end{popfigure}

\begin{attention}[Piège fréquent : Confondre création d'allèles et recombinaison]
\begin{itemize}
    \item \textbf{Mutation} = modification de la séquence d'ADN $\rightarrow$ \textbf{crée un allèle nouveau} (ex. $A \rightarrow a$). Rare, aléatoire, source ultime de la diversité.
    \item \textbf{Brassage (Crossing-over + Ségrégation + Fécondation)} = redistribution des allèles \emph{existants} $\rightarrow$ \textbf{crée de nouveaux génotypes (combinaisons)}. Fréquent, systématique à chaque méiose.
\end{itemize}
\textbf{Au Bac :} On ne dit jamais « le crossing-over crée de nouveaux allèles ». On dit : « Le crossing-over crée de nouvelles \emph{combinaisons d'allèles} sur un même chromosome. »
\end{attention}

\courssub{La fécondation aléatoire : l'ultime brassage}

La rencontre entre un spermatozoïde et un ovocyte est un événement aléatoire parmi des millions de possibilités.

\begin{propriete}[Nombre théorique de zygotes possibles pour un couple]
Si $N_p$ est le nombre de gamètes possibles pour le père et $N_m$ pour la mère, le nombre de zygotes distincts possibles est $N_p \times N_m$.
\par En ne considérant \textbf{que le brassage interchromosomique} :
$N_p = 2^{23}$ et $N_m = 2^{23}$.
\[ N_{zygotes} = 2^{23} \times 2^{23} = 2^{46} \approx 7,0 \times 10^{13} \text{ (70 000 milliards)}. \]
\par En ajoutant les \textbf{crossing-overs} (environ 1 à 3 par chromosome par méiose), la diversité réelle est \emph{astronomiquement plus grande} (virtuellement infinie à l'échelle humaine).
\end{propriete}

\begin{exemplebox}[Exemple chiffré : La « loterie » génétique]
Un couple camerounais a 4 enfants. La probabilité que deux enfants (non jumeaux monozygotes) soient \textbf{génétiquement identiques} pour l'ensemble de leur génome est :
\[ P = \frac{1}{2^{46}} \times \frac{1}{\text{facteur crossing-over}} \approx \mathbf{0}. \]
C'est pourquoi, hors vrais jumeaux (même zygote scindé), \textbf{chaque individu est génétiquement unique}. Cette unicité est la base de la biodiversité intraspécifique et de l'adaptation des populations.
\end{exemplebox}

\courssub{Lien avec les anomalies : révélation des allèles rares}

C'est ici que tout se relie aux sections précédentes (mutations et anomalies).

\begin{aretenir}[Synthèse : Mutation vs Brassage dans l'apparition des anomalies]
\begin{itemize}
    \item \textbf{Étape 1 (Temps long) :} Une \textbf{mutation} survient chez un ancêtre (ex. : apparition de l'allèle $S$ drépanocytaire ou $CFTR$ muté mucoviscidose). Cet allèle est rare.
    \item \textbf{Étape 2 (Temps moyen) :} L'allèle se transmet \textbf{à l'état hétérozygote} (porteur sain : $AS$, $Ff$) pendant des générations. Il est \textbf{masqué} phénotypiquement (récessif) ou peu exprimé. Le brassage le redistribue sans le révéler.
    \item \textbf{Étape 3 (Génération actuelle) :} Par hasard (brassage interchromosomique + fécondation), deux gamètes porteurs de l'allèle rare se rencontrent $\rightarrow$ Enfant \textbf{homozygote} ($SS$, $ff$). L'anomalie « apparaît » soudainement dans la famille.
\end{itemize}
\textbf{Conclusion scientifique :} L'anomalie n'est pas une « malédiction », un « sort » ou une « punition ». C'est la rencontre statistique prévisible (lois de Mendel) de deux allèles mutés hérités d'ancêtres communs ou distincts, rendue possible par le brassage méiotique.
\end{aretenir}

\begin{savaistu}[Le brassage et la résistance au paludisme au Cameroun]
L'allèle $S$ (drépanocytose) est fréquent en zone d'endémie palustre (forêt du Sud, bassin du Lac Tchad) car les hétérozygotes $AS$ résistent mieux au paludisme sévère (\emph{Plasmodium falciparum}). Le brassage maintient cet allèle dans la population : il le mélange constamment à l'allèle $A$ sain. Sans brassage (reproduction clonale), l'allèle $S$ aurait soit disparu, soit fixé, perdant l'avantage hétérozygote. La reproduction sexuée et son brassage sont donc un \textbf{moteur d'adaptation évolutive} majeur pour notre espèce face aux pathogènes locaux.
\end{savaistu}

\courssub{Éradication des préjugés : l'éducation scientifique comme devoir citoyen}

\begin{methode}[Comment expliquer scientifiquement l'apparition d'une anomalie en famille ?]
\begin{enumerate}
    \item \textbf{Identifier le mode de transmission} (arbre généalogique : récessif autosomique, dominant, lié au X...).
    \item \textbf{Rappeler l'origine moléculaire :} Mutation ponctuelle (drépanocytose, phénylcétonurie), délétion (mucoviscidose $\Delta F508$), non-disjonction (trisomie 21).
    \item \textbf{Expliquer le silence des générations précédentes :} Hétérozygotie des porteurs (récessifs) ou mutation \emph{de novo} (trisomie 21, achondroplasie, hémophilie nouveau mutant).
    \item \textbf{Insister sur le hasard méiotique :} « Ce n'est la faute de personne. C'est une combinaison aléatoire lors de la formation des gamètes, comme tirer une carte dans un jeu de 70 000 milliards. »
    \item \textbf{Valoriser le conseil génétique :} Au Cameroun, les centres de dépistage néonatal (ex. : Centre Mère-Enfant de la Fondation Chantal Biya, Hôpital Général de Yaoundé) et les associations (ASDREP pour la drépanocytose) offrent information et prise en charge.
\end{enumerate}
\end{methode}

\begin{attention}[Préjugés à déconstruire absolument]
\begin{itemize}
    \item \textbf{Faux :} « C'est la femme qui a « apporté » la maladie. » $\rightarrow$ \textbf{Vrai :} Les deux parents transmettent un allèle (sauf anomalies liées au X ou \emph{de novo} paternelles).
    \item \textbf{Faux :} « C'est un sort jeté par la belle-famille / un esprit. » $\rightarrow$ \textbf{Vrai :} C'est un événement biophysique (mutation, non-disjonction, crossing-over) soumis aux lois de la chimie et des probabilités.
    \item \textbf{Faux :} « Si on a un enfant malade, les suivants le seront tous. » $\rightarrow$ \textbf{Vrai :} Pour une maladie récessive, le risque est de 1/4 \emph{à chaque grossesse}, indépendant des précédentes (indépendance des événements).
    \item \textbf{Faux :} « Les jumeaux sont identiques donc s'ils sont malades, c'est héréditaire. » $\rightarrow$ \textbf{Vrai :} Jumeaux monozygotes = même génome (sauf mutation somatique). Jumeaux dizygotes = frères ordinaires (50\% gènes communs).
\end{itemize}
\end{attention}

\begin{exoresolu}[Communication médicale : Annonce d'un risque de drépanocytose]
\textbf{Énoncé.} Un couple AS $\times$ AS consulte après la naissance d'un enfant SS. Le père accuse la mère : « C'est ton sang qui est mauvais ». En tant que conseiller en génétique, rédigez une réponse scientifique, empathique et culturellement adaptée.
\tcblower
\textbf{Solution type.}
« Monsieur, Madame, je comprends votre douleur. Scientifiquement, la drépanocytose est une maladie génétique \textbf{récessive autosomique}. Vous êtes \textbf{tous les deux} porteurs sains (génotype AS) : vous avez chacun un allèle normal $A$ et un allèle muté $S$ sur le chromosome 11. Cet allèle $S$ est très fréquent au Cameroun car il protège contre le paludisme grave quand on ne l'a qu'en un seul exemplaire (AS).
\par Lors de la formation de vos gamètes (méiose), le \textbf{hasard} (brassage) a fait que votre spermatozoïde et l'ovocyte de votre épouse portaient \textbf{tous les deux} l'allèle $S$. La probabilité était de 1 sur 4 (25 \%). Ce n'est la faute de personne, ni du père, ni de la mère, ni des ancêtres. C'est une loi de la biologie.
\par Pour les prochaines grossesses, le risque reste 1/4 à chaque fois, mais il y a 3 chances sur 4 d'avoir un enfant non drépanocytaire (AA ou AS). Un diagnostic prénatal est possible. L'enfant SS peut aujourd'hui bien vivre avec un suivi adapté (folates, vaccins, hydroxyurée, hydratation). L'association ASDREP vous accompagnera. »
\end{exoresolu}

\begin{aretenir}[Message final pour le Bac et pour la vie]
\begin{itemize}
    \item \textbf{Mutation} = Source des allèles nouveaux (diversité \emph{de novo}).
    \item \textbf{Brassage (Méiose + Fécondation)} = Moteur de la diversité \emph{combinatoire} (unicité des individus).
    \item \textbf{Anomalies familiales} = Résultat de la rencontre (hasard du brassage) d'allèles mutés préexistants ou \emph{de novo}.
    \item \textbf{Science} = Outil de compréhension, de prévention (conseil génétique, dépistage néonatal) et de lutte contre la stigmatisation.
\end{itemize}
\end{aretenir}

\coursec{VI. Éradication des préjugés : éducation et sensibilisation}

Après avoir exploré les mécanismes moléculaires et chromosomiques à l'origine de la diversité génétique (mutations, brassage allélique, non-disjonctions), nous comprenons désormais que l'apparition d'un caractère nouveau ou d'une anomalie au sein d'une famille n'est pas un événement mystérieux, mais la conséquence naturelle et prévisible de processus biologiques identifiés. Cette section vise à transformer cette rigueur scientifique en levier citoyen : comment, au Cameroun comme ailleurs, faire reculer les préjugés qui pèsent sur les familles touchées par la drépanocytose, l'hémophilie, la trisomie 21 ou le syndrome de Turner ? Comment remplacer la stigmatisation par la solidarité éclairée ?

\courssub{Recensement des préjugés courants au Cameroun : entre tradition et méconnaissance}

Dans de nombreuses communautés camerounaises — qu'elles soient urbaines (Yaoundé, Douala, Bafoussam) ou rurales —, l'annonce d'une anomalie génétique chez un nouveau-né ou un enfant déclenche souvent des interprétations culturelles qui, bien qu'ancrées dans l'histoire sociale, n'ont aucun fondement biologique. Enquête après enquête, les équipes de génétique clinique du Centre Hospitalier Universitaire de Yaoundé et des centres de santé de district rapportent les mêmes récits :

\begin{itemize}
    \item \textbf{La sorcellerie ou la « manguier » :} l'enfant serait « mangé » par un esprit maléfique envoyé par un rival, un oncle maternel jaloux ou une co-épouse.
    \item \textbf{La malédiction ancestrale :} la famille aurait commis un sacrilège (non-respect d'un totem, violation d'un interdit funéraire) qui « retombe » sur la descendance.
    \item \textbf{L'infidélité de la mère :} l'anomalie serait la preuve visible d'un adultère, le père biologique portant le « mauvais gène ».
    \item \textbf{Le « mauvais sang » de la belle-famille :} la lignée maternelle (ou paternelle) serait « porteuse de malheur », justifiant parfois le rejet de l'épouse ou la rupture du mariage.
\end{itemize}

Ces croyances, bien que compréhensibles face à l'angoisse de l'inconnu, ont des conséquences dramatiques : abandon de l'enfant, divorce, exclusion sociale de la mère, refus de soins médicaux au profit de traitements traditionnels inefficaces, voire infanticides déguisés. Elles transforment une épreuve médicale en tragédie familiale et sociale.

\courssub{Réfutation scientifique point par point : la biologie comme antidote}

Chaque préjugé s'effondre dès qu'on lui oppose les faits expérimentaux établis par la génétique moderne, reproductibles dans n'importe quel laboratoire de cytogénétique ou de biologie moléculaire.

\begin{popfigure}
\begin{tikzpicture}[
    node distance=2mm and 4mm,
    prejug/.style={rectangle, draw=poporange, thick, fill=poporangeL, rounded corners, minimum height=1.2cm, text width=6.5cm, align=left, font=\small},
    realite/.style={rectangle, draw=popgreen, thick, fill=popgreenL, rounded corners, minimum height=1.2cm, text width=6.5cm, align=left, font=\small},
    arrow/.style={->, >=Stealth, thick, popdark}
]
% En-têtes
\node[prejug, fill=poporange, font=\bfseries\textcolor{white}, minimum height=0.8cm] (h1) at (0,0) {Préjugé courant};
\node[realite, fill=popgreen, font=\bfseries\textcolor{white}, minimum height=0.8cm] (h2) at (7.5,0) {Réalité scientifique démontrée};

% Ligne 1
\node[prejug, below=of h1] (p1) {Sorcellerie / « manguier » : cause surnaturelle, intentionnelle.};
\node[realite, below=of h2] (r1) {Mutation ponctuelle (ex. \textit{HBB}: Glu6Val) ou non-disjonction méiotique : événement aléatoire, sans intentionnalité.};

% Ligne 2
\node[prejug, below=of p1] (p2) {Malédiction ancestrale : punition pour faute des ancêtres.};
\node[realite, below=of r1] (r2) {Allèles mutés présents dans toute population (fréquence allélique mesurable). Aucune corrélation avec l'histoire familiale « morale ».};

% Ligne 3
\node[prejug, below=of p2] (p3) {Infidélité de la mère : l'anomalie prouve l'adultère.};
\node[realite, below=of r2] (r3) {Transmission autosomique récessive (drépanocytose) ou liée à l'X (hémophilie) : les deux parents sains peuvent être porteurs. Caryotype / séquençage confirment l'origine parentale.};

% Ligne 4
\node[prejug, below=of p3] (p4) {« Mauvais sang » de la belle-famille : lignée maudite.};
\node[realite, below=of r3] (r4) {Chaque individu porte en moyenne 1 à 2 mutations récessives létales ou délétères à l'état hétérozygote. Aucun « sang pur » n'existe.};

% Flèches
\draw[arrow] (p1.east) -- (r1.west) node[midway, above, font=\tiny, popdark] {Électrophorèse / Séquençage};
\draw[arrow] (p2.east) -- (r2.west) node[midway, above, font=\tiny, popdark] {Études de fréquence allélique};
\draw[arrow] (p3.east) -- (r3.west) node[midway, above, font=\tiny, popdark] {Arbre généalogique / Test ADN};
\draw[arrow] (p4.east) -- (r4.west) node[midway, above, font=\tiny, popdark] {Génomique des populations};
\end{tikzpicture}
\legende{Tableau comparatif : préjugés courants au Cameroun et réalités scientifiques établies par la biologie moléculaire et la génétique des populations.}
\end{popfigure}

\begin{experience}[Démonstration par l'électrophorèse d'hémoglobine — TP réalisable en classe]
\textbf{Objectif :} Montrer qu'un couple de parents phénotypiquement sains (AS $\times$ AS) peut avoir un enfant drépanocytaire (SS) sans aucune infidélité ni sorcellerie.

\textbf{Matériel :} Échantillons d'hémoglobine (contrôles AA, AS, SS), gel d'agarose, chambre d'électrophorèse, tampon Tris-EDTA-Borate pH 8,6, colorant (bleu de Coomassie ou rouge de Ponceau).

\textbf{Protocole :}
\begin{enumerate}
    \item Déposer \SI{10}{\micro\liter} de chaque échantillon dans les puits du gel.
    \item Migrer à \SI{100}{\volt} pendant \SI{30}{\minute} (l'hémoglobine S migre plus lentement que l'hémoglobine A vers l'anode).
    \item Colorer, déstacher, photographier.
\end{enumerate}

\textbf{Observations attendues :}
\begin{itemize}
    \item Père (AS) : deux bandes (A et S).
    \item Mère (AS) : deux bandes (A et S).
    \item Enfant 1 (AA) : une bande A.
    \item Enfant 2 (AS) : deux bandes A et S.
    \item Enfant 3 (SS) : une bande S.
\end{itemize}

\textbf{Interprétation :} Le patron de bandes suit exactement les lois de Mendel pour un caractère autosomique récessif. L'enfant SS a reçu l'allèle S de son père \textbf{et} l'allèle S de sa mère. L'électrophorèse est une preuve physique, visible, incontestable de la transmission mendélienne. Elle exclut toute cause surnaturelle ou infidélité.
\end{experience}

\begin{propriete}[Porteurs sains : une universalité biologique]
Tout être humain, quel que soit son origine ethnique, son rang social ou son histoire familiale, est porteur en moyenne de \textbf{1 à 2 allèles mutés récessifs} létaux ou délétères à l'état hétérozygote. Ces allèles sont silencieux phénotypiquement (l'allèle normal assure la fonction) mais peuvent se retrouver homozygotes chez un enfant si les deux parents sont porteurs du même allèle muté. L'apparition d'une anomalie génétique chez un enfant est donc un \textbf{événement probabiliste} (risque de 25 \% pour chaque grossesse si les deux parents sont hétérozygotes), \textbf{pas une punition}, pas une malédiction, pas la preuve d'une faute.
\end{propriete}

\courssub{Les parents ne sont ni coupables ni maudits : la génétique des populations l'atteste}

Prenons l'exemple de la drépanocytose au Cameroun. La fréquence de l'allèle $S$ (\textit{HBB}$^S$) varie selon les régions : environ 10 à 15 \% en zone de forêt (Sud, Est, Centre) et jusqu'à 20 à 25 \% dans certaines zones de savane (Nord, Adamaoua) où le paludisme a maintenu l'allèle par sélection équilibrée (avantage hétérozygote AS face à \emph{Plasmodium falciparum}).

\begin{exoresolu}[Calcul du risque pour un couple AS $\times$ AS au Cameroun]
\textbf{Énoncé :} Dans une population où la fréquence de l'allèle $S$ est $q = 0,15$, on suppose l'équilibre de Hardy-Weinberg. Quelle est la probabilité qu'un couple formé au hasard soit AS $\times$ AS ? Quel est alors le risque d'avoir un enfant SS par grossesse ?

\textbf{Solution détaillée :}
\begin{enumerate}
    \item Fréquence des génotypes (Hardy-Weinberg) :
    \begin{align*}
        f(AA) &= p^2 = (1-q)^2 = 0,85^2 = 0,7225 \\
        f(AS) &= 2pq = 2 \times 0,85 \times 0,15 = 0,255 \\
        f(SS) &= q^2 = 0,15^2 = 0,0225
    \end{align*}
    \item Probabilité qu'un couple soit AS $\times$ AS (choix indépendant) :
    \[
    P(\text{AS} \times \text{AS}) = f(AS)^2 = 0,255^2 \approx 0,065 = 6,5 \%
    \]
    \item Risque d'enfant SS pour un tel couple : $1/4 = 25 \%$ par grossesse.
    \item Conclusion : Environ 1 couple sur 15 au Cameroun est concerné par ce risque, \textbf{sans aucune particularité familiale}. Ce n'est pas une « malédiction » qui vise une famille, c'est une statistique de population.
\end{enumerate}
\end{exoresolu}

\begin{attention}[Piège à éviter : confondre « cause génétique » et « responsabilité parentale »]
Dire « c'est génétique » ne signifie \textbf{pas} « c'est la faute des parents ». Les parents n'ont pas \emph{choisi} leurs allèles ; ils les ont reçus au hasard de la méiose de leurs propres parents. La méiose est un tirage au sort moléculaire (loi de la ségrégation, loi de l'assortiment indépendant, crossing-over). Accuser les parents, c'est accuser la biologie fondamentale de la reproduction. La culpabilité n'a pas de place dans l'explication scientifique.
\end{attention}

\courssub{Rôle de la famille et de la communauté : du rejet au soutien}

La science nous donne les faits ; l'éthique et la culture nous disent comment les vivre. Au Cameroun, la famille élargie (clan, lignage) est traditionnellement le filet de sécurité. Face à une anomalie génétique, ce filet doit se transformer en \textbf{réseau de soin et d'inclusion} :

\begin{itemize}
    \item \textbf{Soutien affectif :} l'enfant atteint (drépanocytaire, hémophile, trisomique 21) a les mêmes besoins d'amour, d'éducation, de jeu que tout enfant. La stigmatisation (le cacher, le priver d'école, le considérer comme « déjà mort ») aggrave son pronostic vital et psychomoteur.
    \item \textbf{Soutien médical :} la drépanocytose se soigne (hydroxyurée, vaccination anti-pneumocoque, folates, transfusion si anémie sévère, greffe de moelle osseuse dans des centres spécialisés comme l'Hôpital Général de Yaoundé). L'hémophilie se traite par substitution du facteur manquant (disponible via la Fédération Mondiale de l'Hémophilie). La trisomie 21 bénéficie d'une prise en charge multidisciplinaire (cardiologie, orthophonie, kinésithérapie, scolarisation adaptée).
    \item \textbf{Soutien social :} les associations de patients (ex. \emph{Association Camerounaise de Lutte contre la Drépanocytose}, \emph{Association des Parents d'Enfants Trisomiques 21}) offrent un espace d'entraide, d'information et de plaidoyer. La communauté religieuse, les chefs traditionnels, les enseignants peuvent relayer des messages de tolérance.
\end{itemize}

\begin{savaistu}[Le « mois de la drépanocytose » au Cameroun]
Chaque année, le 19 juin (Journée mondiale de la drépanocytose) et tout le mois de juin, le Ministère de la Santé Publique (MINSANTE), avec l'appui de l'OMS et d'ONG, organise des campagnes de dépistage gratuit, de sensibilisation dans les marchés, les écoles, les médias (CRTV, radios communautaires). Le slogan « Connaître son statut, c'est protéger ses enfants » replace la responsabilité au niveau de l'information, non de la culpabilité. Depuis 2019, le dépistage néonatal systématique (test de Guthrie sur papier buvard) est en phase de généralisation dans les maternités de référence : il permet une prise en charge précoce (pénicilline prophylactique, éducation des parents) qui divise par 10 la mortalité avant 5 ans.
\end{savaistu}

\courssub{Outils de prévention et d'accompagnement : la génétique clinique au service des familles}

La génétique n'est pas seulement explicative, elle est \textbf{prédictive et préventive}. Trois piliers structurent l'accompagnement des couples à risque :

\begin{popfigure}
\begin{tikzpicture}[
    node distance=1.5cm and 2cm,
    stepbox/.style={rectangle, draw=popblue, thick, fill=popblueL, rounded corners, minimum width=3cm, minimum height=1.2cm, align=center, font=\small},
    arrow/.style={->, >=Stealth, thick, popdark},
    decision/.style={diamond, draw=poporange, thick, fill=poporangeL, aspect=2, minimum width=3.5cm, align=center, font=\small}
]
\node[stepbox, align=center] (depistage){Dépistage\\\textbf{Prénuptial / Prénatal / Néonatal}};
\node[decision, right=of depistage, align=center] (conseil){Conseil génétique\\\textbf{(Information non directive)}};
\node[stepbox, right=of conseil, align=center] (choix){Choix éclairés\\\textbf{Projet parental adapté}};

\draw[arrow] (depistage) -- node[above, font=\tiny] {Si risque identifié} (conseil);
\draw[arrow] (conseil) -- node[above, font=\tiny] {Autonomie du couple} (choix);

% Détails sous chaque étape
\node[below=0.8cm of depistage, font=\scriptsize, text width=3.5cm, align=center, text=popmuted] (d1) {Électrophorèse d'hémoglobine (drépanocytose)\newline Caryotype (trisomie 21)\newline Séquençage ciblé (mucoviscidose, myopathie)};
\node[below=0.8cm of conseil, font=\scriptsize, text width=3.5cm, align=center, text=popmuted] (c1) {Explication du mode de transmission\newline Risques de récurrence (25\%, 50\%...)\newline Options : DPN, DPI, adoption, acceptation};
\node[below=0.8cm of choix, font=\scriptsize, text width=3.5cm, align=center, text=popmuted] (ch1) {Grossesse surveillée + DPN (amniocentèse)\newline FIV + DPI (diagnostic préimplantatoire)\newline Suivi médical spécialisé dès la naissance};
\end{tikzpicture}
\legende{Parcours de prévention génétique : du dépistage au choix éclairé, en passant par le conseil génétique non directif. DPN = Diagnostic Prénatal ; DPI = Diagnostic Préimplantatoire ; FIV = Fécondation In Vitro.}
\end{popfigure}

\begin{methode}[Conseil génétique : les 5 étapes clés]
\begin{enumerate}
    \item \textbf{Recueil de l'histoire familiale :} arbre généalogique sur 3 générations, identification des cas atteints, consanguinité.
    \item \textbf{Confirmation diagnostique :} caryotype (anomalies chromosomiques), électrophorèse / séquençage (anomalies géniques).
    \item \textbf{Calcul du risque de récurrence :} selon le mode de transmission (autosomique récessif/dominant, lié à l'X, mitochondrial, chromosomique \textit{de novo} ou héréditaire).
    \item \textbf{Information sur les options :} diagnostic prénatal (amniocentèse 15-18 SA, trophoblastie 11-13 SA), diagnostic préimplantatoire (FIV + biopsie blastomère), adoption, acceptation du risque avec préparation médicale.
    \item \textbf{Soutien psychologique et suivi :} le conseil est \textbf{non directif} : le généticien informe, le couple décide selon ses valeurs. Aucune pression, aucun jugement.
\end{enumerate}
\end{methode}

\courssub{Message de citoyenneté : la différence génétique, une donnée naturelle ; la solidarité, un devoir social}

La variabilité génétique est le moteur de l'évolution et la condition de la survie de notre espèce face aux pathogènes, aux changements climatiques, aux nouvelles maladies. \textbf{Chaque allèle « anormal » aujourd'hui a pu être un avantage sélectif hier} (ex. allèle $S$ et résistance au paludisme ; allèle $\Delta$F508 de la mucoviscidose : hypothèse d'avantage sélectif hétérozygote face au choléra ou à la typhoïde). Il n'y a pas de « génome parfait », seulement des génomes adaptés à un environnement donné.

\begin{aretenir}[Citoyenneté génétique]
\begin{itemize}
    \item La différence génétique (mutation, anomalie chromosomique) est un \textbf{fait naturel}, universel, inévitable.
    \item La stigmatisation des personnes atteintes et de leurs familles est une \textbf{construction sociale}, injuste, évitable.
    \item L'accès au dépistage, au conseil génétique, aux soins et à l'éducation inclusive est un \textbf{droit fondamental} (Convention relative aux droits des personnes handicapées, ratifiée par le Cameroun en 2021).
    \item Chaque citoyen, chaque professionnel de santé, chaque enseignant, chaque chef traditionnel, chaque journaliste a le devoir de \textbf{relayer l'explication scientifique} pour remplacer la peur par la connaissance, le rejet par l'entraide.
\end{itemize}
\end{aretenir}

\courssub{Savoir-faire : rédiger un message de sensibilisation argumenté}

Le programme exige que vous sachiez construire un argumentaire de sensibilisation à destination d'un public non scientifique (famille, communauté, association de parents d'élèves, groupe religieux). La démarche suit une logique d'investigation inversée : \textbf{partir du fait biologique observable pour déconstruire le préjugé}.

\begin{methode}[Construire un argumentaire de sensibilisation — 4 étapes]
\begin{enumerate}
    \item \textbf{Choisir un fait observable incontestable :} caryotype (47,XX,+21), électrophorèse (bande S), arbre généalogique (transmission récessive), séquençage (mutation ponctuelle).
    \item \textbf{En déduire le mécanisme biologique :} non-disjonction méiotique (hasard), mutation ponctuelle (erreur de réplication), transmission mendélienne (lois statistiques). \emph{Insister sur l'absence d'intentionnalité.}
    \item \textbf{Montrer que les parents ne sont pas responsables :} ils n'ont pas choisi leurs gamètes ; la méiose est un tirage au sort. Tout humain est porteur sain de mutations récessives.
    \item \textbf{Conclure sur l'action solidaire :} dépistage, conseil, soins, inclusion scolaire, lutte contre la stigmatisation. \emph{Le message doit être respectueux des croyances, mais ferme sur les faits.}
\end{enumerate}
\end{methode}

\begin{exemplebox}[Message de sensibilisation type — Drépanocytose, pour une réunion de parents d'élèves]
\textbf{Titre :} « La drépanocytose : comprendre pour mieux accompagner nos enfants »

\textbf{Accroche (fait observable) :} « Chers parents, regardez cette électrophorèse d'hémoglobine. Ce gel montre trois bandes : A (normale), S (drépanocytaire). Un enfant SS n'a que la bande S. Ses parents ont tous les deux les bandes A et S. Ils sont en parfaite santé. »

\textbf{Explication (mécanisme) :} « L'allèle S est une mutation ancienne, présente chez 15 \% de nous au Cameroun. Elle protège contre le paludisme grave quand on en a une seule copie (AS). Mais si un enfant reçoit l'allèle S de son père \textbf{et} de sa mère (risque 1 sur 4), il fait une drépanocytose. Ce n'est pas une sorcellerie, c'est la loi du hasard à la méiose. »

\textbf{Déconstruction du préjugé :} « Accuser la mère d'infidélité ou la belle-famille de "mauvais sang" est faux biologiquement : le père a donné son allèle S, la mère a donné le sien. Tous deux sont des porteurs sains, comme des millions de Camerounais. »

\textbf{Appel à l'action :} « Faisons dépister nos enfants à la naissance (test gratuit). Soutenons les familles touchées : hydroxyurée, vaccins, eau potable, école inclusive. Refusons la stigmatisation. La science nous éclaire ; notre humanité nous guide. »
\end{exemplebox}

\begin{exoresolu}[Analyse critique d'un discours stigmatisant]
\textbf{Énoncé :} Lors d'une réunion de famille, un oncle déclare : « Cet enfant trisomique, c'est la faute de sa mère : elle a trop travaillé au champ pendant sa grossesse, et puis sa famille a toujours eu des "enfants bizarres". »

\textbf{Analyse et réponse argumentée :}
\begin{enumerate}
    \item \textbf{Identification des affirmations non scientifiques :} (1) Travail maternel cause de trisomie 21. (2) Hérédité familiale systématique (« toujours eu »).
    \item \textbf{Réfutation par les faits :}
    \begin{itemize}
        \item Trisomie 21 = 47 chromosomes (trois chromosomes 21) dans \textbf{95 \% des cas par non-disjonction méiotique maternelle \textit{de novo}} (accident aléatoire à la méiose I ou II de l'ovocyte). Aucune corrélation avec l'activité physique, l'alimentation, le stress.
        \item L'âge maternel est le seul facteur de risque établi (risque $\approx$ 1/1500 à 20 ans, 1/350 à 35 ans, 1/30 à 45 ans). Pas de « famille maudite » : la récurrence est exceptionnelle (< 1 \%) sauf translocation robertsonienne héréditaire (rare, détectable par caryotype parental).
    \end{itemize}
    \item \textbf{Message de réorientation :} « Oncle, le caryotype de l'enfant montre 47,XX,+21. C'est un accident de la méiose, comme une erreur de copier-coller dans une cellule. La mère n'y est pour rien. La famille n'y est pour rien. Cet enfant a besoin de notre amour, de suivi cardiaque, de kiné, d'école. Aidons-le. »
\end{enumerate}
\end{exoresolu}

\begin{attention}[Pédagogie du respect : ne pas braquer, mais éclairer]
Combattre un préjugé ne signifie pas nier les croyances culturelles ou religieuses des personnes. Cela signifie : \textbf{« Je comprends que cette explication ait du sens pour vous. Permettez-moi de vous montrer ce que la biologie observe au microscope et au séquenceur. Les deux approches peuvent coexister, mais pour prendre des décisions de santé, les faits biologiques sont la base la plus sûre. »} L'arrogance scientifique referme les portes ; l'humilité pédagogique les ouvre.
\end{attention}

\begin{savaistu}[Phénylcétonurie : la victoire du dépistage néonatal]
La phénylcétonurie (PKU, mutation du gène \textit{PAH} chromosome 12, autosomique récessive) entraîne une accumulation de phénylalanine neurotoxique. Sans traitement : retard mental sévère, épilepsie, troubles du comportement. \textbf{Depuis les années 1970, le test de Guthrie (quelques gouttes de sang sur papier buvard au 3\textsuperscript{e} jour de vie) permet un diagnostic précoce.} Un régime sans phénylalanine (lait maternel dosé, aliments hypoprotidiques spéciaux) débuté avant 3 semaines de vie permet un développement \textbf{strictement normal} (QI normal, scolarité normale, vie professionnelle normale). Au Cameroun, le test n'est pas encore systématique partout, mais des projets pilotes (Hôpital Central Yaoundé, Laquintinie Douala) prouvent sa faisabilité. C'est la preuve ultime : \textbf{la génétique n'est pas une fatalité, la connaissance est un pouvoir d'agir.}
\end{savaistu}

\courssub{Exercice d'application : vous êtes l'ambassadeur de la génétique}

\begin{exercice}[Rédaction d'une affiche de sensibilisation]
Vous êtes chargé(e) de concevoir une affiche A4 pour le hall d'une maternité de district au Cameroun, destinée aux futurs parents. L'affiche doit :
\begin{enumerate}
    \item Présenter \textbf{un fait biologique simple} (ex. « Nous sommes tous porteurs sains de mutations »).
    \item Déconstruire \textbf{un préjugé précis} (ex. « Non, la drépanocytose n'est pas une sorcellerie »).
    \item Donner \textbf{une action concrète} (ex. « Demandez le test d'électrophorèse avant le mariage — gratuit au Centre de Santé »).
    \item Utiliser un \textbf{visuel clair} (schéma de transmission AS $\times$ AS, ou caryotype 47,XX,+21) et un \textbf{slogan positif}.
\end{enumerate}
\textbf{Contraintes :} Français correct, ton respectueux, pas de jargon non expliqué, couleurs vives, numéros de téléphone utiles (Centre de Génétique Médicale, Association de patients).
\end{exercice}

\begin{corrige}[Exercice n°1 — Éléments attendus pour l'affiche]
\begin{itemize}
    \item \textbf{Titre :} « La drépanocytose : une histoire de gènes, pas de sortilèges. »
    \item \textbf{Visuel central :} Carré de Punnett AS $\times$ AS avec pourcentages (25 \% AA, 50 \% AS, 25 \% SS) + photo d'électrophorèse réelle (bandes A, S).
    \item \textbf{Message clé :} « Porter l'allèle S ne rend pas malade. Le transmettre à deux, c'est le hasard. Savoir son statut, c'est protéger son enfant. »
    \item \textbf{Action :} « Test gratuit : Centre de Santé de District (lundi-vendredi 8h-15h). Conseil génétique : Hôpital Central Yaoundé, Laquintinie Douala. »
    \item \textbf{Slogan bas de page :} « Connaître, c'est pouvoir aimer mieux. Ensemble contre la stigmatisation. »
\end{itemize}
\end{corrige}

\courssub{Bilan de la leçon : de la molécule à la citoyenneté}

Cette leçon nous a fait voyager de l'échelle nanométrique (mutation ponctuelle Glu6Val sur le gène \textit{HBB}) à l'échelle de la société camerounaise (préjugés, politiques de santé publique). Nous avons vu que :
\begin{itemize}
    \item Les \textbf{mutations géniques} (substitution, insertion, délétion) et \textbf{chromosomiques} (non-disjonction, translocation) sont les sources ultimes de tout allèle nouveau.
    \item Le \textbf{brassage allélique} (méiose, fécondation) rend chaque individu génétiquement unique et redistribue les allèles mutés à chaque génération.
    \item Les \textbf{anomalies géniques} (drépanocytose, hémophilie, phénylcétonurie, mucoviscidose, myopathie) et \textbf{chromosomiques} (trisomie 21, Turner, Klinefelter) ont des causes identifiées, des mécanismes compris, des diagnostics fiables.
    \item La \textbf{science génétique} n'est pas une connaissance abstraite : elle est un outil de \textbf{prévention} (dépistage), de \textbf{prise en charge} (thérapies), et de \textbf{libération} (lutte contre les préjugés).
\end{itemize}

En tant que futurs citoyens, professionnels de santé, enseignants, parents, vous portez la responsabilité de \textbf{dire le vrai avec humanité}. Chaque fois que vous entendez « c'est la sorcellerie », « c'est la faute de la femme », « cette famille est maudite », vous avez désormais les arguments — le caryotype, l'électrophorèse, l'arbre généalogique, la fréquence allélique — pour répondre : « Regardons les faits ensemble. Cet enfant a besoin de nous, pas de nos accusations. »

\begin{aretenir}[Compétences clés validées — Module I, Leçon 5]
\begin{itemize}
    \item \textbf{Savoir :} Citer les types de mutations, expliquer le brassage allélique, décrire 2 anomalies géniques et 3 anomalies chromosomiques.
    \item \textbf{Savoir-faire :} Identifier une mutation sur séquence d'ADN, identifier une anomalie chromosomique sur caryotype, calculer un risque de récurrence, construire un argumentaire de sensibilisation.
    \item \textbf{Savoir-être :} Rejeter la stigmatisation, promouvoir le dépistage et le conseil génétique, faire preuve d'empathie envers les familles concernées, agir en citoyen responsable face à la diversité génétique.
\end{itemize}
\end{aretenir}

\newpage

\coursec{Activité d'intégration}

\courssub{Objectifs de l'activité}

À l'issue de cette activité d'intégration, tu seras capable de :
\begin{itemize}
\item identifier une \cle{mutation génique} (substitution) par comparaison de séquences d'ADN et relier cette mutation à l'anomalie de l'hémoglobine ;
\item déterminer les \cle{génotypes} d'individus à partir d'un arbre généalogique et de résultats d'électrophorèse ;
\item calculer le \cle{risque génétique} d'une descendance à l'aide d'un échiquier de croisement ;
\item identifier une \cle{anomalie chromosomique} (trisomie 21) sur un caryotype et expliquer son mécanisme (non-disjonction méiotique) ;
\item construire une argumentation scientifique destinée à \cle{éradiquer les préjugés} familiaux et sociaux autour des anomalies génétiques.
\end{itemize}

\courssub{Situation}

Dans le quartier Nkol-Eton à Yaoundé, la famille Ngo Bell est en crise. M. Ngo Bell et Mme Ngo Bell, tous deux en bonne santé apparente, ont mis au monde leur troisième enfant, le petit Junior. Dès l'âge de 8 mois, Junior présente des crises douloureuses répétées (mains et pieds gonflés, fièvres, pâleur). Le diagnostic du Centre Hospitalier d'Essos est sans appel : \cle{drépanocytose} (anémie falciforme). La belle-famille accuse alors Mme Ngo Bell de sorcellerie et exige son renvoi du domicile conjugal. Par ailleurs, un cousin de la famille, âgé de 5 ans, présente un \cle{syndrome de Down} ; sa mère avait 43 ans à sa naissance, et certains membres de la famille parlent de « malédiction ».

Tu es membre du club de sciences du lycée, invité par le médecin du district à participer à une séance de \cle{conseil génétique communautaire}. Tu disposes du dossier suivant.

\textbf{Document 1 : arbre généalogique de la famille Ngo Bell.}

\begin{popfigure}
\begin{center}
\begin{tikzpicture}[scale=1.05, every node/.style={font=\small}]
% Génération I : parents
\node[draw, rectangle, minimum size=0.55cm, label=above:{I-1 M. Ngo Bell (sain)}] (pere) at (0,0) {};
\node[draw, circle, minimum size=0.55cm, label=above:{I-2 Mme Ngo Bell (saine)}] (mere) at (2.2,0) {};
\draw (pere.east) -- (mere.west);
\draw (1.1,0) -- (1.1,-0.7);
\draw (1.1,-0.7) -- (-1.3,-1.2);
\draw (1.1,-0.7) -- (1.1,-1.2);
\draw (1.1,-0.7) -- (3.5,-1.2);
% Génération II
\node[draw, circle, minimum size=0.55cm, label=below:{II-1 fille saine}] at (-1.3,-1.6) {};
\node[draw, rectangle, minimum size=0.55cm, label=below:{II-2 garçon sain}] at (1.1,-1.6) {};
\node[draw, rectangle, fill=popink, minimum size=0.55cm, label=below:{II-3 Junior (drépanocytaire)}] at (3.5,-1.6) {};
\end{tikzpicture}
\end{center}
\legende{Arbre généalogique de la famille Ngo Bell. Carré = homme ; cercle = femme ; symbole rempli = individu atteint de drépanocytose ; trait horizontal = union ; traits verticaux = descendance. Les parents I-1 et I-2 sont phénotypiquement sains.}
\end{popfigure}

\textbf{Document 2 : électrophorèse de l'hémoglobine (pH alcalin).}

\begin{center}
\begin{tabularx}{\linewidth}{|l|X|X|X|}
\hline
\rowcolor{popblueL} \textbf{Individu} & \textbf{Bande HbA} (migration rapide) & \textbf{Bande HbS} (migration lente) & \textbf{Interprétation} \\
\hline
I-1 (père) & présente & présente & à déterminer \\
\hline
I-2 (mère) & présente & présente & à déterminer \\
\hline
II-3 (Junior) & absente & présente & à déterminer \\
\hline
\end{tabularx}
\end{center}

\textbf{Document 3 : séquences partielles du brin codant du gène de la $\beta$-globine (région du codon 6).}

\begin{center}
\begin{tabularx}{\linewidth}{|l|X|}
\hline
\rowcolor{popblueL} \textbf{Individu} & \textbf{Séquence partielle (5' $\rightarrow$ 3')} \\
\hline
Allèle normal (référence) & ... CCT \textbf{GAG} GAG ... (codon 6 = Glu) \\
\hline
I-1 et I-2 (parents) & ... CCT \textbf{GAG} GAG ... \emph{et} ... CCT \textbf{GTG} GAG ... \\
\hline
II-3 (Junior) & ... CCT \textbf{GTG} GAG ... (sur les deux chromosomes) \\
\hline
\end{tabularx}
\end{center}

\textbf{Document 4 : caryotype du cousin.} Le caryotype montre 47 chromosomes : 22 paires d'autosomes normales, \textbf{trois exemplaires du chromosome 21}, et une paire de chromosomes sexuels XY. Formule chromosomique : 47, XY, +21.

\textbf{Document 5 : fréquence de la trisomie 21 à la naissance selon l'âge maternel.}

\begin{popfigure}
\begin{center}
\begin{tikzpicture}
\begin{axis}[
width=0.85\linewidth, height=5.5cm,
xlabel={Âge de la mère (ans)},
ylabel={Fréquence (pour 1000 naissances)},
xmin=15, xmax=50, ymin=0, ymax=35,
xtick={20,25,30,35,40,45},
grid=both, grid style={popline!40},
legend style={at={(0.02,0.98)}, anchor=north west, font=\small}
]
\addplot[mark=*, only marks, poporange, mark size=2pt] coordinates {(20,0.7) (25,0.9) (30,1.2) (35,2.8) (40,9.0) (45,25.0)};
\addlegendentry{Fréquences observées}
\addplot[popcurve, dashed, domain=18:47, samples=100] {0.7*exp(0.145*(x-20))};
\addlegendentry{Tendance (croissance exponentielle)}
\end{axis}
\end{tikzpicture}
\end{center}
\legende{Fréquence de la trisomie 21 à la naissance en fonction de l'âge maternel (ordres de grandeur classiques : environ 1/1500 à 20 ans, 1/350 à 35 ans, 1/100 à 40 ans, 1/30 à 45 ans). La fréquence reste faible et stable avant 30 ans, puis augmente fortement après 35 ans.}
\end{popfigure}

\courssub{Consignes}

\begin{enumerate}
\item En comparant les séquences d'ADN du document 3, identifie le type de mutation responsable de la drépanocytose (nature de la mutation, position, conséquence sur la protéine).
\item À partir des documents 1 et 2, détermine les génotypes des parents I-1 et I-2 et de l'enfant II-3. Justifie pourquoi les parents sont dits \emph{sains} bien que porteurs de l'allèle muté.
\item Construis l'échiquier de croisement I-1 $\times$ I-2 et calcule, pour chaque grossesse future, le risque (en fraction et en pourcentage) d'avoir : (a) un enfant drépanocytaire ; (b) un enfant porteur sain ; (c) un enfant non porteur.
\item À partir du document 4, identifie l'anomalie chromosomique du cousin, nomme-la et explique le mécanisme méiotique qui l'a produite. Le document 5 permet-il de comprendre pourquoi le risque était élevé dans son cas ?
\item Rédige un discours de sensibilisation d'environ 10 lignes, destiné au chef de famille Ngo Bell, pour expliquer scientifiquement l'origine de ces anomalies et éradiquer les accusations de sorcellerie. Ton discours doit utiliser au moins quatre mots-clés de la leçon (mutation, allèle, hétérozygote, chromosome, méiose, hasard, etc.).
\end{enumerate}

\courssub{Ressources à mobiliser}

\begin{itemize}
\item Une \cle{mutation génique} est une modification de la séquence d'un gène : substitution, délétion ou insertion. Une substitution peut changer un codon et donc un acide aminé (mutation faux-sens).
\item La drépanocytose est due à la substitution $\text{A} \rightarrow \text{T}$ dans le codon 6 du gène de la $\beta$-globine : GAG (acide glutamique) devient GTG (valine). L'allèle $S$ est \cle{codominant} pour l'électrophorèse, mais le phénotype « sain » domine à l'échelle de l'individu : le génotype $A/S$ est un \cle{porteur sain}.
\item Échiquier de croisement : chaque parent $A/S$ produit 1/2 de gamètes $A$ et 1/2 de gamètes $S$ (ségrégation des allèles à la méiose).
\item La \cle{trisomie 21} résulte d'une \cle{non-disjonction} du chromosome 21 au cours de la méiose (le plus souvent maternelle, en méiose I), produisant un gamète à 24 chromosomes ; après fécondation, le zygote possède 47 chromosomes.
\item Le brassage génétique (méiose + fécondation) explique la diversité des enfants d'un même couple : chaque naissance est un tirage indépendant, régi par le \cle{hasard}, non par une volonté surnaturelle.
\end{itemize}

\courssub{Démarche et corrigé commenté}

\begin{exoresolu}[Conseil génétique : le cas de la famille Ngo Bell]
Énoncé : identifier la mutation de la drépanocytose, déterminer les génotypes de la famille Ngo Bell, calculer les risques pour les grossesses futures, diagnostiquer l'anomalie du cousin et rédiger un discours de sensibilisation.
\tcblower
\textbf{Étape 1 : identification de la mutation (document 3).}

L'allèle normal porte le codon \textbf{GAG} en position 6 ; l'allèle muté porte \textbf{GTG}. Une seule base a changé : l'adénine (A) est remplacée par la thymine (T). Il s'agit donc d'une \cle{mutation ponctuelle par substitution}. Conséquence sur la protéine : le codon GAG code l'\cle{acide glutamique} (acide aminé hydrophile), le codon GTG code la \cle{valine} (hydrophobe). Cette substitution Glu $\rightarrow$ Val en position 6 de la $\beta$-globine modifie les propriétés de l'hémoglobine : l'HbS polymérise en fibres à faible concentration en dioxygène, ce qui déforme les globules rouges en faucilles (d'où « drépanocytose »). Une seule base modifiée sur les milliards que compte le génome suffit donc à créer un \cle{allèle nouveau} et une protéine anormale.

\textbf{Étape 2 : génotypes de la famille (documents 1 et 2).}

L'électrophorèse distingue les deux formes d'hémoglobine : HbA (normale, migration rapide) et HbS (mutée, migration lente).
\begin{itemize}
\item Junior (II-3) ne possède que la bande HbS : il est \cle{homozygote} $S/S$ : c'est l'individu drépanocytaire.
\item Les parents I-1 et I-2 possèdent les deux bandes HbA et HbS : ils sont \cle{hétérozygotes} $A/S$. Comme ils produisent assez d'HbA fonctionnelle, ils sont cliniquement \textbf{sains} : ce sont des \cle{porteurs sains}. L'allèle $S$ est récessif au niveau du phénotype clinique mais détectable (codominant) à l'électrophorèse.
\end{itemize}
Conclusion cohérente avec l'arbre généalogique : deux parents sains peuvent avoir un enfant atteint si chacun transmet son allèle $S$. Nulle sorcellerie là-dedans : une simple loi de probabilité.

\textbf{Étape 3 : échiquier de croisement et risques (question 3).}

Gamètes de chaque parent $A/S$ : 1/2 $A$ et 1/2 $S$ (ségrégation des allèles à la méiose).

\begin{center}
\begin{tabular}{|c|c|c|}
\hline
\rowcolor{popblueL} \textbf{Gamètes} & $A$ (1/2) & $S$ (1/2) \\
\hline
$A$ (1/2) & $A/A$ sain non porteur (1/4) & $A/S$ porteur sain (1/4) \\
\hline
$S$ (1/2) & $A/S$ porteur sain (1/4) & $S/S$ drépanocytaire (1/4) \\
\hline
\end{tabular}
\end{center}

Pour \textbf{chaque grossesse} (tirages indépendants) :
\begin{align*}
P(\text{enfant atteint } S/S) &= \tfrac{1}{4} = 25\ \% \\
P(\text{porteur sain } A/S) &= \tfrac{2}{4} = 50\ \% \\
P(\text{non porteur } A/A) &= \tfrac{1}{4} = 25\ \%
\end{align*}

\textbf{Étape 4 : anomalie chromosomique du cousin (documents 4 et 5).}

Le caryotype montre \textbf{47 chromosomes avec trois exemplaires du chromosome 21} (47, XY, +21) : c'est la \cle{trisomie 21}, responsable du \cle{syndrome de Down}. Mécanisme : lors de la méiose (le plus souvent chez la mère, en première division), les deux chromosomes 21 ne se séparent pas : c'est une \cle{non-disjonction}. Il en résulte un ovule à 24 chromosomes (deux chromosomes 21) ; sa fécondation par un spermatozoïde normal (23 chromosomes) donne un zygote à 47 chromosomes. Le document 5 montre que la fréquence passe d'environ 1/1500 à 20 ans à environ 1/30 à 45 ans : la mère du cousin avait 43 ans, le risque était donc élevé (de l'ordre de 1/50), ce qui s'explique par le vieillissement des ovocytes, bloqués en prophase I depuis la vie fœtale de la mère. Il s'agit d'un accident méiotique lié à l'âge, pas d'une malédiction.

\textbf{Étape 5 : discours de sensibilisation (proposition de corrigé).}

« Grand-père, membres de la famille, écoutez-moi. La maladie de Junior n'est ni une sorcellerie ni une faute de sa mère : c'est une \cle{mutation}, un petit changement dans un gène, présent chez de très nombreux Camerounais sans qu'ils le sachent. Le père et la mère sont tous deux \cle{hétérozygotes}, porteurs sains d'un \cle{allèle} $S$ : ils sont en bonne santé, mais à chaque grossesse, le \cle{hasard} de la \cle{méiose} et de la fécondation donne une chance sur quatre que l'enfant reçoive l'allèle $S$ des deux parents. Le syndrome de Down de notre cousin vient d'un accident naturel de division d'un \cle{chromosome}, dont le risque augmente avec l'âge de la mère. Ces anomalies ont des causes scientifiques connues, démontrées au laboratoire. Accuser une épouse, c'est blesser une innocente et négliger le vrai combat : le suivi médical, le dépistage et le conseil génétique. Aimons nos enfants, soignons-les, et laissons la science éclairer notre famille. »
\end{exoresolu}

\begin{attention}[Piège classique : l'indépendance des naissances]
Le risque de 25 \% est \textbf{le même à chaque grossesse}, quel que soit le nombre d'enfants déjà atteints. Dire « nous avons déjà un enfant malade, les trois prochains seront sains » est une erreur : les naissances sont des événements indépendants, car chaque fécondation est un tirage au sort nouveau des gamètes. De même, les deux enfants aînés sains (II-1 et II-2) ne « protègent » en rien les grossesses suivantes.
\end{attention}

\begin{attention}[Vocabulaire du Bac : ne pas confondre]
\begin{itemize}
\item \cle{Mutation génique} (modification de la séquence d'un gène, invisible sur caryotype) et \cle{mutation chromosomique} (modification du nombre ou de la structure des chromosomes, visible sur caryotype) : la drépanocytose relève de la première, la trisomie 21 de la seconde.
\item \cle{Porteur sain} ($A/S$, hétérozygote, cliniquement indemne) et \cle{malade} ($S/S$, homozygote) : un porteur sain n'est pas « un peu malade », il est sain mais peut transmettre l'allèle muté.
\item La trisomie 21 n'est \textbf{pas héréditaire} dans sa forme classique : elle résulte d'un accident méiotique survenu chez un parent dont les propres chromosomes sont normaux.
\end{itemize}
\end{attention}

\courssub{Bilan de l'activité}

Cette activité t'a permis de mobiliser de manière intégrée l'ensemble des savoir-faire de la leçon :
\begin{itemize}
\item la \cle{comparaison de séquences d'ADN} a identifié une mutation ponctuelle par substitution (GAG $\rightarrow$ GTG), montrant qu'une seule base modifiée suffit à créer un allèle nouveau et une protéine anormale ;
\item l'\cle{électrophorèse} et l'arbre généalogique ont établi les génotypes $A/S$ des parents porteurs sains et $S/S$ de l'enfant malade ;
\item l'\cle{échiquier de croisement} a chiffré le risque : 25 \% d'enfants atteints à chaque grossesse, de façon indépendante ;
\item l'analyse du \cle{caryotype} a diagnostiqué une trisomie 21, expliquée par une non-disjonction méiotique dont la fréquence croît avec l'âge maternel ;
\item enfin, la rédaction du discours a montré que les anomalies génétiques ont des \cle{causes naturelles, mesurables et démontrables} : la science est l'outil le plus efficace pour combattre les préjugés, protéger les personnes injustement accusées et orienter les familles vers le dépistage et le conseil génétique.
\end{itemize}

\begin{aretenir}[L'essentiel de l'activité]
Les mutations géniques (drépanocytose) et chromosomiques (trisomie 21) sont des événements naturels du hasard biologique. Deux parents sains hétérozygotes ont 25 \% de risque d'enfant atteint à chaque grossesse, indépendamment des naissances précédentes. Comprendre ces mécanismes, c'est éradiquer les accusations de sorcellerie et promouvoir le dépistage et le conseil génétique.
\end{aretenir}

\newpage

\coursec{Méthodes et exercices résolus}

\courssub{Méthodes et exercices résolus}

\begin{methode}[Identifier le type de mutation (génique ou chromosomique) et en prédire la conséquence]
    \textbf{Objectif :} Classer une mutation à l'échelle du gène ou du chromosome et en déduire l'impact sur le phénotype.
    \begin{enumerate}
        \item \textbf{Analyser l'échelle de l'altération :}
              \begin{itemize}
                  \item Si la modification porte sur \textbf{un ou quelques nucléotides} (substitution, insertion, délétion) $\rightarrow$ \cle{Mutation génique} (allélique).
                  \item Si la modification porte sur \textbf{un segment de chromosome} (délétion, duplication, inversion, translocation) ou sur \textbf{le nombre de chromosomes} (non-disjonction) $\rightarrow$ \cle{Mutation chromosomique}.
              \end{itemize}
        \item \textbf{Préciser la nature moléculaire (pour une mutation génique) :}
              \begin{itemize}
                  \item \textbf{Substitution :} un nucléotide remplacé (transition/purine-purine ou pyrimidine-pyrimidine ; transversion/purine-pyrimidine). Conséquence : silencieuse, missens (changement d'acide aminé), non-sens (codon STOP).
                  \item \textbf{Insertion/Délétion :} ajout/perte de nucléotides. Si multiple de 3 $\rightarrow$ ajout/perte d'acides aminés ; sinon $\rightarrow$ \textbf{mutation à décalage de cadre} (frameshift), modification radicale de la protéine.
              \end{itemize}
        \item \textbf{Préciser la nature (pour une mutation chromosomique) :}
              \begin{itemize}
                  \item \textbf{Numérique :} Non-disjonction $\rightarrow$ aneuploïdie (trisomie, monosomie) ou polyploïdie.
                  \item \textbf{Structurale :} Cassure et réparation aberrante (délétion, duplication, inversion, translocation réciproque ou robertsonienne).
              \end{itemize}
        \item \textbf{Conclure sur la conséquence phénotypique :} Relier le type de mutation à l'altération de la protéine (perte de fonction, gain de fonction toxique, absence de protéine) ou au déséquilibre génomique (effet de dose).
    \end{enumerate}
    \textbf{Réflexe Bac :} Toujours citer l'expérience fondatrice (ex: expérience de \textbf{Beadle et Tatum} sur \emph{Neurospora crassa} : ``un gène = une enzyme/protéine'') pour justifier le lien mutation génique $\rightarrow$ protéine altérée.
\end{methode}

\begin{exoresolu}[Classification de mutations et conséquences phénotypiques]
    \textbf{Énoncé :} Le tableau ci-dessous présente quatre anomalies génétiques observées chez l'homme. Pour chacune, précisez s'il s'agit d'une mutation génique ou chromosomique, sa nature précise et sa conséquence principale au niveau moléculaire.
    \begin{center}
    \begin{tabularx}{\linewidth}{|l|X|X|X|}\hline
    \rowcolor{popblueL}
    \textbf{Anomalie} & \textbf{Description moléculaire / caryotypique} & \textbf{Type \& Nature} & \textbf{Conséquence moléculaire principale} \\ \hline
    Drépanocytose & Substitution A $\rightarrow$ T au 6ème codon du gène $\beta$-globine (GAG $\rightarrow$ GTG) & & \\ \hline
    Syndrome de Down (Trisomie 21) & Caryotype 47,XY,+21 (ou 47,XX,+21) & & \\ \hline
    Myopathie de Duchenne & Délétion d'un ou plusieurs exons du gène \emph{DMD} (chromosome X) & & \\ \hline
    Syndrome de Cri du chat & Délétion terminale du bras court du chromosome 5 : 46,XX,del(5)(p15) & & \\ \hline
    \end{tabularx}
    \end{center}

    \tcblower
    \textbf{Solution détaillée :}
    \begin{enumerate}
        \item \textbf{Drépanocytose :}
              \begin{itemize}
                  \item \textbf{Type :} Mutation \textbf{génique} (échelle du nucléotide).
                  \item \textbf{Nature :} \textbf{Substitution} (transversion A $\rightarrow$ T) au premier nucléotide du 6ème codon.
                  \item \textbf{Conséquence :} Mutation \textbf{missens}. Le codon GAG (Glutamate, acide, polaire) devient GTG (Valine, hydrophobe). Cela modifie la structure tertiaire de l'hémoglobine (HbS) qui polymérise en condition d'hypoxie $\rightarrow$ déformation des hématies (drepanocytes).
              \end{itemize}
        \item \textbf{Syndrome de Down (Trisomie 21) :}
              \begin{itemize}
                  \item \textbf{Type :} Mutation \textbf{chromosomique numérique}.
                  \item \textbf{Nature :} \textbf{Trisomie} (aneuploïdie) du chromosome 21 due à une \textbf{non-disjonction} (le plus souvent en méiose I maternelle).
                  \item \textbf{Conséquence :} \textbf{Effet de dose} (surdosage génique). Les gènes du chromosome 21 sont exprimés en 3 exemplaires au lieu de 2 $\rightarrow$ perturbation du développement (retard mental, traits morphologiques caractéristiques, cardiopathies).
              \end{itemize}
        \item \textbf{Myopathie de Duchenne :}
              \begin{itemize}
                  \item \textbf{Type :} \textbf{Anomalie génique} (mutation du gène \emph{DMD}). La délétion d'exons entiers, bien que visible en biologie moléculaire (MLPA, CGH-array), altère la séquence codante d'un gène unique.
                  \item \textbf{Nature :} \textbf{Délétion} d'exons (souvent hors cadre, non multiple de 3).
                  \item \textbf{Conséquence :} \textbf{Mutation à décalage de cadre (frameshift)} $\rightarrow$ apparition précoce d'un codon STOP $\rightarrow$ protéine dystrophine tronquée, non fonctionnelle $\rightarrow$ dégénérescence musculaire progressive.
              \end{itemize}
        \item \textbf{Syndrome de Cri du chat :}
              \begin{itemize}
                  \item \textbf{Type :} Mutation \textbf{chromosomique structurale}.
                  \item \textbf{Nature :} \textbf{Délétion terminale} du bras court du chromosome 5 [del(5p)].
                  \item \textbf{Conséquence :} \textbf{Haploinsuffisance} (perte d'un allèle pour de nombreux gènes contigus) $\rightarrow$ retard de croissance, microcéphalie, cri strident du nourrisson (anomalie laryngée).
              \end{itemize}
    \end{enumerate}
    \textbf{Conclusion :} Les mutations géniques altèrent la séquence codante d'un gène (protéine modifiée/absente), tandis que les mutations chromosomiques modifient la dose génique (nombre ou structure des chromosomes) affectant l'expression de centaines de gènes simultanément.
\end{exoresolu}

\begin{methode}[Identifier une anomalie génique par analyse comparative de séquences d'ADN (ex: Drépanocytose)]
    \textbf{Objectif :} Repérer la mutation ponctuelle responsable d'une maladie en comparant l'allèle normal (sauvage) et l'allèle mutant.
    \begin{enumerate}
        \item \textbf{Aligner les séquences :} Placer la séquence de référence (allèle sauvage, noté $A$ ou $Hb^A$) au-dessus de la séquence mutante (allèle $S$ ou $Hb^S$).
        \item \textbf{Repérer la différence :} Identifier le nucléotide modifié (position, base originelle $\rightarrow$ base mutée).
        \item \textbf{Déterminer le type de mutation :} Substitution (transition/transversion), insertion, délétion.
        \item \textbf{Traduire en conséquence codante :}
              \begin{itemize}
                  \item Utiliser le code génétique (tableau des codons).
                  \item Identifier le codon concerné $\rightarrow$ Acide aminé normal vs Acide aminé mutant.
                  \item Classer : Silencieuse / Missens / Non-sens / Décalage de cadre.
              \end{itemize}
        \item \textbf{Relier au phénotype :} Expliquer comment le changement biochimique de l'acide aminé altère la fonction de la protéine (ex: hydrophobie valine vs charge glutamate $\rightarrow$ polymérisation HbS).
    \end{enumerate}
    \textbf{Code à retenir :} Code génétique universel (ex: GAG=Glu, GTG=Val, GUG=Val). Sens de lecture 5' $\rightarrow$ 3' (brin codant = brin sens, T remplacé par U pour l'ARNm).
\end{methode}

\begin{exoresolu}[Analyse de la mutation responsable de la drépanocytose]
    \textbf{Énoncé :} On donne ci-dessous un extrait de la séquence codante (brin sens, 5' $\rightarrow$ 3') du gène $\beta$-globine (\emph{HBB}) pour l'allèle normal ($Hb^A$) et l'allèle drépanocytaire ($Hb^S$), au niveau du 6ème codon.
    \begin{center}
    \begin{tabular}{|c|l|}\hline
    \rowcolor{popblueL}
    \textbf{Allèle} & \textbf{Séquence ADN (brin codant)} \\ \hline
    $Hb^A$ (Normal) & ... \textbf{CCT GAG GAG} ... \\ \hline
    $Hb^S$ (Drépanocytose) & ... \textbf{CCT GTG GAG} ... \\ \hline
    \end{tabular}
    \end{center}
    \textbf{Questions :}
    \begin{enumerate}
        \item Identifier la mutation ponctuelle (position, bases concernées, type).
        \item En déduire la modification de la séquence protéique (utiliser le code génétique : GAG=Glu, GTG=Val, GAG=Glu).
        \item Expliquer le mécanisme biophysique à l'origine de la drépanocytose en condition d'hypoxie.
    \end{enumerate}

    \tcblower
    \textbf{Solution détaillée :}
    \begin{enumerate}
        \item \textbf{Identification de la mutation :}
              \begin{itemize}
                  \item \textbf{Position :} 1ère base du 2ème codon présenté (6ème codon du gène mature).
                  \item \textbf{Modification :} Substitution de l'adénine (\textbf{A}) par la thymine (\textbf{T}). Écriture : \textbf{A $\rightarrow$ T}.
                  \item \textbf{Type :} \textbf{Transversion} (purine A $\rightarrow$ pyrimidine T). C'est une \textbf{substitution} simple.
              \end{itemize}
        \item \textbf{Conséquence sur la protéine (Traduction) :}
              \begin{itemize}
                  \item Codon normal (Allèle $A$) : \textbf{GAG} $\xrightarrow{\text{Code génétique}}$ \textbf{Glutamate (Glu, E)}. Acide aminé chargé négativement, polaire, hydrophile.
                  \item Codon mutant (Allèle $S$) : \textbf{GTG} $\xrightarrow{\text{Code génétique}}$ \textbf{Valine (Val, V)}. Acide aminé à chaîne latérale hydrophobe, neutre.
                  \item \textbf{Classification :} Mutation \textbf{missens} (changement de sens) non conservative (changement radical de propriétés physico-chimiques).
              \end{itemize}
        \item \textbf{Mécanisme biophysique (Physiopathologie) :}
              \begin{itemize}
                  \item En condition \textbf{normoxique}, l'HbS ($\alpha_2\beta_2^S$) est soluble, identique à l'HbA.
                  \item En condition d'\textbf{hypoxie} (pression partielle en O$_2$ basse), l'HbS libère son O$_2$ $\rightarrow$ changement conformationnel (état T, désoxygéné).
                  \item La Valine en position 6 ($\beta$6Val), hydrophobe, s'expose à la surface de la protéine. Elle s'insère dans une \textbf{poche hydrophobe} complémentaire située sur une autre molécule d'HbS (interaction $\beta$6Val - Poche $\beta$85Phe/$\beta$88Leu).
                  \item Cela entraîne une \textbf{polymérisation} de l'HbS en longs polymères rigides (fibres de 14 brins).
                  \item Ces fibres déforment l'hématie en \textbf{drepanocyte} (forme de faucille), la rigidifient $\rightarrow$ hémolyse (anémie hémolytique) et vaso-occlusion (crises douloureuses, atteintes viscérales).
              \end{itemize}
    \end{enumerate}
    \begin{popfigure}
    \begin{tikzpicture}[scale=0.8, every node/.style={font=\small}]
        % Normal HbA
        \node[draw, rounded corners, fill=popgreenL, minimum width=2.5cm, minimum height=1.2cm] (HbA) at (0,0) {\textbf{HbA (Normoxie / Hypoxie)}};
        \node[circle, draw, fill=popgreen, minimum size=0.6cm] (RBC_A) at (0,-1.8) {};
        \node[below=0.1cm of RBC_A, align=center]{Hématie discoïde\\(souple)};
        \draw[->, thick, popgreen] (HbA.south) -- (RBC_A.north);
        
        % Mutant HbS
        \node[draw, rounded corners, fill=poporangeL, minimum width=2.5cm, minimum height=1.2cm] (HbS) at (6,0) {\textbf{HbS (Normoxie)}};
        \node[circle, draw, fill=poporange, minimum size=0.6cm] (RBC_S_n) at (6,-1.8) {};
        \node[below=0.1cm of RBC_S_n] {Hématie discoïde};
        \draw[->, thick, poporange] (HbS.south) -- (RBC_S_n.north);
        
        \node[draw, rounded corners, fill=popred!25, minimum width=3cm, minimum height=1.2cm] (HbS_hyp) at (12,0) {\textbf{HbS (Hypoxie)}};
        % Polymer fibers
        \foreach \i in {-0.8,-0.4,0,0.4,0.8} {
            \draw[popred, thick, decorate, decoration={snake, amplitude=1pt, segment length=4pt}] (12+\i, -0.6) -- (12+\i, -1.5);
        }
        \node[draw, ellipse, fill=popred, minimum width=1.5cm, minimum height=0.8cm, rotate=-20] (RBC_S_h) at (12,-2) {};
        \node[below=0.1cm of RBC_S_h, align=center]{Drepanocyte\\(rigide)};
        \draw[->, thick, popred] (HbS_hyp.south) -- (RBC_S_h.north);
        
        % Legend
        \node[align=left, anchor=north west] at (-3.5, -3.5) {
            \textbf{Légende :} \\
            1. HbA : Globine $\beta$6Glu (hydrophile) $\rightarrow$ pas d'agrégation. \\
            2. HbS Normoxie : Conformation R (oxygénée), Val6 enterrée $\rightarrow$ soluble. \\
            3. HbS Hypoxie : Conformation T (désoxygénée), Val6 exposée $\rightarrow$ polymérisation en fibres rigides $\rightarrow$ déformation cellulaire.
        };
    \end{tikzpicture}
    \legende{Mécanisme biophysique de la drépanocytose : polymérisation conditionnelle de l'HbS en hypoxie.}
    \end{popfigure}
    \textbf{Conclusion :} Une substitution unique (SNP) change un acide aminé hydrophile en hydrophobe, créant un site d'agrégation intermoléculaire conditionnel (dépendant de l'oxygénation), base moléculaire de la maladie.
\end{exoresolu}

\begin{savaistu}[Drépanocytose au Cameroun : un enjeu de santé publique]
Au Cameroun, la prévalence du trait drépanocytaire (hétérozygote AS) est estimée entre \textbf{20\% et 30\%} de la population générale, avec des pics dans les zones forestières du Sud et de l'Est. La drépanocytose majeure (SS) touche environ \textbf{1 à 2\% des nouveau-nés} (soit 6 000 à 12 000 naissances/an). Depuis 2019, le \textbf{dépistage néonatal systématique} (test de Guthrie sur papier buvard, électrophorèse hémoglobines) est déployé dans les grandes villes (Yaoundé, Douala, Garoua) avec l'appui de l'IEMP (Institut d'Épidémiologie et de Médecine Préventive) et de la Fondation Pierre Fabre. La prise en charge précoce (pénicilline V prophylactique, vaccins anti-pneumocoque/Haemophilus, hydroxyurée, éducation thérapeutique) a fait passer l'espérance de vie de < 5 ans à > 40 ans. Le \textbf{conseil génétique prénuptial} (certificat de génotype Hb) est recommandé par le Ministère de la Santé Publique pour réduire l'incidence des unions à risque (AS $\times$ AS).
\end{savaistu}

\begin{methode}[Identifier une anomalie chromosomique à partir d'un caryotype]
    \textbf{Objectif :} Lire un caryotype (notation ISCN) ou observer une image de chromosomes classés pour diagnostiquer une anomalie numérique ou structurale.
    \begin{enumerate}
        \item \textbf{Compter le nombre total de chromosomes :}
              \begin{itemize}
                  \item 46 = nombre normal (diploïdie humaine).
                  \item 47 = \textbf{Trisomie} (chromosome en excès).
                  \item 45 = \textbf{Monosomie} (chromosome manquant).
              \end{itemize}
        \item \textbf{Analyser les chromosomes sexuels (gonosomes) :}
              \begin{itemize}
                  \item XX = Féminin ; XY = Masculin.
                  \item XO (45,X) = \textbf{Syndrome de Turner} (monosomie X).
                  \item XXY (47,XXY) = \textbf{Syndrome de Klinefelter} (trisomie sexuelle mâle).
                  \item XYY (47,XYY) = Syndrome de l'XYY (souvent asymptomatique).
              \end{itemize}
        \item \textbf{Identifier le chromosome autosomique en excès (si 47,XX ou 47,XY) :}
              \begin{itemize}
                  \item +21 = \textbf{Trisomie 21} (Syndrome de Down). Le plus fréquent (1/700 naissances).
                  \item +18 = Trisomie 18 (Syndrome d'Edwards, létal précoce).
                  \item +13 = Trisomie 13 (Syndrome de Patau, létal précoce).
              \end{itemize}
        \item \textbf{Repérer les anomalies structurales (notation : del, dup, inv, t, r) :}
              \begin{itemize}
                  \item \textbf{del(chr)(bras bande)} : Délétion (perte de matériel). Ex: del(5)(p15) $\rightarrow$ Cri du chat.
                  \item \textbf{t(chr1;chr2)(bande1;bande2)} : Translocation (échange). Si robertsonienne (fusion centromères acrocentriques) $\rightarrow$ risque de trisomie 21 par translocation chez la descendance.
              \end{itemize}
        \item \textbf{Rédiger le diagnostic :} Nom du syndrome + Formule caryotypique complète (ex: 47,XX,+21).
    \end{enumerate}
    \textbf{Réflexe Bac :} La non-disjonction peut survenir en \textbf{Méiose I} (séparation des chromosomes homologues) ou \textbf{Méiose II} (séparation des chromatides). Conséquence : Gamète n+1 ou n-1. Fécondation par gamète normal (n) $\rightarrow$ Zygote 2n+1 (trisomie) ou 2n-1 (monosomie).
\end{methode}

\begin{exoresolu}[Diagnostic caryotypique : Trisomie 21, Turner, Klinefelter]
    \textbf{Énoncé :} Trois caryotypes sont réalisés au Centre de Génétique Médicale de l'Hôpital Central de Yaoundé. La notation ISCN simplifiée est donnée.
    \begin{enumerate}
        \item \textbf{Cas 1 :} 47,XX,+21
        \item \textbf{Cas 2 :} 45,X
        \item \textbf{Cas 3 :} 47,XXY
    \end{enumerate}
    Pour chaque cas :
    \begin{enumerate}
        \item Donner le nombre de chromosomes et le sexe chromosomique.
        \item Identifier l'anomalie (numérique/structurale) et nommer le syndrome.
        \item Préciser l'origine méiotique la plus probable (Méiose I ou II ? Père ou Mère ?) pour la trisomie 21.
        \item Citer deux signes cliniques majeurs pour chaque syndrome.
    \end{enumerate}

    \tcblower
    \textbf{Solution détaillée :}
    \begin{center}
    \begin{tabularx}{\linewidth}{|l|X|X|X|}\hline
    \rowcolor{popblueL}
    \textbf{Critère} & \textbf{Cas 1 : 47,XX,+21} & \textbf{Cas 2 : 45,X} & \textbf{Cas 3 : 47,XXY} \\ \hline
    \textbf{Nombre total} & 47 & 45 & 47 \\ \hline
    \textbf{Sexe chromosomique} & Féminin (XX) & Féminin (XO) & Masculin (XXY) \\ \hline
    \textbf{Type d'anomalie} & Numérique : Trisomie autosomique & Numérique : Monosomie gonosomique & Numérique : Trisomie gonosomique \\ \hline
    \textbf{Syndrome} & \textbf{Syndrome de Down (Trisomie 21)} & \textbf{Syndrome de Turner} & \textbf{Syndrome de Klinefelter} \\ \hline
    \textbf{Origine méiotique probable} & \textbf{Non-disjonction méiotique maternelle} (95\% des cas). \newline Principalement \textbf{Méiose I} (séparation homologues 21). Âge maternel avancé = facteur de risque majeur. & \textbf{Non-disjonction paternelle} (perte du chromosome sexuel Y ou X) $\rightarrow$ spermatozoïde nullo (sans chromosome sexuel) fécondant ovocyte X normal (~70-80\% des cas). \newline Ou non-disjonction maternelle (Méiose I ou II) $\rightarrow$ ovocyte nullo fécondé par spermatozoïde X. & \textbf{Non-disjonction méiotique} : \newline - Maternelle (Méiose I $\rightarrow$ ovocyte XX) + spermatozoïde Y. \newline - Paternelle (Méiose I $\rightarrow$ spermatozoïde XY) + ovocyte X. \newline Âge maternel avancé facteur de risque. \\ \hline
    \textbf{Signes cliniques majeurs (2)} & 1. Retard mental modéré + traits faciaux caractéristiques (pli épicanthique, face plate, langue proéminente). \newline 2. Cardiopathie congénitale fréquente (CIA, CIV, canal atrio-ventriculaire 40-50\%) + hypotonie néonatale. & 1. \textbf{Dysgénie gonadique} (ovaires streaks) $\rightarrow$ infertilité, absence de puberté spontanée, aménorrhée primaire. \newline 2. \textbf{Ptérygion colli} (pli cutané cou), petite taille, lymphedème néonatal (pieds/mains), anomalies cardiovasculaires (coarctation aorte 10-20\%). & 1. \textbf{Hypogonadisme hypergonadotrope} : testicules petits/fermes, azoospermie (infertilité), gynécomastie, taux testostérone bas, FSH/LH élevés. \newline 2. Grande taille (eunuchoïde), troubles du langage/apprentissage légers, risque accru maladies auto-immunes/cancer sein. \\ \hline
    \end{tabularx}
    \end{center}
    \textbf{Conclusion :} L'analyse du nombre total de chromosomes et de la formule des gonosomes permet un diagnostic caryotypique certain. La trisomie 21 est la seule trisomie autosomique viable à l'âge adulte ; les monosomies autosomiques sont létales in utero. Le syndrome de Turner (monosomie X) est la seule monosomie viable.
\end{exoresolu}

\begin{methode}[Expliquer l'origine de la diversité génétique et l'unicité des individus (Brassage allélique)]
    \textbf{Objectif :} Démontrer que la méiose et la fécondation sont les moteurs de la diversité génétique (sauf jumeaux monozygotes).
    \begin{enumerate}
        \item \textbf{Identifier les trois sources de brassage :}
              \begin{enumerate}
                  \item \textbf{Brassage interchromosomique (Indépendance des paires) :} Aléa de la répartition des chromosomes homologues en Méiose I (Plaque équatoriale). $2^n$ combinaisons possibles ($n$ = nombre de paires = 23 chez l'homme $\rightarrow$ $2^{23} \approx 8,4 \times 10^6$).
                  \item \textbf{Brassage intrachromosomique (Crossing-over) :} Échanges de segments entre chromatides non-sœurs de chromosomes homologues (Chiasmes) en Prophase I. Crée des chromosomes \textbf{recombinés} (nouveaux haplotypes alléliques). Quasi-infini.
                  \item \textbf{Fécondation aléatoire :} Rencontre fortuite d'un spermatozoïde (parmi $\sim 10^8$) et d'un ovocyte (1 par cycle). Multiplie les combinaisons : $(2^n \times \text{Crossing-over})^2$.
              \end{enumerate}
        \item \textbf{Calculer le nombre théorique de gamètes différents (sans crossing-over) :} $N = 2^n$.
        \item \textbf{Argumenter l'unicité :} Probabilité que deux individus (non jumeaux) aient le même génotype = quasi nulle ($< 1/70\,000$ milliards). Même les fratries partagent en moyenne 50\% de leurs allèles identiques par descendance.
        \item \textbf{Distinction Mutation / Brassage :} La \textbf{mutation} crée les \textbf{nouveaux allèles} (variabilité \emph{de novo}). Le \textbf{brassage} recombine les \textbf{allèles existants} dans de nouveaux génotypes (variabilité \emph{combinatoire}).
    \end{enumerate}
    \textbf{Schéma mental :} Méiose I (Anaphase I : séparation homologues $\rightarrow$ Brassage inter) + Prophase I (Chiasmes $\rightarrow$ Brassage intra) $\rightarrow$ Gamètes haploïdes uniques $\rightarrow$ Fécondation $\rightarrow$ Zygote unique.
\end{methode}

\begin{exoresolu}[Calcul de la diversité gamétique et argumentation sur l'unicité]
    \textbf{Énoncé :} Chez l'homme, $n = 23$ paires de chromosomes.
    \begin{enumerate}
        \item Calculer le nombre théorique de combinaisons chromosomiques différentes possibles dans les gamètes dus au seul brassage interchromosomique. Exprimer en notation scientifique.
        \item En intégrant le crossing-over (en moyenne 2 à 3 chiasmes par chromosome), expliquer pourquoi le nombre réel de gamètes génétiquement distincts est virtuellement infini.
        \item Un couple a déjà 3 enfants. Quelle est la probabilité qu'un 4ème enfant ait \textbf{exactement} le même génotype nucléaire que l'aîné ? Justifier.
        \item Expliquer pourquoi on dit que « la mutation est la source ultime de la diversité, le brassage en est le mélangeur ».
    \end{enumerate}

    \tcblower
    \textbf{Solution détaillée :}
    \begin{enumerate}
        \item \textbf{Brassage interchromosomique seul :}
              \[ N = 2^n = 2^{23} = 8\,388\,608 \approx \mathbf{8,4 \times 10^6} \text{ combinaisons chromosomiques différentes par gamète.} \]
              (Chaque paire se sépare indépendamment : 2 choix par paire $\times$ 23 paires).
        \item \textbf{Rôle du Crossing-over :}
              Le crossing-over se produit en Prophase I entre chromatides non-sœurs. Il brise le lien physique entre allèles d'un même chromosome.
              \begin{itemize}
                  \item En moyenne 50 à 70 chiasmes par méiose chez l'homme (1 à 3 par bras de chromosome).
                  \item Chaque chiasme crée 2 chromatides recombinées sur 4.
                  \item Le nombre de combinaisons alléliques \emph{intra}-chromosomiques devient astronomique (dépend du nombre de gènes hétérozygotes et de leur distance).
              \end{itemize}
              \textbf{Conclusion :} Le nombre de génotypes gamétiques possibles dépasse largement $10^{10}$, rendant chaque gamète virtuellement unique.
        \item \textbf{Probabilité d'identité génotypique entre frères/sœurs :}
              \begin{itemize}
                  \item Pour un locus donné, deux frères partagent 0, 1 ou 2 allèles identiques par descendance (probabilités 1/4, 1/2, 1/4).
                  \item Probabilité d'avoir le \textbf{même génotype} à un locus = 1/2 (si parents hétérozygotes distincts) ou 1 (si parents homozygotes).
                  \item Sur l'ensemble du génome ($\sim 20\,000$ gènes, dont des milliers hétérozygotes), la probabilité est le produit des probabilités par locus $\rightarrow$ \textbf{quasi nulle} (de l'ordre de $1/2^{k}$ où $k$ est le nombre de locus hétérozygotes informatifs, $k \gg 100$).
              \end{itemize}
              \textbf{Réponse :} Probabilité \textbf{pratiquement nulle} (inférieure à $10^{-30}$). Seuls les jumeaux monozygotes (division précoce du zygote) partagent 100\% de leur ADN nucléaire.
        \item \textbf{Mutation vs Brassage :}
              \begin{itemize}
                  \item \textbf{Mutation :} Événement rare, aléatoire, crée un \textbf{nouvel allèle} (ex: allèle $Hb^S$ apparu par substitution sur allèle $Hb^A$ ancestral). Sans mutation, pas de polymorphisme, pas d'évolution. C'est la \textbf{source ultime}.
                  \item \textbf{Brassage :} Événement systématique à chaque méiose. Ne crée pas de nouveaux allèles, mais de \textbf{nouvelles combinaisons d'allèles} (nouveaux haplotypes, nouveaux génotypes). C'est le \textbf{mélangeur} qui génère la diversité phénotypique sur laquelle la sélection agit.
              \end{itemize}
    \end{enumerate}
\end{exoresolu}

\begin{methode}[Analyser un arbre généalogique pour déterminer le mode de transmission d'une anomalie génique]
    \textbf{Objectif :} Déduire le mode de transmission (Autosomal Récessif/Dominant, Gonosomique Récessif/Dominant) et les génotypes probables.
    \begin{enumerate}
        \item \textbf{Observer la répartition selon le sexe :}
              \begin{itemize}
                  \item Atteints $\approx$ égaux M/F $\rightarrow$ \textbf{Autosomal}.
                  \item Atteints $\gg$ chez les hommes $\rightarrow$ \textbf{Gonosomique (lié à l'X)}.
              \end{itemize}
        \item \textbf{Observer la transmission verticale (génération en génération) :}
              \begin{itemize}
                  \item Maladie présente à chaque génération $\rightarrow$ \textbf{Dominant}.
                  \item Maladie saute des générations (parents sains $\rightarrow$ enfants malades) $\rightarrow$ \textbf{Récessif}.
              \end{itemize}
        \item \textbf{Tester les transmissions critiques (Gonosomique) :}
              \begin{itemize}
                  \item \textbf{Père atteint $\rightarrow$ Fils atteint ?} Si OUI $\rightarrow$ Autosomal (Père donne Y à fils). Si NON (Père atteint $\rightarrow$ Filles toutes atteintes, Fils tous sains) $\rightarrow$ \textbf{X-lié Dominant}.
                  \item \textbf{Mère atteinte (hétérozygote) $\rightarrow$ Fils atteints (50\%) ?} $\rightarrow$ \textbf{X-lié Récessif ou Dominant}.
                  \item \textbf{Père sain + Mère porteuse $\rightarrow$ Fils atteints} $\rightarrow$ \textbf{X-lié Récessif}.
              \end{itemize}
        \item \textbf{Déduire les génotypes :} Utiliser symboles clairs : $A$ (allèle normal), $a$ (allèle mutant autosomique) ; $X^A$, $X^a$ (allèles sur X) ; $Y$.
        \item \textbf{Calculer les risques de récurrence :} Croisement des gamètes parentaux (Tableau de Punnett).
    \end{enumerate}
    \textbf{Piège :} Une maladie autosomique récessive rare peut \emph{paraître} X-liée si peu de cas. Vérifier la transmission père $\rightarrow$ fils.
\end{methode}

\begin{exoresolu}[Analyse de pedigree : Hémophilie A (Facteur VIII)]
    \textbf{Énoncé :} On étudie une famille camerounaise touchée par l'hémophilie A (déficit en facteur VIII). L'arbre généalogique simplifié est le suivant :
    \begin{itemize}
        \item \textbf{Génération I :} Père (I-1) sain, Mère (I-2) saine.
        \item \textbf{Génération II :} 3 enfants : Fille (II-1) saine, Fils (II-2) \textbf{atteint}, Fille (II-3) saine.
        \item \textbf{Génération III :} La fille II-1 (saine) épouse un homme sain (II-4). Ils ont 2 fils : III-1 \textbf{atteint}, III-2 sain. Une fille III-3 saine.
    \end{itemize}
    \textbf{Questions :}
    \begin{enumerate}
        \item Déterminer le mode de transmission (Autosomal/Gonosomique, Dominant/Récessif). Justifier par 3 arguments tirés du pedigree.
        \item Déterminer les génotypes de tous les individus (I-1, I-2, II-1, II-2, II-3, II-4, III-1, III-2, III-3). Noter $X^H$ allèle normal, $X^h$ allèle hémophilie.
        \item Calculer la probabilité qu'un futur fils de II-1 soit atteint.
        \item Expliquer pourquoi le père I-1 ne peut pas transmettre la maladie à son fils II-2.
    \end{enumerate}

    \tcblower
    \textbf{Solution détaillée :}
    \begin{enumerate}
        \item \textbf{Mode de transmission : \textbf{Gonosomique (lié à l'X) Récessif}.}
            \textbf{Arguments :}
            \begin{enumerate}
                \item \textbf{Saut de génération :} Parents I-1 et I-2 sains $\rightarrow$ fils II-2 atteint. Exclut le Dominant (un parent aurait dû être atteint).
                \item \textbf{Biais sexuel :} Seuls des \textbf{hommes} sont atteints (II-2, III-1). Les femmes (I-2, II-1, II-3, III-3) sont toutes phénotypiquement saines. Cela exclut l'autosomal récessif (aurait dû toucher 25\% des filles aussi) et l'X-dominant (filles atteintes).
                \item \textbf{Transmission "en oblique" (Croisée) :} Le grand-père maternel (I-1) est sain. Le père de III-1 (II-4) est sain. Le fils atteint (III-1) a hérité de l'allèle mutant de sa mère (II-1), qui l'a hérité de \textbf{sa mère (I-2)}. I-2 est une femme saine $\rightarrow$ \textbf{Porteuse hétérozygote} ($X^H X^h$). Elle a transmis $X^h$ à son fils II-2 (atteint) et à sa fille II-1 (porteuse). II-1 a transmis $X^h$ à son fils III-1 (atteint). C'est le schéma classique X-récessif.
            \end{enumerate}
        \item \textbf{Génotypes (Allèles : $X^H$ = normal, $X^h$ = hémophilie) :}
            \begin{center}
            \begin{tabular}{|c|l|l|}\hline
            \rowcolor{popblueL}
            \textbf{Individu} & \textbf{Phénotype} & \textbf{Génotype} \\ \hline
            I-1 (Père G1) & Sain & $X^H Y$ \\ \hline
            I-2 (Mère G1) & Sain (Porteuse) & $X^H X^h$ \\ \hline
            II-1 (Fille aînée) & Sain (Porteuse) & $X^H X^h$ (a reçu $X^h$ de sa mère I-2) \\ \hline
            II-2 (Fils aîné) & \textbf{Atteint} & $X^h Y$ (a reçu $X^h$ de sa mère I-2) \\ \hline
            II-3 (Fille cadette) & Sain & $X^H X^H$ ou $X^H X^h$ (probabilité 1/2 porteuse). \textbf{Notation :} $X^H X^{?}$ \\ \hline
            II-4 (Époux II-1) & Sain & $X^H Y$ \\ \hline
            III-1 (Petit-fils 1) & \textbf{Atteint} & $X^h Y$ (a reçu $X^h$ de sa mère II-1) \\ \hline
            III-2 (Petit-fils 2) & Sain & $X^H Y$ \\ \hline
            III-3 (Petite-fille) & Sain & $X^H X^H$ ou $X^H X^h$ (probabilité 1/2 porteuse). \textbf{Notation :} $X^H X^{?}$ \\ \hline
            \end{tabular}
            \end{center}
        \item \textbf{Risque de récurrence pour un fils de II-1 :}
            II-1 est porteuse ($X^H X^h$). Elle produit 50\% gamètes $X^H$, 50\% $X^h$.
            Père II-4 ($X^H Y$) produit 50\% $X^H$, 50\% $Y$.
            Pour un \textbf{fils} : il reçoit $Y$ du père. Il reçoit $X^H$ ou $X^h$ de la mère.
            $P(\text{Fils atteint } | \text{ Fils}) = P(\text{Mère donne } X^h) = \mathbf{1/2 \text{ (soit 50\%)}}$.
        \item \textbf{Pourquoi I-1 ne transmet pas à II-2 :}
            Le père (I-1) est $X^H Y$. Il transmet son chromosome \textbf{Y} à ses \textbf{fils} (détermination du sexe mâle) et son chromosome \textbf{X} ($X^H$) à ses \textbf{filles}.
            Le fils II-2 a reçu le Y de son père. Il a donc \textbf{obligatoirement} reçu son unique chromosome X de sa \textbf{mère (I-2)}. C'est la preuve formelle que l'allèle mutant vient de la mère.
    \end{enumerate}
    \begin{popfigure}
    \begin{tikzpicture}[scale=0.7, every node/.style={font=\small}]
        % Gen I
        \node[draw, rectangle, minimum size=0.8cm] (I1) at (0,3) {I-1};
        \node[draw, circle, minimum size=0.8cm] (I2) at (2,3) {I-2};
        \node[above=0.1cm of I1] {$X^H Y$};
        \node[above=0.1cm of I2] {$X^H X^h$};
        \draw (I1) -- (I2);
        
        % Gen II
        \node[draw, circle, minimum size=0.8cm] (II1) at (-2,1) {II-1};
        \node[draw, rectangle, minimum size=0.8cm, fill=popred!50] (II2) at (0,1) {II-2};
        \node[draw, circle, minimum size=0.8cm] (II3) at (2,1) {II-3};
        \node[above=0.1cm of II1] {$X^H X^h$};
        \node[above=0.1cm of II2] {$X^h Y$};
        \node[above=0.1cm of II3] {$X^H X^{?}$};
        \draw (I1) -- (0,2) -- (II1);
        \draw (0,2) -- (II2);
        \draw (0,2) -- (II3);
        
        % Gen III
        \node[draw, rectangle, minimum size=0.8cm] (II4) at (-4,1) {II-4};
        \node[above=0.1cm of II4] {$X^H Y$};
        \draw (II4) -- (II1);
        
        \node[draw, rectangle, minimum size=0.8cm, fill=popred!50] (III1) at (-4,-1) {III-1};
        \node[draw, rectangle, minimum size=0.8cm] (III2) at (-2,-1) {III-2};
        \node[draw, circle, minimum size=0.8cm] (III3) at (0,-1) {III-3};
        \node[above=0.1cm of III1] {$X^h Y$};
        \node[above=0.1cm of III2] {$X^H Y$};
        \node[above=0.1cm of III3] {$X^H X^{?}$};
        \draw (II4) -- (-3,0) -- (III1);
        \draw (-3,0) -- (III2);
        \draw (-3,0) -- (III3);
        
        % Legend
        \node[align=left, anchor=north west] at (3.5, 3) {
            \textbf{Légende :} \\
            $\square$ = Homme ; $\bigcirc$ = Femme \\
            \textcolor{popred!50}{$\blacksquare$} = Atteint (Hémophilie A) \\
            $\bigcirc$ hachuré = Porteuse (non dessiné ici, voir génotypes) \\
            $X^H$ = Allèle normal ; $X^h$ = Allèle mutant
        };
    \end{tikzpicture}
    \legende{Arbre généalogique de la famille (Hémophilie A, transmission X-récessive).}
    \end{popfigure}
\end{exoresolu}

\begin{methode}[Construire un argumentaire scientifique pour éradiquer les préjugés (Sensibilisation)]
    \textbf{Objectif :} Utiliser les connaissances génétiques (hasard mutation, lois de Mendel, absence de contagion) pour déconstruire les croyances non scientifiques (sorcellerie, malédiction, faute maternelle, contagion).
    \begin{enumerate}
        \item \textbf{Rappeler l'origine biologique :} « Cette anomalie résulte d'une \textbf{erreur aléatoire} (mutation spontanée ou non-disjonction accidentelle) survenue lors de la formation des gamètes ou au tout début du développement embryonnaire. »
        \item \textbf{Exonérer les parents (surtout la mère) :} « Ce n'est la faute de personne. Ni le père, ni la mère n'ont "fait" quelque chose de mal (alimentation, comportement, pensées). L'âge maternel est un \textbf{facteur de risque statistique} pour la non-disjonction (Trisomie 21), pas une cause volontaire. »
        \item \textbf{Refuser la contagion / hérédité "magique" :} « Ce n'est pas une maladie infectieuse (pas de virus/bactérie). On ne l'attrape pas par contact. C'est une constitution génétique fixée à la fécondation. »
        \item \textbf{Expliquer la récurrence (si applicable) :} « Pour les maladies récessives (drépanocytose), les parents sont \textbf{porteurs sains} (hétérozygotes). Le risque de récurrence est de \textbf{25\%} à chaque grossesse, indépendant des grossesses précédentes (indépendance des événements). Ce n'est pas une "malédiction familiale". »
        \item \textbf{Valoriser la prise en charge médicale et sociale :} « Le diagnostic précoce (dépistage néonatal drépanocytose, phénylcétonurie au Cameroun), le conseil génétique, le suivi médical (transfusions, hydroxyurée, kiné, éducation spécialisée) améliorent drastiquement la qualité de vie et l'espérance de vie. L'inclusion scolaire et sociale est un droit. »
    \end{enumerate}
    \textbf{Ton :} Empathique, factuel, non culpabilisant, basé sur des preuves (caryotype, séquençage, biochimie).
\end{methode}

\begin{exoresolu}[Rédaction d'un message de sensibilisation : Cas de la Trisomie 21]
    \textbf{Énoncé :} Dans un village de la région de l'Ouest, une jeune mère de 24 ans accouche d'un enfant porteur de trisomie 21. La belle-famille accuse la mère d'avoir « mangé un aliment interdit » ou d'avoir « été ensorcelée » pendant la grossesse, et exige qu'elle soit renvoyée. En tant qu'agent de santé formé en SVT, vous devez rédiger une note explicative courte (150-200 mots) à l'intention de la famille, utilisant \textbf{uniquement des arguments scientifiques} pour apaiser les tensions et protéger la mère et l'enfant.

    \tcblower
    \textbf{Proposition de note de sensibilisation (Modèle de rédaction Bac) :}

    \begin{definition}[Note d'information médicale : Origine de la Trisomie 21]
    \textbf{Objet :} Explication scientifique de la survenue de la trisomie 21 chez le nouveau-né de Mme [Nom].

    \textbf{Chers membres de la famille,}

    Le diagnostic de trisomie 21 (syndrome de Down) a été confirmé par caryotype (47,XX,+21 ou 47,XY,+21). Il est essentiel de comprendre l'origine biologique de cette anomalie pour écarter toute accusation infondée.

    \textbf{1. Un accident cellulaire aléatoire, non une faute.}
    La trisomie 21 résulte d'une \textbf{non-disjonction} : lors de la formation de l'ovocyte (méiose), les deux chromosomes 21 n'ont pas migré vers des pôles opposés. L'ovocyte a ainsi conservé deux chromosomes 21 au lieu d'un. À la fécondation, le spermatozoïde (apportant un chromosome 21) a restauré le nombre total à 47 au lieu de 46. Cet événement est \textbf{aléatoire} (hasard biophysique), survenant spontanément dans $\approx$ 1/700 grossesses. Il n'est \textbf{lié à aucune alimentation, aucun comportement, aucune pensée, ni aucun "sort"} de la mère pendant la grossesse.

    \textbf{2. L'âge maternel jeune ne protège pas totalement.}
    Bien que la fréquence augmente avec l'âge de la mère (risque accru après 35 ans), \textbf{la majorité des enfants trisomiques 21 naissent de mères jeunes} (car il y a plus de grossesses chez les jeunes femmes). Mme [Nom], âgée de 24 ans, entre dans cette statistique normale.

    \textbf{3. Non-contagiosité et non-hérédité classique (sauf translocation rare).}
    Dans 95\% des cas (non-disjonction standard), le caryotype des parents est normal (46,XX / 46,XY). Le risque de récurrence pour une grossesse ultérieure est faible ($\approx$ 1\%), bien supérieur au risque de population mais sans lien avec une "malédiction familiale".

    \textbf{4. Prise en charge et avenir.}
    Un suivi médical précoce (cardiologie, ORL, ophtalmologie, pédiatrie, stimulation psychomotrice) permet aujourd'hui à ces enfants de scolariser, d'apprendre un métier et de vivre longtemps. L'amour familial et l'inclusion sociale sont les meilleurs traitements.

    \textbf{Conclusion :} La science nous éclaire : cet enfant est un don de la vie, porteur d'un chromosome en plus par le pur hasard de la méiose. Accueillons-le avec la dignité et les soins qu'il mérite.
    \end{definition}
\end{exoresolu}

\begin{methode}[Différencier les conséquences des mutations géniques et chromosomiques sur le phénotype]
    \textbf{Objectif :} Comparer l'impact à l'échelle de l'individu (gravité, étendue des symptômes, viabilité).
    \begin{enumerate}
        \item \textbf{Mutation Génique (Monogénique) :}
              \begin{itemize}
                  \item \textbf{Cible :} Un seul gène $\rightarrow$ Une seule protéine (ou quelques isoformes).
                  \item \textbf{Phénotype :} Souvent \textbf{spécifique}, ciblé (ex: Hb anormale $\rightarrow$ anémie ; CFTR muté $\rightarrow$ mucus épais poumons/pancréas ; Facteur VIII $\rightarrow$ hémophilie).
                  \item \textbf{Viabilité :} Souvent compatible avec la vie (surtout récessives en hétérozygote). Maladies "rares" individuellement, fréquentes collectivement.
                  \item \textbf{Hétérozygote :} Souvent asymptomatique (récessif) ou phénotype intermédiaire (codominance drépanocytose : résistance paludisme).
              \end{itemize}
        \item \textbf{Mutation Chromosomique :}
              \begin{itemize}
                  \item \textbf{Cible :} Des centaines à des milliers de gènes (déséquilibre génomique global).
                  \item \textbf{Phénotype :} \textbf{Pléiotropie sévère, syndromique} (atteinte multi-systèmes : morphologie, cognition, organes vitaux). Ex: Trisomie 21 (cœur, intestin, immunité, faciès, QI) ; Turner (taille, ovaires, cœur, rein) ; Klinefelter (testicules, hormone, apprentissage).
                  \item \textbf{Viabilité :} Monosomies autosomiques = \textbf{Létales in utero}. Trisomies autosomiques viables = seulement 13, 18, 21 (gènes "tolérants" au surdosage). Aneuploïdies gonosomiques (XO, XXY, XYY, XXX) = viables (inactivation X / gènes rares sur Y).
                  \item \textbf{Pénétrance/Expressivité :} Variable mais toujours multisystémique.
              \end{itemize}
        \item \textbf{Synthèse pour le Bac :} « La mutation génique altère une \textbf{fonction précise} (une protéine) ; la mutation chromosomique perturbe l'\textbf{équilibre global du génome} (dosage de centaines de gènes), entraînant un syndrome complexe. »
    \end{enumerate}
\end{methode}

\begin{exoresolu}[Comparaison Drépanocytose (Génique) vs Trisomie 21 (Chromosomique)]
    \textbf{Énoncé :} Complétez le tableau comparatif suivant pour opposer les anomalies géniques (ex: Drépanocytose) et chromosomiques (ex: Trisomie 21).

    \tcblower
    \textbf{Solution détaillée :}
    \begin{center}
    \begin{tabularx}{\linewidth}{|l|X|X|}\hline
    \rowcolor{popblueL}
    \textbf{Critère de comparaison} & \textbf{Anomalie Génique (Ex: Drépanocytose - HbS)} & \textbf{Anomalie Chromosomique (Ex: Trisomie 21)} \\ \hline
    \textbf{Échelle de l'altération} & Moléculaire : 1 nucléotide (SNP) / 1 gène (\emph{HBB} chr11). & Cellulaire : 1 chromosome entier en excès (Chr 21 $\approx$ 300-400 gènes). \\ \hline
    \textbf{Mécanisme moléculaire} & Substitution missens (Glu6Val) $\rightarrow$ Protéine anormale (HbS) polymérisant en hypoxie. & Non-disjonction méiotique $\rightarrow$ Surdosage génique (Effet de dose) $\rightarrow$ Déséquilibre expression de $\sim$300 gènes. \\ \hline
    \textbf{Type d'hérédité} & Mendélienne : \textbf{Autosomique Récessive} (Hétérozygote AS = Porteur sain / Résistant paludisme). & Non-mendélienne (Accident méiotique). \textbf{Risque de récurrence} faible ($\sim$1\%) sauf translocation parentale. \\ \hline
    \textbf{Phénotype (Étendue)} & \textbf{Ciblé / Spécifique :} Hématies (anémie hémolytique, vaso-occlusion), Rate (autosplénectomie), Infections. \newline Atteinte mono-organique dominante (globule rouge). & \textbf{Syndromique / Pléiotrope :} Dysmorphie faciale, Retard mental, Cardiopathie (50\%), Hypotonie, Anomalies digestives, Immunodéficit, Leucémie risque accru. \newline Atteinte multi-systémique. \\ \hline
    \textbf{Gravité / Viabilité} & Variable. Formes majeures (SS) : Espérance de vie améliorée (40-60 ans au Cameroun avec hydroxyurée/vaccins). Formes mineures (AS) : Asymptomatique. & Viable. Espérance de vie $\approx$ 60 ans (progrès chirurgie cardiaque, prise en charge). Retard mental constant (variable). \\ \hline
    \textbf{Dépistage / Diagnostic} & Biochimie : Électrophorèse hémoglobines (HbA, HbS, HbF). \newline Moléculaire : PCR/RFLP ou Séquençage (mutation ponctuelle). & Cytogénétique : \textbf{Caryotype} (47,XX,+21 ou 47,XY,+21). \newline Prénatal : Amniocentèse / Trophobiopsie + QF-PCR / FISH / Séquençage haut débit (NIPT). \\ \hline
    \textbf{Prévention / Conseil} & Dépistage prénuptial / prénatal (statut AS). Conseil génétique : Risque 25\% si couple AS $\times$ AS. & Dépistage combiné 1er trimestre (Marqueurs sériques + Clarté nucale) $\rightarrow$ Diagnostic prénatal invasif si risque élevé. \newline Conseil : Âge maternel, antécédents. \\ \hline
    \end{tabularx}
    \end{center}
    \textbf{Conclusion :} La drépanocytose illustre la pathologie \textbf{monogénique} (protéine unique défaillante), traitée par approche moléculaire/cellulaire. La trisomie 21 illustre la \textbf{pathologie chromosomique} (déséquilibre génomique global), traitée par approche syndromique multidisciplinaire. Les deux nécessitent un \textbf{conseil génétique} rigoureux pour lutter contre les préjugés.
\end{exoresolu}

\begin{attention}[Les 5 erreurs qui coûtent des points au Bac]
    \begin{enumerate}
        \item \textbf{Confondre "Mutation" et "Brassage" :} Écrire que « la méiose crée de nouveaux allèles ». \textbf{Faux.} La méiose \textbf{brasse} les allèles existants (créant de nouveaux génotypes). Seule la \textbf{mutation} (erreur réplication, mutagène) crée un \textbf{nouvel allèle} (nouvelle séquence ADN). \textit{Sanction : -1 à -2 pts sur la définition de la diversité.}
        \item \textbf{Oublier le Crossing-over dans le calcul de la diversité :} Répondre $2^{23}$ pour le nombre de gamètes possibles \textbf{sans préciser} « sans crossing-over » ou « par brassage interchromosomique seul ». \textbf{Correct :} « Le brassage interchromosomique seul donne $2^{23}$ combinaisons ; le crossing-over rend le nombre virtuellement infini. » \textit{Sanction : -0,5 à -1 pt (imprécision scientifique).}
        \item \textbf{Erreur de lecture de pedigree (Transmission Père $\rightarrow$ Fils) :} Conclure à une transmission X-liée (récessive ou dominante) alors qu'un \textbf{père atteint transmet le trait à son fils}. \textbf{Règle d'or :} Le père donne son \textbf{Y} à son fils. Transmission père $\rightarrow$ fils $\rightarrow$ \textbf{Obligatoirement Autosomique}. \textit{Sanction : 0 pt sur la question "Mode de transmission".}
        \item \textbf{Vocabulaire imprécis sur la Trisomie 21 :} Dire « Le chromosome 21 est \textbf{dupliqué} » ou « Il y a une \textbf{mutation du chromosome 21} ». \textbf{Termes exacts :} « \textbf{Non-disjonction} du chromosome 21 $\rightarrow$ \textbf{Trisomie 21} (3 exemplaires au lieu de 2) $\rightarrow$ \textbf{Surdosage génique / Effet de dose} ». « Duplication » = anomalie structurale (segment). « Mutation » = échelle génique. \textit{Sanction : -1 pt (vocabulaire non scientifique).}
        \item \textbf{Accuser la mère / Cause environnementale directe :} Dans la partie "Éradication des préjugés", écrire « La trisomie 21 est due à l'âge de la mère » sans nuance, ou « La drépanocytose est due à une piqûre de moustique / mauvaise alimentation ». \textbf{Exigence Bac :} Argumentation \textbf{strictement biologique} : Hasard méiotique (non-disjonction), Héritage mendélien (allèles parents), Absence de contagion. Toute référence à la sorcellerie, la faute maternelle, l'alimentation (sauf phénylcétonurie \emph{traitement}), vaut \textbf{0 pt} sur la compétence "Savoir-être / Sensibilisation".
    \end{enumerate}
\end{attention}

\newpage

\coursec{Exercices}

\courssub{Je vérifie mes connaissances}

\begin{exercice}[QCM : Types de mutations et conséquences]
Pour chacune des affirmations suivantes, une seule réponse est correcte. Recopie le numéro et la lettre correspondante.
\begin{enumerate}
    \item Une mutation génique est :
    \begin{itemize}
        \item[A.] Une modification du nombre de chromosomes.
        \item[B.] Une modification de la structure d'un chromosome (délétion, inversion, translocation).
        \item[C.] Une modification de la séquence nucléotidique d'un gène (substitution, insertion, délétion de quelques bases).
        \item[D.] Une non-disjonction chromosomique au cours de la méiose.
    \end{itemize}
    \item La drépanocytose (hémoglobine S) résulte d'une mutation ponctuelle dans le gène codant la $\beta$-globine (codon 6 : GAG $\to$ GTG). Cette mutation est une :
    \begin{itemize}
        \item[A.] Substitution non-sens (apparition d'un codon STOP).
        \item[B.] Substitution synonyme (même acide aminé).
        \item[C.] Substitution non-synonyme (changement d'acide aminé : Glu $\to$ Val).
        \item[D.] Insertion d'un nucléotide entraînant un décalage du cadre de lecture.
    \end{itemize}
    \item La trisomie 21 (syndrome de Down) est le plus souvent due à :
    \begin{itemize}
        \item[A.] Une translocation Robertsonnienne du chromosome 21 sur le 14 (forme familiale).
        \item[B.] Une non-disjonction de la paire 21 lors de la méiose I ou II chez la mère (âge maternel avancé).
        \item[C.] Une délétion du bras long du chromosome 21.
        \item[D.] Une mutation du gène \emph{APP} situé sur le chromosome 21.
    \end{itemize}
    \item Le brassage des allèles (interchromosomique et intrachromosomique) est une source majeure de variabilité génétique. Il se produit lors de :
    \begin{itemize}
        \item[A.] La mitose (division cellulaire somatique).
        \item[B.] La fécondation (fusion des gamètes).
        \item[C.] La méiose (répartition aléatoire des chromosomes homologues et crossing-over).
        \item[D.] La réplication de l'ADN (phase S du cycle cellulaire).
    \end{itemize}
\end{enumerate}
\end{exercice}

\begin{exercice}[Vrai / Faux justifié : Anomalies géniques et chromosomiques]
Indique si chaque affirmation est Vraie (V) ou Fausse (F) et justifie brièvement ta réponse par un argument scientifique précis.
\begin{enumerate}
    \item L'hémophilie A (déficit en facteur VIII) est une maladie autosomique récessive touchant indifféremment les deux sexes.
    \item Le syndrome de Turner (monosomie X, 45,X) ne concerne que les femmes et entraîne une stérilité quasi systématique par dysgénésie gonadique.
    \item Un individu porteur sain de la drépanocytose (génotype $Hb^A/Hb^S$) présente une anémie hémolytique sévère en conditions normales de vie.
    \item La mucoviscidose est due à une mutation du gène \emph{CFTR} (chromosome 7) ; la mutation $\Delta$F508 (délétion de 3 bases) est la plus fréquente en population caucasienne.
    \item La non-disjonction peut survenir aussi bien lors de la méiose I (séparation des chromosomes homologues) que lors de la méiose II (séparation des chromatides sœurs).
\end{enumerate}
\end{exercice}

\begin{exercice}[Définitions et notions clés]
Donne une définition précise et rigoureuse (niveau Bac) pour chacun des termes suivants :
\begin{enumerate}
    \item Mutation génique (ou mutation ponctuelle).
    \item Mutation chromosomique (structurale et numérique).
    \item Non-disjonction.
    \item Brassage interchromosomique et brassage intrachromosomique (crossing-over).
    \item Allèle mutant / Allèle sauvage.
    \item Pénétrance et expressivité d'un allèle mutant.
\end{enumerate}
\end{exercice}

\begin{exercice}[Calculs directs : Probabilités de transmission]
\begin{enumerate}
    \item Un couple de porteurs sains pour la drépanocytose (maladie autosomique récessive, allèle $S$ mutant, allèle $A$ normal) désire avoir un enfant. Calcule la probabilité (en pourcentage) que cet enfant soit : a) atteint (drépanocytaire majeur) ; b) porteur sain ; c) non porteur de l'allèle mutant.
    \item Une femme porteuse saine de l'hémophilie A (maladie liée à l'X, récessive, allèle $X^h$ mutant, allèle $X^H$ normal) a un enfant avec un homme sain ($X^H Y$). Calcule la probabilité : a) d'avoir un garçon atteint ; b) d'avoir une fille porteuse ; c) d'avoir un enfant atteint (tous sexes confondus).
    \item Lors de la méiose, un individu diploïde $2n=6$ (3 paires de chromosomes homologues) ne subit aucun crossing-over. Combien de combinaisons chromosomiques différentes (brassage interchromosomique seul) peut-il produire dans ses gamètes ? Applique la formule générale.
\end{enumerate}
\end{exercice}

\courssub{Je m'entraîne}

\begin{exercice}[Analyse comparée de séquences ADN : La drépanocytose]
On donne ci-dessous un extrait de la séquence codante (brin sens) de l'exon 1 du gène $\beta$-globine (\emph{HBB}) pour un individu sain (témoin) et pour un patient suspecté de drépanocytose. Les codons sont séparés par des espaces pour faciliter la lecture.

\textbf{Individu sain (Témoin) :}
\verb|... CTG ACT CCT GAG GAG AAG TCT ...|

\textbf{Patient :}
\verb|... CTG ACT CCT GTG GAG AAG TCT ...|

Tableau du code génétique (extraits utiles) :
\begin{center}
\begin{tabular}{|c|c|c|c|}\hline
Codon & Acide aminé & Codon & Acide aminé \\ \hline
GAG & Glu (Acide glutamique) & GTG & Val (Valine) \\
GAA & Glu & GTA & Val \\
CTG & Leu (Leucine) & ACT & Thr (Thréonine) \\
CCT & Pro (Proline) & AAG & Lys (Lysine) \\
TCT & Ser (Sérine) & & \\ \hline
\end{tabular}
\end{center}

\begin{enumerate}
    \item Identifie le type de mutation ponctuelle survenue (substitution, insertion, délétion). Précise le nucléotide modifié et sa position dans le codon.
    \item Traduis les deux séquences en séquences d'acides aminés (utilise les abréviations à 3 lettres).
    \item Quel est l'acide aminé modifié ? Quelle est la nature physico-chimique du changement (acide $\to$ basique, polaire $\to$ apolaire, etc.) ?
    \item Conclus sur la nature de l'anomalie moléculaire responsable de la drépanocytose. Pourquoi cette mutation entraîne-t-elle une polymérisation de l'hémoglobine S en conditions d'hypoxie ?
\end{enumerate}
\end{exercice}

\begin{exercice}[Interprétation de caryotypes : Identification de syndromes chromosomiques]
Trois caryotypes anonymisés (A, B, C) ont été réalisés à partir de lymphocytes en culture (bande G). Les observations au microscope permettent d'établir les formules caryotypiques suivantes :
\begin{itemize}
    \item \textbf{Caryotype A :} 47 chromosomes au total. Présence de trois chromosomes de la paire 21. Paires sexuelles : XX.
    \item \textbf{Caryotype B :} 45 chromosomes au total. Absence d'un chromosome sexuel. Un seul chromosome X présent. Paires autosomiques normales.
    \item \textbf{Caryotype C :} 47 chromosomes au total. Paire sexuelle composée de deux chromosomes X et d'un chromosome Y (XXY). Paires autosomiques normales.
\end{itemize}

\begin{enumerate}
    \item Pour chaque caryotype, écris la formule caryotypique normalisée (ex : 47,XX,+21).
    \item Identifie le syndrome correspondant à chaque caryotype (Trisomie 21, Turner, Klinefelter).
    \item Pour le caryotype A, explique le mécanisme cytogénétique (non-disjonction) à l'origine de l'anomalie. Précise s'il s'agit d'une non-disjonction méiotique I, méiotique II ou mitotique (mosaïcisme) et justifie ton choix par l'âge maternel comme facteur de risque majeur.
    \item Pour le caryotype C, explique pourquoi le phénotype est masculin malgré la présence de deux chromosomes X. Cite deux signes cliniques majeurs observés à la puberté.
    \item Le caryotype B présente un nombre impair de chromosomes (45). Comment se déroule la méiose chez cette patiente ? Quelle en est la conséquence sur la fertilité ?
\end{enumerate}
\end{exercice}

\begin{exercice}[Arbre généalogique : L'hémophilie A dans une famille camerounaise]
On étudie la transmission de l'hémophilie A (maladie récessive liée à l'X, allèle mutant $h$, allèle normal $H$) dans la famille ci-dessous. Les individus atteints sont symbolisés par des symboles pleins (carrés pour les hommes, cercles pour les femmes). Les individus non atteints sont symbolisés par des symboles vides.

\begin{popfigure}
\begin{tikzpicture}[scale=0.7, transform shape]
% Generation I
\node[draw, circle, minimum size=0.8cm] (I1) at (0,3) {I-1};
\node[draw, rectangle, minimum size=0.8cm, fill=black] (I2) at (3,3) {I-2};
% Marriage line
\draw (I1) -- (I2);
% Vertical down
\draw (1.5,3) -- (1.5,2.2);
% Generation II
\node[draw, circle, minimum size=0.8cm] (II1) at (-2,1) {II-1};
\node[draw, rectangle, minimum size=0.8cm, fill=black] (II2) at (1,1) {II-2};
\node[draw, circle, minimum size=0.8cm] (II3) at (4,1) {II-3};
\node[draw, rectangle, minimum size=0.8cm] (II4) at (7,1) {II-4};
% Sibship line
\draw (-2,1) -- (7,1);
% Verticals from marriage
\draw (1.5,2.2) -- (-2,1);
\draw (1.5,2.2) -- (1,1);
\draw (1.5,2.2) -- (4,1);
\draw (1.5,2.2) -- (7,1);
% Marriages Gen II
\draw (II1) -- (-4.5,1); \node[draw, rectangle, minimum size=0.8cm] (II1m) at (-4.5,1) {};
\draw (II3) -- (6.5,1); \node[draw, rectangle, minimum size=0.8cm] (II3m) at (6.5,1) {};
% Generation III (children of II-1)
\draw (-3.25,1) -- (-3.25,0.2);
\node[draw, rectangle, minimum size=0.8cm, fill=black] (III1) at (-4.5,-0.5) {III-1};
\node[draw, circle, minimum size=0.8cm] (III2) at (-2,-0.5) {III-2};
\draw (-4.5,-0.5) -- (-2,-0.5);
% Generation III (children of II-3)
\draw (5.25,1) -- (5.25,0.2);
\node[draw, rectangle, minimum size=0.8cm] (III3) at (4,-0.5) {III-3};
\node[draw, circle, minimum size=0.8cm] (III4) at (6.5,-0.5) {III-4};
\draw (4,-0.5) -- (6.5,-0.5);
\end{tikzpicture}
\legende{Arbre généalogique d'une famille atteinte d'hémophilie A. Symboles pleins = atteints. I-1 et I-2 sont les grands-parents.}
\end{popfigure}

\begin{enumerate}
    \item Détermine le génotype de chaque individu de la génération I (I-1 et I-2).
    \item Détermine les génotypes probables de tous les individus des générations II et III. Justifie chaque déduction par l'analyse des phénotypes et des règles de transmission liée à l'X.
    \item Calcule la probabilité qu'un futur enfant du couple II-3 $\times$ II-3m (homme sain) soit un garçon atteint.
    \item Explique pourquoi, dans cette famille, seuls les hommes sont atteints (phénotype malade) alors que des femmes portent l'allèle mutant.
    \item Propose une méthode de diagnostic prénatal disponible au Cameroun (ex: Centre Mère-Enfant de la Fondation Chantal Biya, Hôpital Général de Yaoundé) pour détecter cette anomalie chez un fœtus masculin à risque.
\end{enumerate}
\end{exercice}

\begin{exercice}[Brassage des allèles et unicité génétique]
On considère un individu diploïde de génotype $A/a ; B/b ; C/c$ pour trois gènes non liés (sur trois paires de chromosomes homologues différentes). On suppose qu'aucun crossing-over ne se produit pour ces loci.
\begin{enumerate}
    \item Combien de types de gamètes différents cet individu peut-il produire par brassage interchromosomique seul ? Liste-les.
    \item Si cet individu s'autoféconde (autofécondation théorique), combien de génotypes différents sont possibles dans la descendance ? Combien de phénotypes différents (en supposant dominance complète $A>a, B>b, C>c$) ?
    \item On considère maintenant le gène $A/a$ seul. Un crossing-over se produit entre le locus $A$ et le centromère lors de la méiose I. Schématise (schéma simplifié) les chromosomes homologues au stade pachyptène (tétades) et au stade anaphase I, en montrant l'échange de segments. Combien de chromatides recombinantes sur les 4 chromatides de la tétade ?
    \item En réalité, le génome humain compte environ 20 000 gènes répartis sur 23 paires de chromosomes. Explique pourquoi le brassage (inter + intra) rend la probabilité d'avoir deux individus génétiquement identiques (hors vrais jumeaux monozygotes) quasi nulle. Donne un ordre de grandeur du nombre de combinaisons possibles chez l'humain ($n=23$).
\end{enumerate}
\end{exercice}

\courssub{Je me prépare au Bac}

\begin{exercice}[Sujet type Bac : Trisomie 21 et âge maternel — Exploitation de documents]
\textbf{Contexte :} Au Cameroun, le dépistage prénatal de la trisomie 21 (syndrome de Down) est proposé aux femmes enceintes, notamment par la mesure de la clarté nucale (échographie du 1er trimestre) et le dosage des marqueurs sériques (PAPP-A, $\beta$-hCG libre), complété si besoin par l'amniocentèse (caryotype fœtal). La figure ci-dessous représente l'évolution de la prévalence de la trisomie 21 à la naissance en fonction de l'âge maternel (données épidémiologiques internationales).

\begin{popfigure}
\begin{tikzpicture}
\begin{axis}[
    width=\linewidth,
    height=8cm,
    xlabel={Âge maternel (années)},
    ylabel={Prévalence de la trisomie 21 (pour 1000 naissances)},
    xmin=15, xmax=50,
    ymin=0, ymax=35,
    xtick={15,20,25,30,35,40,45,50},
    ytick={0,5,10,15,20,25,30,35},
    grid=major,
    grid style={dashed, gray!50},
    axis lines=left,
    tick label style={font=\small},
    label style={font=\small},
]
% Data points approximate: (20, 0.5), (25, 0.7), (30, 1.0), (35, 2.5), (37, 4), (40, 10), (42, 16), (45, 30), (48, 35)
\addplot[popcurve, thick, mark=*, mark size=2, color=popblue] coordinates {
    (20,0.5) (25,0.7) (30,1.0) (32,1.5) (35,2.5) (37,4) (40,10) (42,16) (45,30) (48,35)
};
\end{axis}
\end{tikzpicture}
\legende{Évolution de la prévalence de la trisomie 21 à la naissance en fonction de l'âge maternel.}
\end{popfigure}

\begin{enumerate}
    \item \textbf{Analyse de la courbe :} Décris l'évolution de la prévalence de la trisomie 21 en fonction de l'âge maternel. Cite des valeurs chiffrées précises (avec unités) pour étayer ta description (ex: à 25 ans, à 35 ans, à 45 ans).
    \item \textbf{Interprétation biologique :} Formule une hypothèse explicative reliant l'âge maternel à l'augmentation du risque de non-disjonction chromosomique lors de l'oogenèse. Rappelle le moment où l'oocyte I est bloqué en prophase I de la méiose et ce qui se passe à la puberté puis à chaque cycle ovarien.
    \item \textbf{Calcul de risque :} Une femme de 38 ans consulte pour une grossesse. D'après la courbe, quel est son risque approximatif d'avoir un enfant trisomique 21 (exprimé en \%) ? Si 500 femmes de 38 ans accouchent dans une maternité de Yaoundé, combien d'enfants trisomiques 21 attend-on statistiquement (arrondi à l'entier) ?
    \item \textbf{Prévention et dépistage :} Cite les deux principales techniques de diagnostic prénatal invasif permettant d'obtenir le caryotype fœtal. Donne l'avantage et l'inconvénient majeur de chacune (risque de fausse couche, précocité). Quelle technique non invasive (ADN fœtal circulant) se développe actuellement ?
    \item \textbf{Éthique et éducation :} En tant que futur professionnel de santé ou citoyen éclairé, comment expliquerais-tu à une famille camerounaise que la trisomie 21 n'est \textbf{pas} une « malédiction » ou une « sorcellerie » mais un accident chromosomique aléatoire dont la probabilité augmente avec l'âge maternel ? Quels arguments scientifiques (mécanisme, hasard) et humains (inclusion, prise en charge précoce) utiliserais-tu ?
\end{enumerate}
\end{exercice}

\begin{exercice}[Sujet type Bac : Synthèse argumentée — Mutation vs Brassage]
\textbf{Sujet :} « L'apparition d'un caractère nouveau dans une famille est-elle toujours due à une mutation ? »

Rédige un texte argumenté structuré (introduction, développement en deux parties distinctes, conclusion) d'environ 300 à 400 mots pour répondre à cette question. Tu dois obligatoirement :
\begin{itemize}
    \item Définir rigoureusement ce qu'est une mutation (génique et chromosomique) et son rôle dans la création \textbf{de novo} d'allèles nouveaux.
    \item Expliquer le mécanisme du brassage des allèles (interchromosomique et intrachromosomique/crossing-over) lors de la méiose et son rôle dans la recombinaison d'allèles \textbf{préexistants} dans la population.
    \item Illustrer par un exemple concret (ex: groupe sanguin ABO, phénotype drépanocytaire, couleur des yeux) la différence entre « apparition d'un allèle nouveau par mutation » et « apparition d'un phénotype nouveau dans une fratrie par brassage ».
    \item Conclure sur la complémentarité de ces deux moteurs de la biodiversité : la mutation comme source ultime de nouveauté, le brassage comme générateur immédiat de diversité chez l'individu.
\end{itemize}
La qualité de la rédaction scientifique (vocabulaire précis : allèle, locus, gamète, zygote, crossing-over, non-disjonction), la logique du raisonnement et la pertinence des exemples seront évaluées.
\end{exercice}

\begin{exercice}[Mini-projet d'éducation à la santé : Sensibilisation sur la drépanocytose au Cameroun]
\textbf{Contexte :} La drépanocytose est la maladie génétique la plus fréquente au Cameroun (prévalence du trait drépanocytaire $AS$ : 20 à 25\% de la population générale, jusqu'à 30\% dans certaines zones). Pourtant, de forts préjugés persistent : « maladie de sorcellerie », « enfant sorcier », stigmatisation des couples $AS \times AS$, refus du dépistage prénuptial.

\textbf{Consigne :} Tu es chargé(e) de concevoir le contenu scientifique d'une \textbf{affiche de sensibilisation} (format A3) destinée aux jeunes en âge de se marier et aux femmes enceintes, à afficher dans les centres de santé (ex: CSI de district, Hôpital de District, Centre Pasteur du Cameroun).

L'affiche doit contenir les éléments suivants (organise-les logiquement avec des titres clairs) :
\begin{enumerate}
    \item \textbf{Titre accrocheur} (ex: « La drépanocytose : Connaître son statut, c'est protéger ses enfants »).
    \item \textbf{Origine scientifique :} Explication simple (2-3 phrases) de la mutation du gène $\beta$-globine (chromosome 11), transmission autosomique récessive. Schéma simplifié des génotypes : $AA$ (sain), $AS$ (porteur sain / trait drépanocytaire), $SS$ (drépanocytaire majeur).
    \item \textbf{Risques pour la descendance :} Tableau de croisement $AS \times AS$ (carré de Punnett) avec les probabilités (25\% $SS$, 50\% $AS$, 25\% $AA$). Message clair : « Si les deux partenaires sont $AS$, 1 enfant sur 4 risque d'être malade à chaque grossesse. »
    \item \textbf{Moyens de dépistage disponibles au Cameroun :}
    \begin{itemize}
        \item Électrophorèse de l'hémoglobine (sur hémolysat) : examen de référence, peu coûteux, disponible dans la plupart des laboratoires d'hôpitaux de district et régionaux.
        \item Test de solubilité (Emmel test) : test rapide de terrain (dépistage de masse), à confirmer par électrophorèse.
        \item Dépistage néonatal (talon) : mis en place progressivement (Programme National de Lutte contre la Drépanocytose).
        \item Diagnostic prénatal (biopsie des villosités choriales ou amniocentèse + biologie moléculaire PCR) : centres spécialisés (Yaoundé, Douala).
    \end{itemize}
    \item \textbf{Prise en charge et prévention :} Vaccination anti-pneumocoque (Prevenar 13), prophylaxie par pénicilline V (jusqu'à 5 ans), acide folique, hydratation, hydroxyurée (traitement modificateur de la maladie), suivi en Centre de Prise en Charge de la Drépanocytose (ex: Hôpital Mère-Enfant de la Fondation Chantal Biya, Hôpital Laquintinie Douala).
    \item \textbf{Message de lutte contre les préjugés :} « La drépanocytose n'est ni une malédiction, ni une sorcellerie. C'est une maladie génétique héréditaire. Le dépistage avant le mariage (certificat de génotype) est un acte d'amour et de responsabilité. Ne stigmatisons pas les malades et les porteurs. »
    \item \textbf{Contacts utiles :} Numéros verts / Associations (ex: ASCAD - Association Camerounaise de Lutte contre la Drépanocytose, Programme National de Lutte contre la Drépanocytose - PNLD).
\end{enumerate}
\textbf{Note :} Tu n'as pas à dessiner l'affiche, mais à rédiger le \textbf{contenu textuel complet et structuré} (titres, textes, tableaux, schémas décrits) prêt à être mis en page par un graphiste. Sois précis sur les génotypes, les pourcentages, les noms des structures de santé camerounaises.
\end{exercice}

\newpage

\coursec{Corrigés des exercices}

\begin{corrige}[Exercice 1 — QCM : Types de mutations et conséquences]
\begin{enumerate}
    \item \textbf{Réponse C.} Une mutation génique (ou ponctuelle) est une modification de la séquence nucléotidique d'un gène : substitution, insertion ou délétion d'une ou quelques bases. Les réponses A, B et D décrivent des mutations \emph{chromosomiques} (numériques ou structurales).
    \item \textbf{Réponse C.} Le codon 6 passe de GAG (Glutamate) à GTG (Valine) : il s'agit d'une substitution non-synonyme (mutation faux-sens), car l'acide aminé est modifié (Glu $\to$ Val). Ce n'est ni un codon STOP (non-sens), ni une substitution synonyme, ni un décalage du cadre de lecture (le nombre de bases est inchangé).
    \item \textbf{Réponse B.} Dans environ 95\,\% des cas, la trisomie 21 résulte d'une non-disjonction de la paire 21 lors de la méiose maternelle (I ou II), dont la fréquence augmente avec l'âge maternel. La translocation Robertsonnienne (forme familiale, environ 4\,\%) et le mosaïcisme (environ 1\,\%) sont minoritaires.
    \item \textbf{Réponse C.} Le brassage interchromosomique (répartition aléatoire des chromosomes homologues en anaphase I) et le brassage intrachromosomique (crossing-over en prophase I) ont lieu au cours de la \cle{méiose}. La fécondation amplifie ensuite la diversité par rencontre aléatoire des gamètes, mais ce n'est pas le lieu du brassage proprement dit.
\end{enumerate}
\textbf{Bilan :} 1-C ; 2-C ; 3-B ; 4-C.
\end{corrige}

\begin{corrige}[Exercice 2 — Vrai / Faux justifié]
\begin{enumerate}
    \item \textbf{Faux.} L'hémophilie A est une maladie \emph{récessive liée au chromosome X} (gène \emph{F8} en Xq28). Les garçons ($XY$), hémizygotes, sont atteints dès qu'ils portent l'allèle mutant ; les filles sont le plus souvent porteuses saines ($X^H X^h$). Les deux sexes ne sont donc pas touchés indifféremment.
    \item \textbf{Vrai.} Le syndrome de Turner (45,X) ne concerne que des individus de sexe féminin (un seul X, pas de Y). L'absence du second chromosome sexuel provoque une dysgénésie gonadique (ovaires fibreux, non fonctionnels), d'où une aménorrhée primaire et une stérilité quasi constante.
    \item \textbf{Faux.} Le porteur sain $Hb^A/Hb^S$ (trait drépanocytaire) est asymptomatique en conditions normales : son hémoglobine A suffit à assurer le transport de l'\ce{O2}. Seules des conditions extrêmes (hypoxie sévère, forte altitude, effort intense) peuvent provoquer des signes. L'anémie hémolytique sévère caractérise la forme majeure $Hb^S/Hb^S$.
    \item \textbf{Vrai.} La mucoviscidose est due à des mutations du gène \emph{CFTR} (chromosome 7, transmission autosomique récessive). La mutation $\Delta$F508 (délétion de 3 nucléotides entraînant la perte d'une phénylalanine en position 508) représente environ 70\,\% des allèles mutants en population caucasienne.
    \item \textbf{Vrai.} La non-disjonction peut affecter la méiose I (les deux chromosomes homologues migrent vers le même pôle) ou la méiose II (les deux chromatides sœurs ne se séparent pas). Dans les deux cas, on obtient des gamètes à $n+1$ ou $n-1$ chromosomes.
\end{enumerate}
\end{corrige}

\begin{corrige}[Exercice 3 — Définitions et notions clés]
\begin{enumerate}
    \item \textbf{Mutation génique (ponctuelle) :} modification aléatoire et héritable de la séquence nucléotidique d'un gène (substitution, insertion ou délétion d'une ou quelques bases), pouvant créer un allèle nouveau et modifier la protéine codée.
    \item \textbf{Mutation chromosomique :} modification héritable de la structure d'un chromosome (délétion, duplication, inversion, translocation) — mutations \emph{structurales} — ou du nombre de chromosomes (aneuploïdies comme les trisomies et monosomies, polyploïdies) — mutations \emph{numériques}.
    \item \textbf{Non-disjonction :} défaut de séparation des chromosomes homologues (méiose I) ou des chromatides sœurs (méiose II ou mitose), produisant des cellules filles avec un chromosome en surnombre ($n+1$) ou en défaut ($n-1$).
    \item \textbf{Brassage interchromosomique :} répartition aléatoire et indépendante des chromosomes homologues de chaque paire lors de l'anaphase I de méiose, générant $2^n$ combinaisons de gamètes. \textbf{Brassage intrachromosomique (crossing-over) :} échange réciproque de segments entre chromatides non sœurs de chromosomes homologues en prophase I (stade pachytène), créant des chromosomes recombinés portant de nouvelles combinaisons d'allèles.
    \item \textbf{Allèle mutant :} version d'un gène issue d'une mutation, souvent responsable d'un phénotype modifié ou pathologique. \textbf{Allèle sauvage :} version la plus fréquente dans la population naturelle, assurant la fonction normale.
    \item \textbf{Pénétrance :} proportion (en \%) des individus porteurs d'un génotype donné qui expriment effectivement le phénotype correspondant. \textbf{Expressivité :} degré d'intensité variable avec lequel le phénotype s'exprime chez les individus porteurs.
\end{enumerate}
\end{corrige}

\begin{corrige}[Exercice 4 — Calculs directs : Probabilités de transmission]
\begin{enumerate}
    \item Croisement $A/S \times A/S$ (drépanocytose, autosomique récessive). Carré de Punnett :
    \begin{center}
    \begin{tabular}{|c|c|c|}\hline
    \rowcolor{popblueL} Gamètes & $A$ (1/2) & $S$ (1/2) \\ \hline
    $A$ (1/2) & $A/A$ (1/4) & $A/S$ (1/4) \\ \hline
    $S$ (1/2) & $A/S$ (1/4) & $S/S$ (1/4) \\ \hline
    \end{tabular}
    \end{center}
    a) Enfant atteint $S/S$ : $\dfrac{1}{4} = \mathbf{25\,\%}$. \quad b) Porteur sain $A/S$ : $\dfrac{1}{2} = \mathbf{50\,\%}$. \quad c) Non porteur $A/A$ : $\dfrac{1}{4} = \mathbf{25\,\%}$.
    \item Croisement $X^H X^h \times X^H Y$ (hémophilie A, récessive liée à l'X) :
    \begin{center}
    \begin{tabular}{|c|c|c|}\hline
    \rowcolor{popblueL} & $X^H$ (1/2) & $Y$ (1/2) \\ \hline
    $X^H$ (1/2) & $X^H X^H$ fille saine (1/4) & $X^H Y$ garçon sain (1/4) \\ \hline
    $X^h$ (1/2) & $X^H X^h$ fille porteuse (1/4) & $X^h Y$ garçon atteint (1/4) \\ \hline
    \end{tabular}
    \end{center}
    a) Garçon atteint : $\dfrac{1}{4} = \mathbf{25\,\%}$ de tous les enfants (soit 1 garçon sur 2). \quad b) Fille porteuse : $\dfrac{1}{4} = \mathbf{25\,\%}$ (soit 1 fille sur 2). \quad c) Enfant atteint tous sexes confondus : seuls les garçons $X^h Y$ sont atteints, soit $\mathbf{25\,\%}$.
    \item Le nombre de combinaisons gamétiques par brassage interchromosomique seul est $2^n$, où $n$ est le nombre de paires de chromosomes homologues. Ici $n=3$ :
    \[ 2^3 = \mathbf{8\ combinaisons\ différentes}. \]
\end{enumerate}
\end{corrige}

\begin{attention}[Exercice 4 — Points de vigilance]
Pour les maladies liées à l'X, ne jamais oublier que le père transmet son X à toutes ses filles et son Y à tous ses fils. Distinguer « probabilité parmi tous les enfants » et « probabilité parmi les garçons ».
\end{attention}

\begin{corrige}[Exercice 5 — Analyse comparée de séquences ADN : la drépanocytose]
\begin{enumerate}
    \item \textbf{Type de mutation :} comparaison codon par codon : seul le 4\up{e} codon \emph{de l'extrait fourni} (correspondant au codon 6 du gène \emph{HBB}) diffère : GAG (témoin) $\to$ GTG (patient). Il s'agit d'une \cle{substitution} : l'adénine (A) en 2\up{e} position du codon 6 du gène \emph{HBB} est remplacée par une thymine (T) : transversion A $\to$ T. Aucune base n'est ajoutée ni supprimée : pas de décalage du cadre de lecture.
    \item \textbf{Traduction :}
    \begin{itemize}
        \item Témoin : Leu -- Thr -- Pro -- \textbf{Glu} -- Glu -- Lys -- Ser
        \item Patient : Leu -- Thr -- Pro -- \textbf{Val} -- Glu -- Lys -- Ser
    \end{itemize}
    \item \textbf{Acide aminé modifié :} le glutamate (Glu) en position 6 de la chaîne $\beta$-globine est remplacé par une valine (Val). Le Glu est un acide aminé \emph{acide, polaire, hydrophile, chargé négativement} ; la Val est un acide aminé \emph{apolaire, hydrophobe, neutre}. C'est donc un changement \textbf{polaire chargé $\to$ apolaire}, en surface de la protéine.
    \item \textbf{Conclusion :} la drépanocytose est une anomalie moléculaire due à une \textbf{mutation ponctuelle non-synonyme} (faux-sens) du gène $\beta$-globine : substitution A $\to$ T au codon 6, provoquant la substitution Glu6 $\to$ Val. La valine apolaire en surface crée un « patch » hydrophobe qui, en conditions d'\textbf{hypoxie} (hémoglobine désoxygénée, conformation T), se fixe sur une poche hydrophobe d'une molécule d'HbS voisine : les molécules d'hémoglobine S s'agrègent en longues fibres (polymérisation), déformant le globule rouge en faucille (drépanocyte), d'où hémolyse et vaso-occlusions.
\end{enumerate}
\end{corrige}

\begin{attention}[Exercice 5 — Points de vigilance]
Toujours lire la séquence codon par codon et repérer la position exacte de la base modifiée. Ne pas confondre « substitution » (remplacement d'une base) et « délétion » (perte d'une base, qui décale le cadre de lecture si elle ne porte pas sur un multiple de 3). Préciser le numéro du codon dans la protéine mature (codon 6) et non seulement dans l'extrait.
\end{attention}

\begin{corrige}[Exercice 6 — Interprétation de caryotypes]
\begin{enumerate}
    \item \textbf{Formules normalisées :}
    \begin{itemize}
        \item Caryotype A : \textbf{47,XX,+21} (47 chromosomes, femme, chromosome 21 surnuméraire).
        \item Caryotype B : \textbf{45,X} (45 chromosomes, un seul chromosome sexuel X).
        \item Caryotype C : \textbf{47,XXY} (47 chromosomes, deux X et un Y).
    \end{itemize}
    \item \textbf{Syndromes :} A = trisomie 21 (syndrome de Down) chez une fille ; B = syndrome de Turner (monosomie X) ; C = syndrome de Klinefelter.
    \item \textbf{Caryotype A :} la trisomie 21 libre (non mosaïque) résulte d'une \textbf{non-disjonction méiotique} de la paire 21, le plus souvent chez la mère (environ 90\,\% des cas), en méiose I (majoritaire) ou en méiose II. Le gamète à 24 chromosomes ($n+1$, deux chromosomes 21) fécondé par un spermatozoïde normal donne un zygote 47,+21. Ce n'est pas un mosaïcisme (qui résulterait d'une non-disjonction \emph{mitotique} post-zygotique et donnerait deux lignées cellulaires, souvent atténuées). Le facteur de risque majeur est l'\textbf{âge maternel} : les ovocytes I restent bloqués en prophase I depuis la vie fœtale jusqu'à l'ovulation, parfois plus de 40 ans, ce qui fragilise le fuseau méiotique et les cohésines, favorisant les non-disjonctions.
    \item \textbf{Caryotype C (47,XXY) :} le phénotype est masculin car le chromosome \textbf{Y porte le gène \emph{SRY}} (déterminant testiculaire) : dès qu'un Y est présent, les gonades indifférenciées évoluent en testicules. Un des deux X est inactivé (corpuscule de Barr). Signes pubertaires majeurs : \textbf{testicules de petite taille et fermes} (hypogonadisme, azoospermie, infertilité) et \textbf{gynécomastie} (développement mammaire), avec grande taille et déficit androgénique.
    \item \textbf{Caryotype B (45,X) :} chez cette patiente, l'unique chromosome X n'a \textbf{pas d'homologue} : en prophase I, il ne peut s'apparier (pas de synapsis, pas de crossing-over). La méiose est bloquée ou anarchique ; de plus, l'absence du second X provoque une dégénérescence massive des ovocytes dès la vie fœtale (dysgénésie gonadique : ovaires en « stries » fibreuses). Conséquence : \textbf{aménorrhée primaire et stérilité quasi systématique} (grossesses spontanées très rares, $<5\,\%$).
\end{enumerate}
\end{corrige}

\begin{corrige}[Exercice 7 — Arbre généalogique : l'hémophilie A]
\begin{enumerate}
    \item \textbf{Génération I :} I-2 est un homme atteint : son unique X porte l'allèle mutant : génotype $X^h Y$. Le couple a des fils atteints (II-2) : ceux-ci ont reçu leur X de leur mère I-1, qui est saine mais a transmis $X^h$ : I-1 est donc \textbf{porteuse obligée} : génotype $X^H X^h$.
    \item \textbf{Générations II et III :}
    \begin{itemize}
        \item II-1 (femme saine, fille d'un père atteint $X^h Y$) : elle reçoit obligatoirement $X^h$ de son père et $X^H$ de sa mère porteuse ou non ; étant saine, elle est \textbf{porteuse obligée} : $X^H X^h$. (Confirmation : son fils III-1 est atteint.)
        \item II-2 (homme atteint) : $X^h Y$.
        \item II-3 (femme saine, fille de père atteint) : \textbf{porteuse obligée} : $X^H X^h$.
        \item II-4 (homme sain) : $X^H Y$.
        \item Conjoints II-1m et II-3m (hommes sains) : $X^H Y$.
        \item III-1 (garçon atteint) : $X^h Y$ (a reçu $X^h$ de sa mère II-1).
        \item III-2 (fille saine de II-1 $X^H X^h$ $\times$ $X^H Y$) : génotype $X^H X^H$ ou $X^H X^h$ avec probabilité 1/2 chacun (porteuse potentielle).
        \item III-3 (garçon sain de II-3 porteuse) : $X^H Y$ (il a reçu le $X^H$ maternel).
        \item III-4 (fille saine de II-3) : $X^H X^H$ ou $X^H X^h$ (1/2 chacun).
    \end{itemize}
    \item \textbf{Couple II-3 ($X^H X^h$) $\times$ II-3m ($X^H Y$) :} probabilité d'un garçon atteint $X^h Y$ : $P(\text{garçon}) \times P(X^h \text{ maternel}) = \dfrac{1}{2} \times \dfrac{1}{2} = \dfrac{1}{4} = \mathbf{25\,\%}$ de tous les enfants (soit 1 garçon sur 2).
    \item \textbf{Pourquoi seuls les hommes sont atteints :} la maladie est récessive liée à l'X. L'homme, hémizygote ($XY$), exprime la maladie dès que son unique X porte $h$. La femme, avec deux X, n'est malade que si elle est $X^h X^h$ (père atteint \emph{et} mère au moins porteuse), situation très rare ; les femmes $X^H X^h$ sont porteuses saines car l'allèle normal $H$ suffit à assurer une coagulation correcte.
    \item \textbf{Diagnostic prénatal au Cameroun :} prélèvement de cellules fœtales par \textbf{amniocentèse} (15--18 SA) ou \textbf{biopsie de villosités choriales} (11--13 SA), suivi d'une étude moléculaire (PCR/séquençage du gène \emph{F8}) ou d'un caryotype pour déterminer le sexe fœtal ; réalisable dans des centres de référence (Centre Mère-Enfant de la Fondation Chantal Biya, Hôpital Général de Yaoundé), avec conseil génétique préalable. On peut aussi citer le dosage du facteur VIII sur sang de cordon.
\end{enumerate}
\end{corrige}

\begin{attention}[Exercice 7 — Points de vigilance]
Toute fille d'un père hémophile est \emph{obligatoirement} porteuse. Un homme atteint ne transmet jamais l'allèle à ses fils (il leur donne le Y). Bien justifier chaque génotype par la filiation des chromosomes X.
\end{attention}

\begin{corrige}[Exercice 8 — Brassage des allèles et unicité génétique]
\begin{enumerate}
    \item Trois gènes indépendants, deux allèles chacun : nombre de types de gamètes $= 2^3 = \mathbf{8}$ :
    \[ ABC,\quad ABc,\quad AbC,\quad Abc,\quad aBC,\quad aBc,\quad abC,\quad abc. \]
    \item \textbf{Autofécondation :} chaque parent produit 8 types de gamètes, d'où $8 \times 8 = 64$ combinaisons de fécondation. Nombre de génotypes distincts : pour chaque gène, 3 génotypes possibles ($AA$, $Aa$, $aa$), donc $3^3 = \mathbf{27\ génotypes}$. Nombre de phénotypes (dominance complète) : 2 par gène, donc $2^3 = \mathbf{8\ phénotypes}$.
    \item \textbf{Crossing-over entre le locus $A$ et le centromère :} en prophase I (pachytène), la tétade comporte 4 chromatides ; deux chromatides non sœurs échangent réciproquement un segment. Résultat : \textbf{2 chromatides recombinantes sur 4} (les deux autres restent parentales). Schéma simplifié :
    \begin{popfigure}
    \begin{tikzpicture}[scale=0.9, transform shape]
    % Tétrade avant crossing-over
    \draw[very thick, popblue] (0,0) -- (0,3);
    \draw[very thick, popblue] (0.35,0) -- (0.35,3);
    \draw[very thick, poporange] (0.8,0) -- (0.8,3);
    \draw[very thick, poporange] (1.15,0) -- (1.15,3);
    \fill[popdark] (0.175,1.5) circle (0.07);
    \fill[popdark] (0.975,1.5) circle (0.07);
    \node[left, font=\small] at (0,2.6) {$A$};
    \node[left, font=\small] at (0,0.4) {$A$};
    \node[right, font=\small] at (1.15,2.6) {$a$};
    \node[right, font=\small] at (1.15,0.4) {$a$};
    \node[font=\small\bfseries] at (0.6,-0.6) {Tétrade (pachytène)};
    % Flèche
    \draw[->, very thick, popdark] (2,1.5) -- (3.2,1.5) node[midway, above, font=\small] {crossing-over};
    % Après crossing-over
    \draw[very thick, popblue] (4,0) -- (4,3);
    \draw[very thick, popblue] (4.35,1.8) -- (4.35,3);
    \draw[very thick, poporange] (4.35,0) -- (4.35,1.8);
    \draw[very thick, poporange] (4.8,1.8) -- (4.8,3);
    \draw[very thick, popblue] (4.8,0) -- (4.8,1.8);
    \draw[very thick, poporange] (5.15,0) -- (5.15,3);
    \fill[popdark] (4.175,1.5) circle (0.07);
    \fill[popdark] (4.975,1.5) circle (0.07);
    \node[left, font=\small] at (4,2.6) {$A$};
    \node[left, font=\small] at (4,0.4) {$A$};
    \node[right, font=\small] at (5.15,2.6) {$a$};
    \node[right, font=\small] at (5.15,0.4) {$a$};
    \node[font=\small\bfseries] at (4.6,-0.6) {Après échange : 2 chromatides recombinées};
    \end{tikzpicture}
    \legende{Crossing-over entre deux chromatides non sœurs d'une paire de chromosomes homologues (bleu = chromosome paternel portant $A$ ; orange = chromosome maternel portant $a$ ; points noirs = centromères). Après l'échange, 2 chromatides sur 4 sont recombinées.}
    \end{popfigure}
    \item \textbf{Unicité génétique :} chez l'humain ($n=23$ paires), le brassage interchromosomique seul produit $2^{23} \approx \num{8,4}$ millions de types de gamètes par individu ; un couple peut donc générer $(2^{23})^2 = 2^{46} \approx 7 \times 10^{13}$ combinaisons zygotiques, \emph{sans} compter le brassage intrachromosomique. Or chaque paire de chromosomes subit en moyenne 1 à 3 crossing-overs par méiose, à des positions variables : le nombre de gamètes possibles devient astronomique (bien supérieur à $10^{12}$). La probabilité que deux fécondations indépendantes donnent le même génotype est donc quasi nulle : seuls les vrais jumeaux monozygotes (issus d'un même zygote) partagent un patrimoine génétique identique.
\end{enumerate}
\end{corrige}

\begin{corrige}[Exercice 9 — Sujet type Bac : Trisomie 21 et âge maternel]
\textbf{Barème indicatif (sur 20) :} Q1 : 4 pts ; Q2 : 5 pts ; Q3 : 3 pts ; Q4 : 4 pts ; Q5 : 4 pts.
\begin{enumerate}
    \item \textbf{Analyse de la courbe :} la prévalence reste faible et presque stable avant 30 ans : environ 0,5 pour 1000 naissances à 20 ans, 0,7 pour 1000 à 25 ans et 1 pour 1000 à 30 ans. Elle augmente ensuite lentement (2,5 pour 1000 à 35 ans, 4 pour 1000 à 37 ans) puis \emph{exponentiellement} après 38--40 ans : 10 pour 1000 à 40 ans, 16 pour 1000 à 42 ans et environ 30 pour 1000 à 45 ans (soit 3\,\%). Entre 30 et 45 ans, le risque est multiplié par 30 environ.
    \item \textbf{Interprétation :} hypothèse : le vieillissement de l'ovocyte augmente la fréquence des non-disjonctions de la paire 21. Rappel : tous les ovocytes I entrent en méiose pendant la vie fœtale et restent \textbf{bloqués en prophase I (stade dictyotène)} de la naissance jusqu'à l'ovulation. À chaque cycle ovarien après la puberté, un ovocyte reprend et achève la méiose I. Une femme de 40 ans ovule donc un ovocyte dont les chromosomes sont appariés depuis 40 ans : les \textbf{cohésines} maintenant les chromatides sœurs et les protéines du fuseau se dégradent avec le temps, ce qui favorise les erreurs de ségrégation (non-disjonction en méiose I surtout).
    \item \textbf{Calcul de risque :} à 38 ans, la courbe indique environ 6 pour 1000 (interpolation entre 4 pour 1000 à 37 ans et 10 pour 1000 à 40 ans), soit $\approx \mathbf{0,6\,\%}$. Pour 500 femmes de 38 ans : $500 \times 0{,}006 = \mathbf{3}$ enfants trisomiques 21 attendus statistiquement (toute valeur entre 2 et 5 est acceptable si cohérente avec la lecture graphique).
    \item \textbf{Diagnostic prénatal invasif :}
    \begin{itemize}
        \item \textbf{Biopsie de villosités choriales (BVC)} : réalisée tôt (11--13 SA), résultat rapide ; inconvénient : risque de fausse couche légèrement plus élevé (environ 1\,\%).
        \item \textbf{Amniocentèse} : réalisée à partir de 15 SA, très fiable (caryotype fœtal sur cellules du liquide amniotique) ; inconvénient : plus tardive, risque de fausse couche d'environ 0,5 à 1\,\%.
    \end{itemize}
    Technique \textbf{non invasive} en développement : le dépistage par \textbf{ADN fœtal libre circulant} dans le sang maternel (DPNI, dès 10 SA), sans risque pour le fœtus, mais restant un test de \emph{dépistage} à confirmer par caryotype.
    \item \textbf{Éthique et éducation :} arguments attendus : la trisomie 21 est un \textbf{accident chromosomique aléatoire} (non-disjonction méiotique) qui peut survenir dans toute famille, sans lien avec le comportement des parents, ni faute morale, ni « sorcellerie » ; sa probabilité augmente simplement avec l'âge maternel, ce qui est mesurable et démontré par l'épidémiologie. Sur le plan humain : l'enfant trisomique bénéficie d'une prise en charge précoce (suivi médical, éveil, scolarisation adaptée) et d'une pleine inclusion familiale et sociale ; le dépistage prénatal relève d'un choix éclairé des parents, encadré par un conseil génétique respectueux.
\end{enumerate}
\end{corrige}

\begin{attention}[Exercice 9 — Points de vigilance]
Exigences du correcteur : valeurs chiffrées \emph{avec unités} (pour 1000 naissances) ; rappel explicite du blocage de l'ovocyte I en prophase I ; distinction dépistage/diagnostic ; calcul de risque cohérent avec la lecture graphique ; ton respectueux et non stigmatisant dans la question d'éthique.
\end{attention}

\begin{corrige}[Exercice 10 — Sujet type Bac : Synthèse argumentée « Mutation vs Brassage »]
\textbf{Barème indicatif (sur 20) :} introduction et définitions : 4 pts ; partie I (mutation) : 5 pts ; partie II (brassage) : 5 pts ; exemples et articulation : 3 pts ; conclusion : 2 pts ; qualité rédactionnelle et vocabulaire : 1 pt.

\textbf{Corrigé rédigé (modèle) :}

\emph{Introduction.} L'apparition d'un caractère nouveau dans une famille — groupe sanguin inattendu, maladie héréditaire inédite — suscite souvent des interrogations, voire des préjugés. Or la génétique distingue deux mécanismes fondamentaux de la variabilité : la \cle{mutation}, qui crée des allèles nouveaux, et le \cle{brassage} des allèles, qui recompose des allèles déjà existants. Le caractère nouveau est-il toujours le produit d'une mutation ? Nous montrerons que non : la mutation est la source ultime de la nouveauté, mais le brassage méiotique suffit le plus souvent à expliquer l'apparition d'un phénotype inédit dans une fratrie.

\emph{I. La mutation, créatrice \emph{de novo} d'allèles nouveaux.} Une mutation est une modification aléatoire et héritable du matériel génétique. Les mutations \emph{géniques} (substitutions, insertions, délétions de nucléotides) altèrent la séquence d'un gène : ainsi, la substitution A $\to$ T du codon 6 du gène $\beta$-globine a créé l'allèle $Hb^S$ responsable de la drépanocytose. Les mutations \emph{chromosomiques} modifient la structure (délétions, translocations) ou le nombre des chromosomes (trisomie 21 par non-disjonction). Survenant au hasard lors de la réplication ou de la méiose, la mutation est le seul mécanisme capable de faire apparaître un allèle qui n'existait chez aucun des deux parents : c'est la source première de toute nouveauté génétique.

\emph{II. Le brassage, recombinant d'allèles préexistants.} Cependant, la grande majorité des « caractères nouveaux » observés dans les familles ne doit rien à une mutation récente. Lors de la méiose, le \textbf{brassage interchromosomique} (répartition aléatoire des homologues : $2^{23}$ combinaisons chez l'humain) et le \textbf{brassage intrachromosomique} (crossing-over en prophase I) génèrent des gamètes portant des combinaisons d'allèles inédites, puis la fécondation associe au hasard deux gamètes. Ainsi, deux parents de groupe sanguin A (génotype $I^A/i$) peuvent avoir un enfant de groupe O ($i/i$) : l'allèle $i$ existait, masqué, chez les parents ; seule sa \emph{recombinaison} est nouvelle. De même, un enfant drépanocytaire $SS$ né de parents porteurs sains $AS$ ne résulte pas d'une mutation nouvelle, mais de la rencontre de deux allèles $S$ hérités.

\emph{Conclusion.} L'apparition d'un caractère nouveau dans une famille n'est donc pas toujours due à une mutation : le plus souvent, elle traduit le brassage d'allèles anciens, transmis silencieusement depuis des générations. Mutation et brassage sont complémentaires : la première crée, à l'échelle de l'évolution, la matière première de la biodiversité ; le second la décline, à chaque génération, en une diversité qui rend chaque individu génétiquement unique. Comprendre ces mécanismes, c'est aussi déconstruire les préjugés qui frappent les familles.
\end{corrige}

\begin{attention}[Exercice 10 — Points de vigilance]
Le correcteur attend : des définitions rigoureuses (allèle, locus, gamète, crossing-over, non-disjonction employés à bon escient), une distinction claire « allèle nouveau » (mutation) vs « combinaison nouvelle » (brassage), au moins un exemple chiffré ou concret correctement exploité, et une conclusion qui répond explicitement à la question posée.
\end{attention}

\begin{corrige}[Exercice 11 — Mini-projet : affiche de sensibilisation sur la drépanocytose]
\textbf{Contenu textuel complet et structuré de l'affiche (prêt à mettre en page) :}

\textbf{TITRE (bandeau supérieur, grands caractères) :} « La drépanocytose : connaître son statut, c'est protéger ses enfants ! »

\textbf{BLOC 1 — D'où vient la maladie ? (origine scientifique).} La drépanocytose est une maladie \emph{génétique héréditaire} du sang, causée par une mutation du gène de la $\beta$-globine, situé sur le chromosome 11. Cette mutation transforme l'hémoglobine normale (HbA) en hémoglobine anormale (HbS) qui déforme les globules rouges en faucilles. Elle se transmet sur le mode \emph{autosomique récessif} : il faut recevoir l'allèle $S$ des deux parents pour être malade. Trois statuts possibles : $AA$ = non porteur (sain) ; $AS$ = porteur sain (trait drépanocytaire, en bonne santé) ; $SS$ = drépanocytaire majeur (malade). [Schéma à dessiner : trois silhouettes étiquetées $AA$, $AS$, $SS$ avec globules rouges ronds ou en faucille.]

\textbf{BLOC 2 — Quels risques pour nos enfants ? (carré de Punnett $AS \times AS$).}
\begin{center}
\begin{tabular}{|c|c|c|}\hline
\rowcolor{popblueL} & $A$ (père) & $S$ (père) \\ \hline
$A$ (mère) & $AA$ : enfant non porteur (25\,\%) & $AS$ : porteur sain (25\,\%) \\ \hline
$S$ (mère) & $AS$ : porteur sain (25\,\%) & $SS$ : enfant malade (25\,\%) \\ \hline
\end{tabular}
\end{center}
Message central : « Si les deux partenaires sont $AS$, à \textbf{chaque grossesse} : 1 enfant sur 4 (25\,\%) risque d'être malade ($SS$), 1 sur 2 (50\,\%) sera porteur sain ($AS$), 1 sur 4 (25\,\%) sera non porteur ($AA$). »

\textbf{BLOC 3 — Comment connaître son statut ? (dépistage au Cameroun).}
\begin{itemize}
    \item \textbf{Électrophorèse de l'hémoglobine} : examen de référence, fiable et peu coûteux, disponible dans les laboratoires des hôpitaux de district et régionaux.
    \item \textbf{Test de solubilité (test d'Emmel)} : test rapide de dépistage de masse, à confirmer par électrophorèse.
    \item \textbf{Dépistage néonatal} (goutte de sang au talon) : déployé progressivement par le Programme National de Lutte contre la Drépanocytose (PNLD).
    \item \textbf{Diagnostic prénatal} (biopsie de villosités choriales ou amniocentèse + analyse moléculaire PCR) : centres spécialisés de Yaoundé et Douala.
\end{itemize}

\textbf{BLOC 4 — Vivre avec la drépanocytose : prévenir et soigner.} Vaccination anti-pneumococcique ; prophylaxie par pénicilline V chez le nourrisson jusqu'à 5 ans ; acide folique quotidien ; hydratation abondante ; traitement de fond par hydroxyurée ; suivi régulier dans un Centre de Prise en Charge de la Drépanocytose (Hôpital Mère-Enfant de la Fondation Chantal Biya à Yaoundé, Hôpital Laquintinie de Douala).

\textbf{BLOC 5 — Brisons les préjugés ! (bandeau coloré).} « La drépanocytose n'est ni une malédiction, ni une sorcellerie : c'est une maladie génétique héréditaire, écrite dans les chromosomes, que personne ne choisit. Faire le dépistage avant le mariage (certificat de génotype) est un acte d'amour et de responsabilité. Ne stigmatisons ni les malades ni les porteurs sains : soutenons-les. »

\textbf{BLOC 6 — Contacts utiles (pied d'affiche).} ASCAD (Association Camerounaise de Lutte contre la Drépanocytose) ; Programme National de Lutte contre la Drépanocytose (PNLD, Ministère de la Santé Publique) ; centre de santé le plus proche.
\end{corrige}

\begin{attention}[Exercice 11 — Points de vigilance]
Le jury évalue : l'exactitude des génotypes et des pourcentages (25/50/25), le caractère autosomique récessif correctement formulé, la mention de structures de santé camerounaises réelles, et un message de lutte contre les préjugés fondé sur des arguments scientifiques (mutation du chromosome 11, hasard de la transmission) et non sur des croyances.
\end{attention}

\newpage

\coursec{Fiche bilan}

\textbf{Carte mentale de la leçon}\par
\begin{popfigure}
\begin{tikzpicture}[
    mindmap,
    grow cyclic,
    every node/.style={concept, circular drop shadow, execute at begin node=\hskip0pt, font=\small\bfseries},
    root concept/.append style={concept color=popdark, minimum size=2.8cm, text width=2.6cm, align=center, font=\normalsize\bfseries},
    level 1/.append style={level distance=4.2cm, sibling angle=60, minimum size=2.2cm, text width=2.0cm, align=center, font=\small\bfseries},
    level 2/.append style={level distance=3cm, sibling angle=40, minimum size=1.6cm, text width=1.5cm, align=center, font=\scriptsize}
]
\node[root concept, align=center]{Anomalies \&\\ Nouveaux Caractères}
    child[concept color=popblue] { node[concept, align=center]{I. Mutations\\Origine allèles} }
    child[concept color=popgreen] { node[concept, align=center]{II. Anomalies géniques\\Drépanocytose} }
    child[concept color=poporange] { node[concept, align=center]{III. Autres anomalies\\géniques} }
    child[concept color=poppurple] { node[concept, align=center]{IV. Anomalies\\chromosomiques} }
    child[concept color=poppink] { node[concept, align=center]{V. Brassage allèles\\Diversité} }
    child[concept color=popgold] { node[concept, align=center]{VI. Lutte anti-préjugés\\Éducation} };
\end{tikzpicture}
\legende{Carte mentale récapitulative de la leçon 05 : des mutations à l'éradication des préjugés.}
\end{popfigure}

\begin{aretenir}[Bilan synthétique de la leçon 05]
\textbf{Définitions clés}\par
\begin{center}
\begin{tabularx}{\linewidth}{|>{\bfseries}l|X|}\hline
\rowcolor{popblueL} \textbf{Terme} & \textbf{Définition rigoureuse} \\\hline
\cle{Mutation génique} & Modification de la séquence nucléotidique d'un gène (substitution, insertion, délétion) créant un \cle{nouvel allèle}. Source de la variabilité allélique. \\\hline
\cle{Mutation chromosomique} & Modification de la structure (délétion, duplication, inversion, translocation) ou du \cle{nombre de chromosomes} (anéuploïdie, polyploïdie). \\\hline
\cle{Drépanocytose} & Anomalie génique autosomique récessive : substitution A$\to$T (codon 6 $\beta$-globine) $\Rightarrow$ Glu$\to$Val $\Rightarrow$ HbS polymérise en désoxygénation. \\\hline
\cle{Hémophilie A/B} & Anomalie génique \cle{liée au X} récessive : déficit en facteur VIII (A) ou IX (B) $\Rightarrow$ coagulation défaillante. \\\hline
\cle{Trisomie 21} & Anomalie chromosomique numérique : 3 chromosomes 21 (non-disjonction méiotique maternelle $\approx$ 95\% cas). Syndrome de Down. \\\hline
\cle{Syndrome de Turner} & Monosomie X (45,X) : non-disjonction paternelle fréquente. Phénotype féminin, stérilité, dysgénésie gonadique. \\\hline
\cle{Syndrome de Klinefelter} & 47,XXY (non-disjonction maternelle/paternelle). Phénotype masculin, hypogonadisme, infertilité. \\\hline
\cle{Brassage interchromosomique} & Assortiment indépendant des paires homologues en méiose I $\Rightarrow$ $2^n$ combinaisons gamétiques ($n$=nombre de paires). \\\hline
\cle{Brassage intrachromosomique} & \cle{Crossing-over} (chiasmes) en prophase I : échange de segments entre chromatides homologues non-sœurs $\Rightarrow$ allèles recombinés. \\\hline
\end{tabularx}
\end{center}

\textbf{Principes \& Repères à connaître par cœur}\par
\begin{center}
\begin{tabularx}{\linewidth}{|>{\bfseries}l|X|}\hline
\rowcolor{popgreenL} \textbf{Notion} & \textbf{Énoncé / Valeur} \\\hline
Code génétique & Universel, dégénéré, sans chevauchement. Une mutation change le sens $\Leftrightarrow$ nouveau phénotype possible. \\\hline
Non-disjonction & Échec de séparation des chromosomes (anaphase I) ou chromatides (anaphase II) $\Rightarrow$ gamètes $n+1$ et $n-1$. \\\hline
Risque trisomie 21 & Croît avec l'âge maternel : $\approx$ 1/1500 (20 ans), 1/900 (30 ans), 1/100 (40 ans), 1/30 (45 ans). \\\hline
Transmission récessive autosomique & Parents hétérozygotes sains $\times$ sains $\Rightarrow$ 25\% atteints, 50\% porteurs, 25\% sains non porteurs. \\\hline
Transmission liée au X récessive & Mère porteuse $\times$ père sain $\Rightarrow$ 50\% fils atteints, 50\% filles porteuses. Père atteint $\times$ mère saine $\Rightarrow$ toutes filles porteuses, fils sains. \\\hline
Diversité gamétique & $N = 2^n \times (\text{crossing-over})$ ; $n=23$ chez l'homme $\Rightarrow$ $2^{23} \approx 8,4 \times 10^6$ sans crossing-over. \\\hline
Unicité individuelle & Brassage méiotique + fécondation aléatoire $\Rightarrow$ probabilité génotype identique $\approx$ 0 (sauf vrais jumeaux). \\\hline
\end{tabularx}
\end{center}

\textbf{Méthodes express}\par
\begin{enumerate}
\item \textbf{Analyse de séquence ADN :} Aligner séquence normale/mutée $\rightarrow$ identifier type (substitution/délétion/insertion) $\rightarrow$ traduire en AA $\rightarrow$ conclure (synonyme, non-sens, faux-sens, décalage de cadre).
\item \textbf{Analyse de caryotype :} Compter chromosomes, identifier paires homologues $\rightarrow$ repérer excès/défaut (trisomie/monosomie) ou anomalie structurelle (bande G).
\item \textbf{Arbre généalogique :} Repérer sexe (carré/cercle), phénotype (plein/vide), générations $\rightarrow$ tester hypothèses (auto/récessif, auto/dominant, X/récessif, X/dominant, Y) $\rightarrow$ déduire génotypes probables.
\item \textbf{Calcul risque :} Déterminer génotypes parents $\rightarrow$ croiser (Punnett) $\rightarrow$ fréquence phénotypes/génotypes attendus.
\item \textbf{Argumentation anti-préjugés :} 1) Hasard méiotique (non-disjonction, mutation \emph{de novo}) 2) Aucune responsabilité parentale (comportement, alimentation, « sort ») 3) Conseil génétique et diagnostic prénatal (DPI, DPNI) pour choix éclairés.
\end{enumerate}
\end{aretenir}

\begin{aretenir}[Checklist avant le Bac]
$\square$ Je définis sans ambiguïté \textbf{mutation génique} (échelle gène) et \textbf{mutation chromosomique} (échelle chromosome : nombre/structure).\\
$\square$ J'explique au niveau moléculaire l'origine de la \textbf{drépanocytose} : substitution A$\to$T (codon 6 $\beta$-globine) $\Rightarrow$ Glu$\to$Val $\Rightarrow$ HbS $\Rightarrow$ falciformation en désoxygénation.\\
$\square$ Je cite \textbf{deux anomalies géniques} (ex. hémophilie A, phénylcétonurie, myopathie de Duchenne, mucoviscidose) et je précise leur \textbf{mode de transmission} (autosome/X, récessif/dominant).\\
$\square$ J'identifie sur un \textbf{caryotype} une trisomie 21 (3 chr 21), un syndrome de Turner (45,X) et un syndrome de Klinefelter (47,XXY).\\
$\square$ J'explique la \textbf{non-disjonction} (anaphase I ou II) comme cause majeure des anomalies numériques et je sais qu'elle est \textbf{aléatoire} et \textbf{non héréditaire} (sauf rare translocation).\\
$\square$ Je distingue \textbf{brassage interchromosomique} ($2^n$ combinaisons, assortiment indépendant) et \textbf{intrachromosomique} (crossing-over, chiasmes, allèles recombinés).\\
$\square$ Je calcule le nombre théorique de combinaisons gamétiques ($2^{23}$) et j'explique pourquoi chaque individu (sauf jumeaux monozygotes) est \textbf{génétiquement unique}.\\
$\square$ Je sais \textbf{lire un arbre généalogique} : symboles, déduire le mode de transmission, calculer un risque de récidive pour une grossesse à venir.\\
$\square$ J'argumente scientifiquement contre les préjugés : \textbf{hasard biologique} (mutation, non-disjonction), \textbf{absence de cause comportementale}, rôle du \textbf{conseil génétique} et du \textbf{dépistage prénatal} (DPNI, amniocentèse) pour l'autonomie reproductive.\\
$\square$ Je connais l'intérêt du \textbf{dépistage néonatal} (talon, drépanocytose, phénylcétonurie) pour la prise en charge précoce (ex. centre de drépanocytose au Cameroun).\\
$\square$ Je relie \textbf{biodiversité} et \textbf{santé} : hétérozygote drépanocytaire (AS) confère résistance au paludisme (avantage sélectif en zone endémique camerounaise).\\
$\square$ Je rédige une conclusion structurée : Observation $\rightarrow$ Hypothèse $\rightarrow$ Expérience/Document $\rightarrow$ Interprétation $\rightarrow$ Conclusion, en vocabulaire précis (allèle, locus, gamète, zygote, phénotype, génotype).
\end{aretenir}

\findecours

\end{document}
